\documentclass[11pt,a4paper]{article}
\usepackage[T1]{fontenc}
\usepackage[utf8]{inputenc}
\usepackage[english]{babel}
\usepackage{textcomp}
\usepackage{amsmath, amssymb, amsfonts, amsthm}
\usepackage{geometry}
\geometry{top=2.0cm, bottom=2.0cm, left=2.0cm, right=2.0cm}
\usepackage{booktabs}
\usepackage{tabularx}
\usepackage{hyperref}
\usepackage{float}
\usepackage{xltabular}
\usepackage{array}
\usepackage{underscore}

\providecommand{\lambdabar}{\bar{\lambda}}
\hypersetup{colorlinks=true, linkcolor=blue, citecolor=blue, urlcolor=blue}

\emergencystretch=2em

\newtheorem{theorem}{Theorem}
\newtheorem{lemma}[theorem]{Lemma}
\theoremstyle{definition}
\newtheorem{definition}{Definition}
\newtheorem{postulate}{Postulate}
\theoremstyle{remark}
\newtheorem{remark}{Remark}

\begin{document}
	
\title{\textbf{\Large VACUUM TOPOLOGICAL ELASTODYNAMICS}\\
\large Theoretical model of fundamental structures and interactions\\[2mm]
\small \texttt{Version 5.7: Deductive Directed Acyclic Graph (DAG) of the continuum.}}

\author{\textbf{Dmitriy Mykhilenko}\\[1mm]
\small DOI: ---\\
\small Open project repository and Lean 4 core: \href{https://github.com/ainjohansen/-}{\texttt{https://github.com/ainjohansen/-}}\\[0.5mm]
\small Interactive 3D laboratories: \href{https://ainjohansen.github.io/-/}{\texttt{https://ainjohansen.github.io/-/}}}

\date{\today}

\maketitle

\begin{abstract}
\noindent The present work formulates a closed theoretical model of physical space and fundamental interactions, based on the mechanics of continuous media. The physical vacuum is considered not as an empty geometric manifold with external fields, but as a material three-dimensional quasi-incompressible elastic continuum with a spinor order parameter. Elementary particles are stable vortex defects of the medium (filaments, tori and knots), bounded by a finite radius of plastic saturation.

In the model the empirical arbitrariness of gauge constants is eliminated: the spectral mass ratios of leptons and bosons, the coupling constants, the quark and neutrino mixing angles, the parameters of the nuclear forces and the gravitational hierarchy are derived through an oriented Directed Acyclic Graph (DAG) from two material anchors --- the coefficient of wave impedance matching and the mass scale of the main baryon resonator. The speed of light and the Planck constant express the kinematics of transverse shear waves and the topological holonomy of the medium.

Electrodynamics is described through the isomorphism of elastic displacements, which excludes the existence of magnetic monopoles and defines the electric charge as the surface momentum flux of the cavity. Newtonian gravitation and the effects of general relativity are derived as a result of the shear tension of the matrix around Eshelby inclusions. The equality of the phase velocities of transverse and longitudinal waves is proven, and the predominant transverse character of gravitational radiation and the high degree of suppression of the scalar mode are explained by the electrostriction of the pre-stressed Coulomb tail of the defect.

The anomalies of the rotation curves of galaxies and the Baryon asymmetry of the Universe ($\eta_B$) find a natural mechanical explanation in the elasticity of the sphere of the Hubble horizon and in the chiral crystallization of knots, without invoking hypotheses about dark matter or hidden mass. The mathematical structure of the model has been verified by methods of interactive theorem proving in the Lean 4 system without external axioms, and the numerical consequences are consistent with the experimental databases of fundamental constants.
\end{abstract}
	
\tableofcontents
\newpage

\section{Ontological basis and equations of state of the elastic continuum}
	
\subsection{Methodological status of the model and limits of applicability}

The present work formulates not a speculative ``theory of everything'', but a theoretical model of a continuous medium describing the interaction of coupled self-oscillatory processes of various degrees of quality factor ($Q$) in a single elastic substrate of the vacuum.

The model is deliberately restricted to a verifiable range of physical scales (exactly 60 orders of magnitude), within which experimental verification is available:

\begin{enumerate}
\item \textbf{The infrared boundary (IR) --- the Hubble sphere ($R_H \sim 10^{26}\text{ m}$):}
The material continuum at the cosmological horizon is closed by the Hubble sphere of radius $R_H = c / H_0$. The background radial tension of the medium at this macroscopic boundary, at the reference cosmological Planck value ($H_0 \approx 67.4\text{ km/s/Mpc} \approx 2.184 \times 10^{-18}\text{ s}^{-1}$), sets the minimum threshold of elastic acceleration (the MOND scale):
\begin{equation}
a_0 = \frac{c H_0}{2\pi} \approx 1.042 \times 10^{-10}\text{ m/s}^2.
\label{eq:mond_acceleration_exact}
\end{equation}
All hypothetical constructions concerning the physical reality ``beyond the Hubble sphere'' are recognized as speculative owing to the absence of a causal connection with our instruments.

\item \textbf{The ultraviolet boundary (UV) --- Planck cavitation ($l_P \sim 10^{-35}\text{ m}$):}
The continuous medium of the vacuum does not admit mathematical point singularities with infinite energy density. At the scale of the Planck core $l_P \sim \sqrt{\hbar G_N / c^3}$ elastodynamics is cut off by the physical limit of the plastic yield of Huber--von Mises:
\begin{equation}
\sigma_{\mathrm{yield}} = \alpha G_0 \approx 6.356 \times 10^{72}\text{ Pa},
\label{eq:mises_yield_cutoff}
\end{equation}
which sets the saturation of the maximum volumetric density of elastic energy $\mathcal{E}_{\mathrm{max}} = \frac{1}{2}\alpha^2 G_0$. The extrapolation of the field equations into the sub-Planckian region is recognized as non-operational.
\end{enumerate}

Within this 60-order resonator all observed phenomenology of the microworld reduces to the spectroscopy of self-oscillations: stable particles are non-damping topological modes ($Q \to \infty$), while unstable particles are stressed high-frequency resonators of finite quality factor, shedding excess stress back into the elastic matrix.

\subsection{Rejection of emptiness and the relational paradigm}
Twentieth-century fundamental physics relies on the postulate of geometric emptiness in which point material particles are placed. This approach generates mathematical divergences: the infinite self-energy of a point charge ($E \propto e^2/r \to \infty$ as $r \to 0$), the necessity of ultraviolet cutoffs, and the introduction of 26 a posteriori fitting constants of the Standard Model.
	
Topological elastodynamics of the vacuum (TEV) returns to the materialist ontology of the mechanics of continuous media:
\begin{enumerate}
\item \textbf{Space is a material medium.} The physical world is a homogeneous elastic continuum. Space does not contain matter --- it is itself the substance.
\item \textbf{A particle is a topological knot of the medium.} Matter is not imported into the vacuum from outside. An elementary particle is a topologically stable vortex defect (a knot or a filament) tied out of the fabric of the vacuum itself.
\item \textbf{The principle of absence of singularities.} In a continuous continuum there are no infinite quantities ($K_0 \ne \infty$, $\sigma \ne \infty$) and no absolutely empty regions. At the center of any defect, at the scale of the core $r_0$ the continuous elastic stresses reach the limit of plastic yield $\sigma_{\mathrm{yield}} = \alpha G_0$, bringing the medium into a state of phase saturation with a finite energy density $\mathcal{E}_{\mathrm{core}} \sim \alpha^2 G_0$.
\end{enumerate}
	
\subsection{Equations of state, solenoidal condition and the Poisson bound}
	
\begin{postulate}[Material continuum of the vacuum]
	The physical vacuum is a three-dimensional isotropic elastic continuum $\mathbb{R}^3$, characterized by the shear modulus $G_0$, the volumetric mass density $\rho_0$ and the bulk modulus $K_0$.
\end{postulate}
	
In early mechanical models the quasi-incompressibility of the medium was frequently associated with the limiting passage of the phase velocity of longitudinal acoustic compression waves to infinity ($c_L \gg c$), which led to a violation of causality. In TEV this approach is replaced by a kinematic constraint.
	
\begin{postulate}[Solenoidal condition of the wave spectrum]
	In the spectrum of physical propagating wave perturbations of the vacuum the longitudinal dynamical modes are absent:
\begin{equation}
	N_{\mathrm{longitudinal}} = 0.
	\label{eq:n_longitudinal_zero}
\end{equation}
	The incompressibility condition of the medium:
\begin{equation}
	\nabla \cdot \mathbf{u}(\mathbf{x}, t) = 0
	\label{eq:solenoidal_constraint}
\end{equation}
	is a kinematic constraint of the continuum, isomorphic to the Coulomb gauge of the vector potential in electrodynamics ($\nabla \cdot \mathbf{A} = 0$).
\end{postulate}
	
The dynamical equation of motion of the elastic medium~\cite{Navier1827,Cauchy1828} in the presence of the constraint (\ref{eq:solenoidal_constraint}) is formulated variationally with a Lagrange multiplier $p(\mathbf{x}, t)$ having the physical meaning of the hydrostatic pressure of the constraint reactions:
\begin{equation}
	\rho_0 \frac{\partial^2 \mathbf{u}}{\partial t^2} = -\nabla p + G_0 \nabla^2 \mathbf{u}.
	\label{eq:navier_cauchy_lagrange}
\end{equation}
\noindent\small\textbf{Formalized:} \texttt{TevSolenoidalBulk.lean}, theorems \texttt{solenoidal_smul}, \texttt{solenoidal_pressure_integral_kernel}, \texttt{divergence\_const\_smul}.

Taking the divergence of both sides of (\ref{eq:navier_cauchy_lagrange}) with $\nabla \cdot \mathbf{u} = 0$ yields an elliptic equation for the Lagrange multiplier:
\begin{equation}
	\nabla^2 p(\mathbf{x}, t) = 0 \quad (\text{outside the vortex cores of the defects}).
	\label{eq:laplace_pressure}
\end{equation}
	
The pressure field $p(\mathbf{x}, t)$ specifies the distribution of the internal reaction stresses of the medium (analogously to the Coulomb potential in the Coulomb gauge of QED), but it is not an independent wave degree of freedom and does not transport dynamical energy. The velocity of a physical perturbation equals the phase velocity of the transverse shear waves:
\begin{equation}
	c \equiv c_T = \sqrt{\frac{G_0}{\rho_0}}.
	\label{eq:speed_of_light_def}
\end{equation}
	
\begin{theorem}[Estimate of the Poisson ratio in terms of the elastic moduli]
\label{thm:poisson_bound_exact}
In a continuum with shear modulus $G_0 > 0$ and bulk modulus $K_0 > 0$ the deviation of the Poisson ratio $\nu$ from the limiting value of an incompressible medium $1/2$ is bounded by an analytic inequality:
\begin{equation}
	\left| \nu - \frac{1}{2} \right| \le \frac{G_0}{2 K_0} \to 0 \quad (\text{for } K_0 \gg G_0).
	\label{eq:poisson_bound_ineq}
\end{equation}
\end{theorem}

\noindent\small\textbf{Formalized:} \texttt{TevElasticityBoundary.lean}, theorem \texttt{poisson_quasi_incompressible_bound} (also \texttt{poisson_deviation_exact}, \texttt{lame_to_bulk_identity}).

\begin{proof}
The Poisson ratio of an isotropic body is expressed through the elastic moduli:
\begin{equation}
	\nu = \frac{3K_0 - 2G_0}{2(3K_0 + G_0)} = \frac{1}{2} - \frac{3G_0}{2(3K_0 + G_0)}.
\end{equation}
Since $G_0 > 0$ and $3K_0 + G_0 > 3K_0$, the absolute value of the difference is estimated as:
\begin{equation}
	\left| \nu - \frac{1}{2} \right| = \frac{3G_0}{2(3K_0 + G_0)} < \frac{3G_0}{2(3K_0)} = \frac{G_0}{2K_0}.
\end{equation}
\end{proof}
	
\subsection{Spinor order parameter of the continuum}
	
\begin{postulate}[Internal spinor orientation]
The local state of the shear orientation of the micro-elements of the continuum is described by a normalized spinor field $\boldsymbol{\Phi}(\mathbf{x}) \in \mathrm{SU}(2) \cong S^3$:
\begin{equation}
	\boldsymbol{\Phi}(\mathbf{x}) = \begin{pmatrix} z_1(\mathbf{x}) \\ z_2(\mathbf{x}) \end{pmatrix}, \quad |z_1|^2 + |z_2|^2 = 1.
	\label{eq:spinor_order_parameter}
\end{equation}
\end{postulate}
	
The spinor field defines a three-dimensional rotation of the matrix micro-elements through the covering homomorphism $\mathrm{SU}(2) \to \mathrm{SO}(3)$. A complete topological revolution in physical space $\mathbb{R}^3$ by an angle $2\pi$ inverts the sign of the spinor: $\boldsymbol{\Phi} \to -\boldsymbol{\Phi}$. The periodicity of a closed vortex perturbation requires a double revolution by $4\pi$ (the Dirac belt trick), which naturally generates the kinematics of fermions with spin $S = 1/2$.
	
\subsection{Metrological audit: architecture of the two material inputs}
	
To avoid parametric over-determination and the mixing of the physical properties of the medium with the units of measurement, we perform a classification of all quantities of the model in an oriented acyclic graph:
	
\begin{enumerate}
\item \textbf{Two fundamental material inputs (2 degrees of freedom):}
\begin{itemize}
	\item \textit{Wave impedance-matching coefficient} (Huber--von Mises yield strain threshold):
\begin{equation}
	\alpha \equiv \gamma_{\mathrm{yield}} = \frac{\sigma_{\mathrm{yield}}}{G_0} = \frac{Z_0}{2 R_K} \approx \frac{1}{137.035999177}.
	\label{eq:alpha_impedance_def}
\end{equation}
	It is a continuous input for the dimensionless spectrum (particle mass ratios, mixing angles, coupling constants and the gravitational hierarchy $\alpha_G$).
	\item \textit{Gauge anchor of the defect scale:} the rest mass of the primary stable trefoil knot T(3,2) $M_p = 938.27208943\text{ MeV}$ (proton mass, CODATA 2022), fixing the spatial radius of the core:
\begin{equation}
	\lambda_p = \frac{\hbar}{M_p c} \approx 0.21030891\text{ fm},
	\label{eq:lambda_p_anchor_def}
\end{equation}
	converting the dimensionless eigenvalues into the energy scale.
\end{itemize}
		
\item \textbf{Topological and kinematic invariants:}
\begin{itemize}
	\item Dimensionality of the space $d = 3$, giving the number of independent components of the shear deviator $N_{3D} = \frac{d(d+1)}{2} - 1 = 5$ and the dimensionality of the moving Frenet--Serret frame of the filament $N_{1D} = 3$;
	\item Winding numbers of the knot on the Clifford torus $(p, q) = (3, 2)$, corresponding to the minimum of the Dirac--Laplace functional ($I(3,2) = 67/5 = 13.4$) and determining the dimensionality of the basis $\mathcal{N}_{\mathrm{basis}} = p^2 + q^2 = 13$;
	\item Parameters of the two-sheeted covering $\mathrm{Spin}(3) \to \mathrm{SO}(3)$: multiplicity $N_{\mathrm{rev}} = 2$ (period $4\pi$) and real dimensionality $N_{\mathrm{spin}} = 4$;
	\item Holonomic solenoidal condition $\nabla \cdot \mathbf{u} = 0$, giving the Poisson incompressibility limit $\nu = 1/2$.
\end{itemize}
		
	\item \textbf{Metrological SI basis:}
\begin{itemize}
	\item Speed of the transverse waves $c = 299\,792\,458\text{ m/s}$ — the conversion coefficient between the units of length and time ($c \equiv \sqrt{G_0 / \rho_0}$);
	\item Quantum of action $\hbar = 1.054571817 \times 10^{-34}\text{ J}\cdot\text{s}$ — the intrinsic mechanical angular momentum of the vortex core of the base soliton ($\hbar \equiv M_p c \lambda_p$), linking the oscillation frequency with the macroscopic energy;
	\item Elementary charge $e = 1.602176634 \times 10^{-19}\text{ C}$ (magnetic permeability $\mu_0 = \frac{2\alpha h}{c e^2}$) — the conversion bridge between the mechanical stresses of the matrix $[\text{Pa}]$ and the electrodynamical units $[\text{V}, \text{A}]$.
\end{itemize}
\item \textbf{Derived macro-parameters of the elastic matrix:}
	The mass density of the medium $\rho_0$, the elastic shear modulus $G_0$ and the fundamental tension of the vortex filament $\mu_{\mathrm{str}}$ are not independent input parameters. They are deductively computed through the two-block transfer cascade $13 \oplus 5 = 18$:
\begin{align}
	\alpha_G &= \sqrt{3}\,\alpha^{18}\left(1 - \frac{21}{5}\frac{\alpha}{\pi}\right) \approx 5.905792 \times 10^{-39}, \label{eq:alpha_G_master_def} \\
	G_N^{\mathrm{pred}} &= \alpha_G \frac{\hbar c}{M_p^2} \approx 6.67395 \times 10^{-11}\text{ m}^3/(\text{kg}\cdot\text{s}^2), \label{eq:G_N_master_def} \\
\mu_{\mathrm{str}} &= \frac{c^2}{\pi G_N} = \frac{M_p^2 c}{\pi \hbar \alpha_G} \approx 4.2865 \times 10^{26}\text{ kg/m}, \label{eq:mu_str_master_def} \\
T_0 &= \mu_{\mathrm{str}} c^2 \approx 3.8524 \times 10^{43}\text{ N}, \label{eq:T_0_master_def} \\
\rho_0 &= \frac{\mu_{\mathrm{str}}}{\lambda_p^2} = \frac{M_p^2 c}{\pi \hbar \alpha_G \lambda_p^2} \approx 9.6915 \times 10^{57}\text{ kg/m}^3, \label{eq:rho_0_master_def} \\
G_0 &= \rho_0 c^2 \approx 8.710 \times 10^{74}\text{ Pa}, \quad \tau_c = \alpha G_0 \approx 6.356 \times 10^{72}\text{ Pa}. \label{eq:G_0_master_def}
\end{align}
\end{enumerate}

\begin{theorem}[Relation of the vacuum yield threshold to the Frenkel theoretical shear strength limit]
\label{thm:frenkel_to_mises_exact}
Local shear failure of the elastic matrix at the wall of the vortex core ($r = r_0$) occurs upon reaching the classical limit of the Frenkel theoretical shear strength for a $2\pi$-periodic potential:
\begin{equation}
    \tau_{\mathrm{local}}(r_0) = \tau_{\mathrm{Frenkel}} = \frac{G_0}{2\pi}.
    \label{eq:frenkel_local_limit}
\end{equation}
Owing to the toroidal geometry of the soliton with the core-to-ring aspect ratio $r_0 / R_e = \alpha$ and phase circulation $2\pi$, the effective macroscopic cavitation threshold of the vacuum matrix is:
\begin{equation}
    \sigma_{\mathrm{yield}} = \tau_{\mathrm{Frenkel}} \times (2\pi \alpha) = \left(\frac{G_0}{2\pi}\right) \times 2\pi \alpha \equiv \alpha G_0.
    \label{eq:frenkel_cancellation_exact}
\end{equation}
\end{theorem}

\noindent\small\textbf{Formalized:} \texttt{TevBridge.lean}, theorems \texttt{nambu_no_leak_H1}, \texttt{nambu_no_leak_H2}. \texttt{TevLeptonMonopole.lean}, theorems \texttt{div_curl_identically_zero}, \texttt{magnetic_monopole_charge_zero}.

\begin{proof}
The period of the sinusoidal shear barrier equals $2\pi$, which gives the Frenkel denominator $2\pi$. At the same time, the phase circulation of the velocity around the vortex tube along a closed contour $\oint \mathbf{v} \cdot \mathrm{d}\mathbf{l} = 2\pi r_0 c$ contributes a phase factor $2\pi$ to the numerator of the volume stress integral. The $2\pi$ factors cancel, fixing the macroscopic yield threshold $\sigma_{\mathrm{yield}} = \alpha G_0$.
\end{proof}
Thus, the mathematical system of the TEV relies on two material inputs $(\alpha, M_p)$.

\begin{table}[htbp]
\centering
\small
\setlength{\tabcolsep}{4pt}
\renewcommand{\arraystretch}{1.2}
\caption{Comparative audit of the fundamental degrees of freedom: Standard Model vs TEV}
\label{tab:parameter_accounting_sm_vs_tev}
\vspace{2mm}
\begin{tabularx}{\textwidth}{>{\raggedright\arraybackslash}p{4.2cm} >{\centering\arraybackslash}p{2.5cm} >{\centering\arraybackslash}p{2.5cm} >{\raggedright\arraybackslash}X}
\toprule
Physics sector & SM parameters & TEV inputs & Deduction mechanism in the TEV \\
\midrule
Lepton masses ($e, \mu, \tau$) & 3 Yukawa constants & \textbf{0} & Mode spectrum of the Clifford torus $W(n) = (2n-1)/n$ \\
Quark masses ($u, d, s, c, b, t$) & 6 Yukawa constants & \textbf{0} & Nambu quantum $E_0 = \frac{5}{67}M_p$ \\
CKM mixing angles & 3 angles + 1 $CP$ phase & \textbf{0} & Lode geometry $\theta_0 = 2/9$, $4\pi$ covering, $\delta_{\mathrm{CP}}=\arctan(3/\sqrt{2})$ \\
PMNS neutrino mixing & 3 angles + 2 phases & \textbf{0} & Torus basis $\sin^2(2\theta_{13}) = 1/12$, reduction $5 \to 3$ \\
Neutrino masses ($\Delta m^2$) & 3 parameters & \textbf{0} & 1D-string scale $\mathcal{M}_\nu \propto \alpha^5 (M_p/3)$ \\
Gauge couplings ($g_1, g_2, g_3$) & 3 coupling constants & \textbf{0} & $\sin^2\theta_W = 3/13$, $\alpha_s = 2/17(1 + 1.5\alpha/\pi)$ \\
Electroweak vacuum ($v, M_H$) & 2 parameters ($\mu^2, \lambda$) & \textbf{0} & Plastic failure $E_{\mathrm{EW}} = M_p/\alpha$, $v^2/E_{\mathrm{EW}}^2 = 11/3$ \\
Strong $CP$ problem ($\theta_{\mathrm{QCD}}$) & 1 parameter & \textbf{0} & Isochoricity of plasticity ($\operatorname{Tr}\dot{\boldsymbol{\varepsilon}}^p = 0 \implies \theta \equiv 0$) \\
Gravitation ($G_N, \Lambda$) & 2 GR constants & \textbf{0} & Cascade $\alpha_G = \sqrt{3}\alpha^{18}(1-4.2\alpha/\pi)$, $\Lambda \equiv 0$ \\
\midrule
\textbf{Continuous material inputs} & \textbf{26 free parameters} & \textbf{2 INPUTS} & $\alpha = Z_0/(2R_K)$ and the proton mass $M_p$ \\
\midrule
\textit{Discrete topological basis} & \textit{$SU(3)\times SU(2)\times U(1)$, 3 generations, $\theta$-vacuum} & \textit{$d=3$, knot $T(3,2)$, $N_{\mathrm{spin}}=4, N_{\mathrm{rev}}=2$} & In both models the discrete basis is not a tunable continuous parameter \\
\bottomrule
\end{tabularx}
\end{table}

\subsection{Ontological status of the fine-structure constant}
	
\begin{theorem}[Material nature of the constant $\alpha$]
\label{thm:alpha_materiality_proof}
The dimensionless constant $\alpha$ is the ratio of the material responses of an isotropic elastic continuum:
\begin{equation}
\alpha \equiv \frac{Z_0}{2 R_K} = \frac{\mu_0 c}{2 (h/e^2)} \equiv \frac{e^2}{4\pi\varepsilon_0 \hbar c}
\end{equation}
\end{theorem}

\noindent\small\textbf{Formalized:} \texttt{TevLeptonMonopole.lean}, theorems \texttt{gauss_cavitation_boundary_flux}, \texttt{electron_charge_is_one}, \texttt{charge_is_unit_iff_two_sheeted}.

\begin{proof}
Consider the emptiness predicate (Minkowski space without a material substrate).
\begin{enumerate}
\item The wave impedance $Z_0 = \mu_0 c$~\cite{Gordon1923} expresses the dynamic resistance of the medium to the propagation of a transverse shear deformation. In a continuum with density $\rho_0$ and shear modulus $G_0$ it equals $Z_0 = \sqrt{\rho_0 G_0}$. A space devoid of substance possesses neither inertia density ($\rho_0 = 0$) nor elastic stiffness ($G_0 = 0$), so the notion of wave resistance loses its physical meaning ($Z_0 \equiv \emptyset$).
\item The von Klitzing constant $R_K = h/e^2$ sets the quantum of Hall resistance of a one-dimensional channel with phase circulation $\oint \nabla\theta \cdot \mathrm{d}\mathbf{l} = 2\pi$. In the absence of a continuous medium with an order parameter $\boldsymbol{\Phi} \in S^3$, formation of the topological vortex core is impossible ($R_K \equiv \emptyset$).
\item The quantity $\alpha$ is the matching coefficient of the wave impedances of the three-dimensional volume and the one-dimensional vortex resonator. The impedance ratio exists exclusively in a medium with elastic and phase properties.
\end{enumerate}
Hence, the experimentally registered finite value $\alpha \approx 1/137.036$ serves as confirmation of the materiality of the vacuum.
\end{proof}

\begin{remark}[Ontological status of the constant]
The number $1/137$ expresses the ratio of the compliance of the elastic body of the vacuum upon excitation of a one-dimensional vortex quantum channel to the resistance of the volumetric shear deformation of the matrix.
\end{remark}
	
\subsection{Lorentz invariance and the dynamical theory of relativity}
	
Introducing a material medium of high density does not contradict the precision tests of special relativity thanks to the Lorentz--Poincaré dynamical concept:
	
\begin{enumerate}
	\item \textbf{Principle of internal kinematics:} Material observers, rods, and reference clocks consist of topological solitons and transverse waves of the same elastic medium.
	\item \textbf{Dynamic Lorentz contraction:} The Navier--Cauchy wave equation for a stable wave packet moving through the medium with velocity $\mathbf{v}$ undergoes deformation along the direction of motion:
\begin{equation}
	\left(1 - \frac{v^2}{c^2}\right) \frac{\partial^2 \mathbf{u}}{\partial x^2} + \frac{\partial^2 \mathbf{u}}{\partial y^2} + \frac{\partial^2 \mathbf{u}}{\partial z^2} = \frac{1}{c^2} \frac{\partial^2 \mathbf{u}}{\partial t^2}.
	\label{eq:navier_cauchy_moving}
\end{equation}
The spatial scale of the packet contracts by the relativistic factor $\gamma = 1/\sqrt{1 - v^2/c^2}$. The interatomic bonds of the measuring rod shorten in the direction of motion:
\begin{equation}
	\Delta L = \Delta L_0 \sqrt{1 - \frac{v^2}{c^2}}.
	\label{eq:lorentz_length_contraction}
\end{equation}
\item \textbf{Slowing of the internal clock rate:} The proper frequency of self-oscillations of the moving soliton decreases proportionally to $\sqrt{1 - v^2/c^2}$ owing to the elongation of the Zitterbewegung phase trajectory.
\end{enumerate}
	
	An internal observer performing a measurement of the speed of light in both directions (``there and back'') with deformed rods and slowed clocks obtains a constant value $c$ in any inertial frame of reference. The principle of relativity acts as a dynamical symmetry of the equations of motion of the elastic continuum for internal observers.
	
\section{Electro-elastodynamic isomorphism\\ of MacCullagh--Whittaker}
	
\subsection{Relation between elastic displacements and electromagnetic fields}
	
In 1839--1848 J.~MacCullagh showed that the equations of a rotational elastic medium are formally isomorphic to the equations of optics; in 1910 E.~Whittaker established their correspondence with Maxwell's electrodynamics~\cite{MacCullagh1848,Whittaker1910}.
	
In a quasi-incompressible continuum ($\nabla \cdot \mathbf{u} = 0$) we define the vector and scalar electrodynamic potentials through the displacement field $\mathbf{u}(\mathbf{x}, t)$:
\begin{equation}
	\mathbf{A}(\mathbf{x}, t) = \sqrt{\mu_0 G_0} \, \mathbf{u}(\mathbf{x}, t), \quad \varphi \equiv 0 \quad (\text{Coulomb gauge } \nabla \cdot \mathbf{A} = 0).
	\label{eq:potentials_def}
	\end{equation}
	
The electromagnetic field strengths are defined through the gauge relations:
\begin{align}
\mathbf{B} &= \nabla \times \mathbf{A} = \sqrt{\mu_0 G_0} \, (\nabla \times \mathbf{u}), \label{eq:field_B_def} \\
\mathbf{E} &= -\frac{\partial \mathbf{A}}{\partial t} = -\sqrt{\mu_0 G_0} \, \frac{\partial \mathbf{u}}{\partial t}. \label{eq:field_E_def}
\end{align}
	
Substituting these definitions into Hooke's law for the volume energy density of pure shear:
\begin{equation}
	\mathcal{E}_{\mathrm{elast}} = \frac{1}{2} \rho_0 \|\dot{\mathbf{u}}\|^2 + \frac{1}{2} G_0 \|\nabla \times \mathbf{u}\|^2,
	\label{eq:elastic_energy_density_raw}
\end{equation}
and taking into account the relation $c^2 = G_0 / \rho_0 = 1/(\varepsilon_0 \mu_0)$, we obtain:
\begin{equation}
	\mathcal{E}_{\mathrm{elast}} = \frac{1}{2} \rho_0 \left\| -\frac{\mathbf{E}}{\sqrt{\mu_0 G_0}} \right\|^2 + \frac{1}{2} G_0 \left\| \frac{\mathbf{B}}{\sqrt{\mu_0 G_0}} \right\|^2 = \frac{\varepsilon_0 \mathbf{E}^2}{2} + \frac{\mathbf{B}^2}{2\mu_0} \equiv \mathcal{E}_{\mathrm{EM}}.
	\label{eq:em_energy_density_match}
\end{equation}
The energy density of the elastic shear deformation coincides with the volume density of Maxwell's electromagnetic energy.
	
\subsection{Theorem on the absence of magnetic monopoles}
	
In field theory the absence of magnetic monopoles is postulated by the equation $\nabla \cdot \mathbf{B} = 0$. In topological elastodynamics this property is derived as a theorem of continuum mechanics.
	
\begin{theorem}[Absence of magnetic charges in a smooth continuum]
\label{thm:no_monopoles}
If the physical vacuum is a continuous connected medium whose displacement vector satisfies the smoothness condition $\mathbf{u} \in C^2(\mathbb{R}^3)$, then the magnetic charge $Q_m$ of any closed volume $V \subset \mathbb{R}^3$ vanishes:
\begin{equation}
	\nabla \cdot \mathbf{B} \equiv 0 \implies Q_m = \oint_{\partial V} \mathbf{B} \cdot \mathrm{d}\mathbf{S} \equiv 0.
	\label{eq:div_B_zero_theorem}
\end{equation}
\end{theorem}

\noindent\small\textbf{Formalized:} \texttt{TevMaxwell.lean}, theorems \texttt{no\_magnetic\_monopoles}, \texttt{curl\_curl\_eq\_neg\_laplacian\_of\_solenoidal}, \texttt{ampere\_maxwell\_bridge}.


\begin{proof}
We compute the divergence of the magnetic field (\ref{eq:field_B_def}) in tensor components:
\begin{equation}
	\nabla \cdot \mathbf{B} = \partial_i B_i = \sqrt{\mu_0 G_0} \, \partial_i (\epsilon_{ijk} \partial_j u_k) = \sqrt{\mu_0 G_0} \, \epsilon_{ijk} \, (\partial_i \partial_j u_k).
\end{equation}
By Schwarz's lemma on the equality of mixed partial derivatives for a twice differentiable field $\mathbf{u} \in C^2$:
\begin{equation}
	\partial_i \partial_j u_k = \partial_j \partial_i u_k.
\end{equation}
The contraction of the antisymmetric Levi-Civita tensor $\epsilon_{ijk} = -\epsilon_{jik}$ with the symmetric operator of second derivatives $\partial_i \partial_j$ vanishes:
\begin{equation}
	\epsilon_{ijk} \, (\partial_i \partial_j u_k) = \frac{1}{2} (\epsilon_{ijk} + \epsilon_{jik}) \partial_i \partial_j u_k \equiv 0.
\end{equation}
Applying the Gauss--Ostrogradsky theorem:
\begin{equation}
	Q_m = \oint_{\partial V} \mathbf{B} \cdot \mathrm{d}\mathbf{S} = \int_V (\nabla \cdot \mathbf{B}) \, \mathrm{d}^3\mathbf{x} \equiv 0.
\end{equation}
\end{proof}
	
\begin{remark}[Geometric derivation]
A magnetic monopole is impossible in a continuous space. Its existence would require a singular topological break of the connectivity of the medium (a Dirac string with a break of continuity), which is excluded by the finiteness of the yield limit of the elastic vacuum.
\end{remark}
	
\subsection{Topological nature of the electric charge}
	
Since the vacuum is quasi-incompressible ($\nabla \cdot \mathbf{u} = 0$), the divergence of the electric field in the volume, according to (\ref{eq:field_E_def}), equals:
\begin{equation}
	\nabla \cdot \mathbf{E} = -\frac{\partial}{\partial t}(\nabla \cdot \mathbf{A}) = -\sqrt{\mu_0 G_0} \, \frac{\partial}{\partial t}(\nabla \cdot \mathbf{u}) = 0,
\label{eq:div_E_gauge}
\end{equation}
where $[\mathbf{u}] = \text{m}$, the divergence $\nabla \cdot \mathbf{u}$ is dimensionless, and both sides have the dimension $[\text{V/m}^2]$.
	
The Coulomb electric charge $\oint \mathbf{E} \cdot \mathrm{d}\mathbf{S} = q/\varepsilon_0$ arises as a consequence of the topology of the internal cavitation boundary of the core of the defect (the Eshelby inclusion mechanism).
	
Consider a region $\Omega \subset \mathbb{R}^3$ containing a vortex soliton. At its center the stresses exceed the threshold $\sigma_{\mathrm{yield}} = \alpha G_0$, forming a plastic core of radius $r_0$. The region of linear elasticity is a manifold with an internal boundary:
\begin{equation}
	\Omega_{\mathrm{elastic}} = \Omega \setminus V_{\mathrm{core}}, \quad \partial \Omega_{\mathrm{elastic}} = S_\infty \cup S_{\mathrm{core}},
\label{eq:manifold_boundaries}
\end{equation}
where $S_\infty$ is the outer sphere of large radius, while $S_{\mathrm{core}} = \partial V_{\mathrm{core}}$ is the internal surface of the vortex core.
	
\begin{theorem}[Electric charge as a surface flux of the cavity momentum]
\label{thm:gauss_charge_boundary}
The electric charge of the soliton equals the integral flux of the radial dynamic displacement of the elastic matrix through the internal boundary of the cavitation core $S_{\mathrm{core}}$:
\begin{equation}
	\Phi_E = \oint_{S_\infty} \mathbf{E} \cdot \mathrm{d}\mathbf{S} = \oint_{S_{\mathrm{core}}} \mathbf{E} \cdot \mathrm{d}\mathbf{S} = \frac{q}{\varepsilon_0}.
	\label{eq:gauss_charge_flux}
\end{equation}
\end{theorem}

\noindent\small\textbf{Formalized:} \texttt{TevBridge.lean}, theorem \texttt{koide_exact_from_mises_balance}; \texttt{TevQuantum.lean}, theorem \texttt{koide_exact_from_mises_balance}.


\begin{proof}
	We apply the Gauss--Ostrogradsky theorem to the region $\Omega_{\mathrm{elastic}}$ between the outer sphere and the core:
\begin{equation}
	\int_{\Omega_{\mathrm{elastic}}} (\nabla \cdot \mathbf{E}) \, \mathrm{d}^3\mathbf{x} = \oint_{S_\infty} \mathbf{E} \cdot \mathrm{d}\mathbf{S} - \oint_{S_{\mathrm{core}}} \mathbf{E} \cdot \mathrm{d}\mathbf{S}.
\end{equation}
By the solenoidal character of the continuum outside the core $\nabla \cdot \mathbf{E} = 0$. Consequently:
\begin{equation}
	\oint_{S_\infty} \mathbf{E} \cdot \mathrm{d}\mathbf{S} = \oint_{S_{\mathrm{core}}} \mathbf{E} \cdot \mathrm{d}\mathbf{S}.
\end{equation}
The field $\mathbf{E} \propto -\dot{\mathbf{u}}$ is generated by the circulation of the phase of the spinor field $\boldsymbol{\Phi} \in S^3$. Since the spinor covering has degree $N_{\mathrm{rev}} = 2$ (period $4\pi$), the circulation of the displacement vector along the M\"obius strip generates a radial kinematic recoil momentum of the matrix on the cavity wall. The integral of this flux is quantized by the topological index of the boundary map $\pi_3(S^3) \cong \mathbb{Z}$:
\begin{equation}
	q = \varepsilon_0 \oint_{S_{\mathrm{core}}} \mathbf{E} \cdot \mathrm{d}\mathbf{S} = \pm 1 \, e.
	\label{eq:charge_quantization_exact}
\end{equation}
\end{proof}
	
	Electric charge is the reaction of an incompressible medium to the continuous vortex motion of the inner wall of the cavitation cavity.
	
\section{Compactification hierarchy: the dimension ladder of defects}
\noindent\small\textbf{Formalized in the kernel:} \texttt{TevTopology.lean}, theorems \texttt{basis\_dimension\_eq\_13}, \texttt{deviator\_dimension\_3d\_eq\_5}, \texttt{cascade\_exponent\_eq\_18}, \texttt{three\_generations\_bound}.
	
	The emergence of rest mass in a continuous medium is a process of geometric localization of the transverse shear wave momentum. In a continuum with a spinor order parameter $\boldsymbol{\Phi} \in \mathrm{SU}(2)$ this process forms a dimension ladder $d = 0 \to 1 \to 2 \to 3$:
	
\begin{enumerate}
\item \textbf{Photon ($d = 0$ topological localization / open plane wave):}
A free wavepacket of transverse Navier--Cauchy shear. The wavefront is not closed, the phase topology is trivial ($\pi_1 = 0$), the rest mass is absent:
\begin{equation}
	m_\gamma = 0, \quad v = c = \sqrt{\frac{G_0}{\rho_0}}.
	\label{eq:photon_mass_zero}
\end{equation}
Energy transport is carried out at the velocity of transverse shear waves without formation of a cavitation volume.
		
\item \textbf{Neutrino ($d = 1$, open one-dimensional vortex thread):}
Upon reaching the shear threshold the wavefront rolls up along the propagation line into a solenoidal circulation cord with the topological quantum of phase $\oint \nabla\theta \cdot \mathrm{d}\mathbf{l} = 2\pi$. The translational momentum is conserved along the thread axis, the transverse motion is localized by the Frenet--Serre trihedron ($N_{1D} = 3$). Due to the absence of a closed cavitation volume the thread does not retain an electrostatic charge:
\begin{equation}
	Q_\nu = 0.
	\label{eq:neutrino_charge_zero}
\end{equation}
A small rest mass is formed by the resistance to bending and torsion of the thread~\cite{Mindlin1964}:
\begin{equation}
	\mathcal{M}_\nu \propto \alpha^5 \left(\frac{M_p}{3}\right) \approx 12.29\text{ meV}, \quad v \approx c.
	\label{eq:neutrino_mass_scale_concept}
\end{equation}
		
\item \textbf{Charged lepton ($d = 2$, closed vortex ring / torus $T^2$):}
The linear momentum flux of the vortex thread closes into a ring, forming a stationary vortex torus $T^2$. Translational motion passes into self-oscillations of the Zitterbewegung limit cycle:
\begin{equation}
	\omega_0 = \frac{M c^2}{\hbar}.
	\label{eq:zitterbewegung_freq_def}
\end{equation}
At the center of the torus an excavated cavitation cavity $\Delta V = M / \rho_0$ is formed. The radial gradient of the kinematic pressure of the matrix generates the electric charge $q = \pm 1e$, while the balance of volumetric and shear stresses fixes the Koide invariant $Q = 2/3$.
		
The degree of tension of the vortex tube on the torus surface forms a discrete generation spectrum with winding $W(n) = (2n-1)/n$:
\begin{itemize}
	\item $n=1$: the electron $(1,1)$ — an unstressed configuration ($\Delta E = 0$, stable);
	\item $n=2$: the muon $(3,2)$ — a compressed spring ($\Delta E = 1/2$, metastable);
	\item $n=3$: the tau-lepton $(5,3)$ — a strongly compressed spring ($\Delta E = 2/3$, decay).
\end{itemize}
		
\item \textbf{Baryon ($d = 3$, self-knotted three-dimensional knot $T(3,2)$):}
The vortex tube performs 3 poloidal turns and 2 toroidal encirclements, closing into a Trefoil knot ($T(3,2)$) on the Clifford torus with topological index $\pi_3(S^3) \cong \mathbb{Z}$ (baryon number $B=1$).
		
		Translational motion is closed in three spatial dimensions. Three standing-wave antinodes form internal quasiparticle centers (``quarks'') with charges ($+1, +1, -1$), acting as a three-phase transformer with a shear of $120^\circ$. The vector balance of phases ensures the radiative stability of the proton ($\tau_p > 10^{34}$ yr).
\end{enumerate}
\noindent\small\textbf{Formalized in the kernel:} \texttt{TevLeptonTorus.lean}, theorems \texttt{zero_baryon_iff_torus}, \texttt{leptons_are_not_knots}, \texttt{proton_distinct_from_lepton}.

\section{Huber--von Mises yield strength, the stress deviator and the Koide invariant}

\subsection{Wave-impedance matching and the scale of shear failure}

In quantum electrodynamics the fine-structure constant $\alpha \approx 1/137.035999$ is postulated as a fundamental dimensionless coupling constant $e^2 / (4\pi \varepsilon_0 \hbar c)$. In topological elastodynamics this quantity represents the ratio of the characteristic wave impedances of the volumetric matrix of the vacuum and of the one-dimensional vortex channel.

A free three-dimensional continuum possesses a wave impedance for transverse shear waves:
\begin{equation}
Z_0 = \sqrt{\frac{\mu_0}{\varepsilon_0}} = \mu_0 c \approx 376.730313668\ \Omega.
\label{eq:wave_impedance_z0}
\end{equation}

A quantized vortex deformation thread with phase circulation $\oint \nabla\theta \cdot \mathrm{d}\mathbf{l} = 2\pi$ is characterized by the von Klitzing quantum of resistance:
\begin{equation}
R_K = \frac{h}{e^2} \approx 25812.80745930\ \Omega.
\label{eq:klitzing_rk_def}
\end{equation}

\begin{definition}[Coefficient of wave-impedance matching]
The coupling coefficient of the one-dimensional circulation channel with the three-dimensional volume is defined by the ratio of the matrix impedance to twice the vortex quantum of resistance:
\begin{equation}
\alpha \equiv \frac{Z_0}{2 R_K} = \frac{\mu_0 c}{2 (h/e^2)} = \frac{e^2}{4\pi \varepsilon_0 \hbar c} \approx \frac{1}{137.035999177}.
\label{eq:alpha_def}
\end{equation}
The factor 2 in the denominator expresses the two-sheeted spinor covering $N_{\mathrm{rev}} = 2$ (period $4\pi$).
\end{definition}

\subsection{Deductive derivation of the form of the Huber---von Mises criterion}

In linear elasticity theory the medium is infinitely strong, and the formation of solitons is impossible. The formation of a stable defect requires the specification of a criterion of plastic yield of the matrix.

Unlike the phenomenological theory of metals, where the Mises criterion is postulated empirically, in TEV its form is derived deductively from Nambu kinematics:

\begin{theorem}[Isochoric nature of Nambu plastic failure]
\label{thm:isochoric_mises_form}
In a quasi-incompressible continuum with a spinor order parameter ($\nabla \cdot \mathbf{u} = 0$, $K_0 \gg G_0$) the kinematics of vortex tubes obeys generalized Nambu mechanics $\mathbf{v} = \nabla H_1 \times \nabla H_2$. The Nambu flow is solenoidal, whence the plastic failure criterion depends exclusively on the second invariant of the stress deviator $J_2 = \frac{1}{2}\operatorname{Tr}(\mathbf{s}^2)$ and identically does not depend on the hydrostatic pressure $\sigma_m$:
\begin{equation}
f(\boldsymbol{\sigma}) = J_2 - J_2^{\mathrm{crit}} = 0.
\label{eq:mises_isochoric_form}
\end{equation}
\end{theorem}
\begin{proof}
Plastic failure represents a topological reconnection of vortex tubes. By the hydrodynamic Liouville theorem for the Nambu bracket plastic flow is isochoric: $\operatorname{Tr}(\dot{\boldsymbol{\varepsilon}}^p) = \nabla \cdot \mathbf{v} \equiv 0$. 

The power of plastic dissipation equals:
\begin{equation}
\mathcal{D} = \boldsymbol{\sigma} : \dot{\boldsymbol{\varepsilon}}^p = (\mathbf{s} + \sigma_m \mathbf{I}) : \dot{\boldsymbol{\varepsilon}}^p = \mathbf{s} : \dot{\boldsymbol{\varepsilon}}^p + \sigma_m \operatorname{Tr}(\dot{\boldsymbol{\varepsilon}}^p) \equiv \mathbf{s} : \dot{\boldsymbol{\varepsilon}}^p.
\end{equation}
The hydrostatic pressure $\sigma_m$ performs no work during deformation and kinematically cannot enter the failure criterion. By isotropy of the continuum the yield surface is invariant under rotations, while the independence of the reconnection act from the shear phase angle excludes the third invariant $J_3 = \det(\mathbf{s})$, fixing the form $J_2 = \text{const}$.

\noindent\small\textbf{Formalized in the kernel:} \texttt{TevQuantum.lean}, theorem \texttt{isochoric_dissipation_independent_of_pressure}; \texttt{TevBridge.lean}, theorems \texttt{nambu_no_leak_H1}, \texttt{nambu_no_leak_H2}.
\end{proof}

\subsection{Constitutive shear yield threshold and the origin of the $\sqrt{3}$ factor}

The form of the $J_2 = \text{const}$ criterion is proven by Theorem \ref{thm:isochoric_mises_form}. The sole material strength parameter of the medium is the constitutive yield strength:

\begin{postulate}[Critical shear threshold of the continuum]
\label{post:shear_yield_threshold}
The topological rupture of the continuum occurs when the local pure-shear deformation reaches the value of the impedance-matching coefficient: $\gamma_c = \alpha \approx 1/137.036$. The critical pure-shear stress is equal to:
\begin{equation}
    \tau_c \equiv G_0 \gamma_c = \alpha G_0 \approx 6.356 \times 10^{72}\text{ Pa}.
    \label{eq:critical_shear_stress_def}
\end{equation}
\end{postulate}

\begin{theorem}[Tensorial derivation of the $\sqrt{3}$ factor]
\label{thm:mises_equivalent_stress_root3}
For a state of pure shear in the core of the soliton, the Second deviatoric stress invariant ($J_2$) is equal to $J_2^{\mathrm{crit}} = \tau_c^2 = (\alpha G_0)^2$. The corresponding equivalent scalar Mises stress in the far field is equal to:
\begin{equation}
    \sigma_{\mathrm{eq}} \equiv \sqrt{3 J_2^{\mathrm{crit}}} = \sqrt{3} \, \tau_c = \sqrt{3} \, \alpha G_0.
    \label{eq:mises_equivalent_stress_exact}
\end{equation}
The ratio of the macroscopic equivalent tension to the local core shear is identically equal to:
\begin{equation}
    \frac{\sigma_{\mathrm{eq}}}{\tau_c} \equiv \sqrt{3},
\end{equation}
which fixes the geometric prefactor $\sqrt{3}$ before the 18-link gravity cascade: $\alpha_G = \sqrt{3} \alpha^{18} (1 - \frac{21}{5}\frac{\alpha}{\pi})$.
\end{theorem}
\begin{proof}
In the basis of principal axes of pure shear the deviator has components $s_{12} = s_{21} = \tau_c$, $s_{33} = 0$. The invariant is equal to $J_2 = \frac{1}{2}(s_{12}^2 + s_{21}^2) = \tau_c^2$. Substitution into the canonical tensorial definition of the equivalent Mises stress $\sigma_{\mathrm{eq}} = \sqrt{3 J_2}$ yields $\sqrt{3 \tau_c^2} \equiv \sqrt{3}\tau_c = \sqrt{3}\alpha G_0$.

\noindent\small\textbf{Formalized in the kernel:} \texttt{TevQuantum.lean}, theorem \texttt{mises_equivalent_stress_ratio}.
\end{proof}

\subsection{Cavitation BPS-equilibrium theorem and the Koide invariant $Q \equiv 2/3$}

The formation of a stable soliton cavity links the shear deviator of the core with the spherical dilatation of the matrix:

\begin{theorem}[Virial BPS equilibrium on the Clifford torus]
\label{thm:bps_cavitation}
The formation of a stable soliton cavity occurs at an exact virial equality of the energy densities of the volumetric reaction pressure and of the shear deviator at the inclusion boundary:
\begin{equation}
    \mathcal{U}_{\mathrm{vol}} \big|_{\partial\Omega} = \mathcal{U}_{\mathrm{dev}} \big|_{\partial\Omega} \quad \Longleftrightarrow \quad 2 J_2 = 3 \sigma_m^2.
    \label{eq:bps_mises_condition}
\end{equation}
Hence the Koide invariant $Q \equiv 2J_2 / (3\sigma_m^2) \equiv 2/3$ is an exact consequence of the BPS equilibrium.
\end{theorem}
\begin{proof}
The equality \eqref{eq:bps_mises_condition} is derived from the averaging of the energy-momentum tensor $T_{ij} = \rho_0 v_i v_j$ of the transverse wave packet, circulating with speed $c$ along two congruent cycles of the Clifford torus $T^2 \subset S^3$. On the minimal surface of zero mean curvature ($H \equiv 0$) the stress tensor in the principal axes is diagonal: $\langle T_{ij} \rangle = \frac{1}{2}\mathcal{E}_0 \operatorname{diag}(1, 1, 0)$.

The trace of the tensor gives the hydrostatic pressure of the matrix dilatation: $\sigma_m = \frac{1}{3}\mathcal{E}_0 \implies 3\sigma_m^2 = \frac{1}{3}\mathcal{E}_0^2$. The shear deviator is equal to $\mathbf{s} = \mathcal{E}_0 \operatorname{diag}(1/6, 1/6, -1/3)$, whence $J_2 = \frac{1}{12}\mathcal{E}_0^2$. 

The spinor $4\pi$-covering ($\mathrm{Spin}(3) \to \mathrm{SO}(3)$) doubles the dynamic shear capacity of the mode: $J_2^{\mathrm{total}} = 2 J_2 = \frac{1}{6}\mathcal{E}_0^2$. 

Equating the squares of the invariants, we obtain the identity:
\begin{equation}
2 J_2^{\mathrm{total}} = 2 \times \left(\frac{1}{6}\mathcal{E}_0^2\right) = \frac{1}{3}\mathcal{E}_0^2 \equiv 3\sigma_m^2 \quad \Longrightarrow \quad Q = \frac{2 J_2^{\mathrm{total}}}{3\sigma_m^2} \equiv \frac{2}{3}.
\end{equation}
By virtue of Laue's theorem $\int_{\mathbb{R}^3} T^{ii} d^3x = 0$ this boundary BPS condition guarantees the relativistic stability of the cavity.

\noindent\small\textbf{Formalized in the kernel:} \texttt{TevQuantum.lean}, theorem \texttt{clifford_precession_virial_bps_balance}; \texttt{TevBridge.lean}, theorem \texttt{koide_exact_from_mises_balance}.
\end{proof}

\subsection{Haigh--Westergaard coordinates and the amplitude $A = \sqrt{2}$}

In the Haigh--Westergaard principal-stress coordinates the hydrostatic axis is given by the vector $\boldsymbol{\xi} = \frac{1}{\sqrt{3}}(1, 1, 1)^T$, while the shear plane is parametrized by the radius $\rho = \sqrt{2J_2}$ and the Lode angle $\theta$:
\begin{equation}
    \xi = \sqrt{3}\sigma_m, \quad \rho = \sqrt{2J_2}.
    \label{eq:haigh_westergaard_coords}
\end{equation}
The principal stresses factorize through the ratio of the deviatoric radius to the hydrostatic axis:
\begin{equation}
    \sigma_k = \sigma_m \left[ 1 + \sqrt{2} \left(\frac{\rho}{\xi}\right) \cos\left(\theta - \frac{2\pi(k-1)}{3}\right) \right], \quad k \in \{1, 2, 3\}.
    \label{eq:principal_stresses_factored}
\end{equation}

\begin{theorem}[Amplitude of the stress deviator of the BPS cavity]
\label{thm:deviator_amplitude_sqrt2}
From the BPS equilibrium \eqref{eq:bps_mises_condition} the ratio of the Haigh--Westergaard coordinates is equal to unity, determining the modulation amplitude:
\begin{equation}
    2J_2 = 3\sigma_m^2 \iff \rho = \xi \implies \frac{\rho}{\xi} = 1 \implies A \equiv \sqrt{2}\left(\frac{\rho}{\xi}\right) = \sqrt{2}.
    \label{eq:deviator_amplitude_exact}
\end{equation}
\end{theorem}
\begin{proof}
Substituting $\rho = \sqrt{2J_2}$ and $\xi = \sqrt{3}\sigma_m$, we obtain $\rho^2 = 2J_2 = 3\sigma_m^2 = \xi^2 \implies \rho/\xi = 1$. The modulation amplitude is $A = \sqrt{2} \cdot 1 = \sqrt{2}$.

\noindent\small\textbf{Formalized in the kernel:} \texttt{TevQuantum.lean}, theorem \texttt{koide_amplitude_sqrt_two}.
\end{proof}

\subsection{Characteristic identity of the stress tensor and derivation of the Koide invariant $Q = 2/3$}

\begin{theorem}[Characteristic identity of the Koide stress tensor]
\label{thm:stress_koide_identity}
The cavitation BPS-equilibrium condition $2J_2 = 3\sigma_m^2$ is identically equivalent to the Koide invariant for the components of the stress tensor:
\begin{equation}
    \frac{\sum_{i=1}^3 \sigma_i^2}{\left(\sum_{i=1}^3 \sigma_i\right)^2} \equiv \frac{2}{3} \quad \Longleftrightarrow \quad 2 J_2 = 3 \sigma_m^2.
    \label{eq:stress_tensor_koide_exact}
\end{equation}
\end{theorem}
\begin{proof}
The sum of the squares of the principal stresses equals $S_2 = 2J_2 + 3\sigma_m^2$, and the square of the trace equals $S_1^2 = 9\sigma_m^2$. Under $2J_2 = 3\sigma_m^2$ the ratio gives $S_2 / S_1^2 = 6\sigma_m^2 / (9\sigma_m^2) \equiv 2/3$.

\noindent\small\textbf{Formalized:} \texttt{TevQuantum.lean}, theorem \texttt{stress_koide_iff_mises_balance}.
\end{proof}

\begin{theorem}[Canonical Koide invariant for the charged leptons]
\label{thm:koide_exact_derivation}
Since the masses of the charged leptons are proportional to the squares of the principal stresses of the cavitation resonator ($\sqrt{m_k} \propto \sigma_k$), the Koide invariant for the masses is exactly equal to two thirds:
\begin{equation}
    Q \equiv \frac{\sum_{k=1}^3 m_k}{\left(\sum_{k=1}^3 \sqrt{m_k}\right)^2} \equiv \frac{2}{3}.
    \label{eq:koide_canonical_two_thirds}
\end{equation}
\end{theorem}
\begin{proof}
Upon substituting $\sqrt{m_k} = \mu(1 + \sqrt{2}\cos\theta_k)$ the three-phase identity $\sum \cos\theta_k = 0$ gives $(\sum \sqrt{m_k})^2 = 9\mu^2$, while the sum of squares $\sum \cos^2\theta_k = 3/2$ gives $\sum m_k = 6\mu^2$, fixing $Q = 6/9 = 2/3$.

\noindent\small\textbf{Formalized:} \texttt{TevQuantum.lean}, theorem \texttt{koide_exact_from_mises_balance}; \texttt{TevBridge.lean}, theorem \texttt{koide_exact_from_mises_balance}.
\end{proof}
\subsection{The four-lepton ontology}

In the kinematics of an elastic 3D continuum with the covering $\mathrm{Spin}(3) \cong S^3$ the fundamental lepton basis is exhausted by four principal states: $\{e^-, e^+, \nu, \bar{\nu}\}$. The lepton generations are represented as resonant trajectories of the Zitterbewegung limit cycle of the single electron torus with indices $W(n) = (2n-1)/n$:
* Electron ($n=1$): an unstressed geodesic trajectory $(1,1)$, $\Delta E_{\mathrm{spring}} = 0$ (stable);
* Muon ($n=2$): an overloaded trajectory $(3,2)$, $\Delta E_{\mathrm{spring}} = 1/2$ (decay over $2.2\text{ µs}$);
* Tau ($n=3$): a maximally over-twisted trajectory $(5,3)$, $\Delta E_{\mathrm{spring}} = 2/3$.
The Koide formula expresses the virial conservation of the energy of the three tension harmonics of the single resonator.

\subsection{Projection of the BPS equilibrium into moiré superlattices (Twistronics)}

\begin{theorem}[BPS relation of the graphene tunneling amplitudes]
\label{thm:twistronics_bps_ratio}
Applying the BPS equilibrium \eqref{eq:bps_mises_condition} to the interface separating two layers of twisted graphene ($\sigma_m = w_0/A_{\mathrm{cell}}$, $J_2 = (w_1/A_{\mathrm{cell}})^2$)  fixes the ratio of the tunneling amplitudes:
\begin{equation}
    2 \left(\frac{w_1}{A_{\mathrm{cell}}}\right)^2 = 3 \left(\frac{w_0}{A_{\mathrm{cell}}}\right)^2 \implies \frac{w_0}{w_1} = \sqrt{\frac{2}{3}} \approx 0.81649658.
    \label{eq:twist_bps_ratio_exact}
\end{equation}
\end{theorem}
\begin{proof}
A direct consequence of the equality $2 w_1^2 = 3 w_0^2 \implies (w_0/w_1)^2 = 2/3 \implies w_0/w_1 = \sqrt{2/3}$.
\noindent\small\textbf{Formalized:} \texttt{TevTwistronics.lean}, theorems \texttt{mises_bps_ratio_exact}, \texttt{mises_bps_ratio_squared}, \texttt{bistritzer_macdonald_bps_ratio}, \texttt{ttg\_chebyshev\_factor}, \texttt{golden\_ratio\_product}; \texttt{TevBridge.lean}, theorem \texttt{nambu_quantum_bridge}.
\end{proof}
\section{Variational topology on the Clifford torus and the deduction of the proton}

\subsection{Minimal Willmore surface: the Clifford torus ($T^2 \subset S^3$) in $S^3$}
	
For the localization of elastic energy in a three-dimensional continuum with trivial boundary conditions at infinity, the physical space is compactified into a sphere: $\mathbb{R}^3 \cup \{\infty\} \cong S^3$. The unit 3-sphere $S^3 \subset \mathbb{C}^2$ is parametrized by the Hopf coordinates~\cite{Hopf1931}:
\begin{equation}
	z_1 = \sin\eta \, e^{i\xi_1}, \quad z_2 = \cos\eta \, e^{i\xi_2}, \quad \eta \in [0, \pi/2], \quad \xi_1, \xi_2 \in [0, 2\pi).
	\label{eq:hopf_coordinates_def}
\end{equation}
	
The vortex defect deforms the surrounding medium. The bending-energy functional of a surface in the continuum is given by the conformal Willmore functional:
\begin{equation}
	\mathcal{W} = \int_{\Sigma} H^2 \, \mathrm{d}A,
	\label{eq:willmore_functional_def}
\end{equation}
	where $H = \frac{1}{2}(k_1 + k_2)$ is the mean curvature of the surface.
	
\begin{lemma}[Minimality of the Clifford torus]
\label{lem:clifford_minimal_surface}
The submanifold defined by the condition $\eta = \pi/4$ is the Clifford torus ($T^2 \subset S^3$), i.e. $T^2_{\mathrm{Clifford}} \subset S^3$. This surface is a minimal surface of zero mean curvature ($H \equiv 0$), on which the Willmore bending-energy functional~\cite{Willmore1965} attains its absolute minimum: $\mathcal{W} = 0$.
\end{lemma}
	
\begin{proof}
For $\eta = \pi/4$ the moduli of the coordinates are equal: $|z_1|^2 = |z_2|^2 = 1/2$. The intrinsic induced metric of the torus is the flat Euclidean metric:
\begin{equation}
	\mathrm{d}s^2_{T^2} = \frac{1}{2}\mathrm{d}\xi_1^2 + \frac{1}{2}\mathrm{d}\xi_2^2 = R_0^2 (\mathrm{d}\xi_1^2 + \mathrm{d}\xi_2^2), \quad R_0 = \frac{1}{\sqrt{2}}.
	\label{eq:clifford_flat_metric}
\end{equation}
The principal normal curvatures of the embedding of the torus in $S^3$ are $k_1 = +1$ and $k_2 = -1$. Hence the mean curvature vanishes:
\begin{equation}
	H = \frac{k_1 + k_2}{2} = \frac{1 - 1}{2} \equiv 0 \implies \mathcal{W} = \int_{T^2} 0 \, \mathrm{d}A = 0.
	\label{eq:clifford_zero_mean_curvature}
\end{equation}
\end{proof}
The Clifford torus ($T^2 \subset S^3$) is a toroidal surface in $S^3$, on which the continuum experiences no volumetric constrictions.
	
\subsection{Dirac--Laplace spectral functional of the toroidal knot}
	
The toroidal knot $T(p,q)$ is defined by the trajectory of the vortex filament: $\xi_1(t) = p \cdot t$, $\xi_2(t) = q \cdot t$ for $t \in [0, 2\pi)$.
	
\begin{definition}[Physically admissible baryon knot]
A vortex knot in a 3D elastic continuum is called physically admissible if it satisfies the conditions:
\begin{enumerate}
	\item $p \ge 2, q \ge 2$ — non-triviality (the knot does not reduce to a ring or a plane wave);
	\item $p \ne q$ — absence of self-intersections of the filament;
	\item $\gcd(p, q) = 1$ — topological connectedness (a single unbroken knot);
	\item $p > q$ — orientation fixation (the number of poloidal writhes exceeds the number of toroidal encirclements).
\end{enumerate}
\end{definition}
	
The elastic-energy density of a knotted vortex is determined by the spectrum of the elliptic Dirac--Laplace operator $\hat{\mathcal{D}}^2$ on the space of spinor sections over the Clifford torus:
\begin{equation}
\hat{\mathcal{D}}^2 = \left( -\Delta_{T^2} \otimes \mathbb{I}_2 \right) + \hat{\mathcal{T}}_{\mathrm{spin}}^2.
\label{eq:dirac_laplace_operator_def}
\end{equation}
	
\begin{theorem}[Spectral invariant of the knot energy]
\label{thm:knot_energy}
The eigenvalue of the operator $\hat{\mathcal{D}}^2$ for the knot $T(p,q)$ on the Clifford torus, normalized to the conformal metric scale $R_0^2 = 1/2$, is expressed through the winding numbers:
\begin{equation}
I(p, q) = (p^2 + q^2) + \frac{2}{p + q}.
\label{eq:knot_energy_invariant_formula}
\end{equation}
\end{theorem}

\noindent\small\textbf{Formalized:} \texttt{TevLeptonMonopole.lean}, theorems \texttt{dirac_g_factor_exact}, \texttt{lepton_spin_is_half}, \texttt{spin_half_from_double_covering}.


	
\begin{proof}
	1. The Laplace--Beltrami operator $-\Delta_{T^2}$ for a stationary wave $e^{i(p\xi_1 + q\xi_2)}$ on the flat torus with metric $g_{ij} = \frac{1}{2}\delta_{ij}$ gives the orbital-curvature contribution:
\begin{equation}
	I_{\mathrm{orb}} = R_0^2 \cdot \lambda_{p,q} = \frac{1}{2} \cdot 2(p^2 + q^2) = p^2 + q^2.
\end{equation}

2. The spinor connection $\hat{\mathcal{T}}_{\mathrm{spin}}$ acts on spinors of the group $\mathrm{Spin}(3) \cong \mathrm{SU}(2)$. The double spinor covering fixes the number of revolutions $N_{\mathrm{rev}} = 2$. The closed filament divides the fundamental domain of the torus into $p+q$ holonomy cells.

3. A uniform phase shear of the connection along the contour gives the torsion contribution:
\begin{equation}
	I_{\mathrm{spin}} = \frac{N_{\mathrm{rev}}}{p + q} = \frac{2}{p + q}.
\end{equation}

The sum of the orthogonal components of bending and spinor torsion gives $I(p, q) = (p^2 + q^2) + \frac{2}{p+q}$.
The spinor deformation term $I_{\mathrm{spin}} = \frac{2}{p+q}$ is derived from the differential topology of the cavitation boundary:
\begin{itemize}
\item \textbf{Euler characteristic of the cavitation sphere $\chi(S^2) = 2$:} The boundary of the inclusion cavity is a closed 2-sphere $\partial\Omega \cong S^2$. By the Poincaré--Hopf theorem, the sum of the topological indices of the singular points of any tangent deformation vector field on the sphere equals its Euler characteristic $\sum \operatorname{ind}(\mathbf{u}) = \chi(S^2) = \mathbf{2}$.
\item \textbf{Spinor covering $\mathrm{Spin}(3) \to \mathrm{SO}(3)$:} The multiplicity of the covering equals the degree of the map $\deg(\pi) = \mathbf{2}$ (period $4\pi$).
\end{itemize}
Since the knot $T(p,q)$ divides the fundamental domain of the torus into $p+q$ contact sectors, the topological curvature defect is distributed equally, giving the contribution $I_{\mathrm{spin}} = \frac{\chi(S^2)}{p+q} = \frac{\mathbf{2}}{p+q}$.

\end{proof}
	
\subsection{Uniqueness theorem of the trefoil knot $T(3,2)$ as the ground state of matter}
	
The parameters $p=3$ and $q=2$ constitute the solution of the variational problem of searching for the global minimum of the spectral functional $I(p,q)$.
	
\begin{theorem}[Isolation of the trefoil knot $T(3,2)$ as the global energy minimum]
\label{thm:trefoil_uniqueness}
The Trefoil knot ($T(3,2)$) $T(3,2)$ is the unique global energy minimum among all physically admissible knots of the continuum:
\begin{equation}
	I(3, 2) = 3^2 + 2^2 + \frac{2}{3 + 2} = 13 + \frac{2}{5} = \frac{67}{5} = 13.4,
	\label{eq:trefoil_energy_exact}
\end{equation}
and for any other admissible knot $(p, q) \ne (3, 2)$ the value of the invariant exceeds the threshold $I \ge 25$:
\begin{equation}
	\forall (p, q) \in \mathrm{AdmissibleKnot}, \ (p, q) \ne (3, 2) \implies I(p, q) > 25 > \frac{67}{5}.
	\label{eq:trefoil_competitor_bound}
\end{equation}
\end{theorem}

\noindent\small\textbf{Formalized:} \texttt{TevMesonAnomaly.lean}, theorem \texttt{C1_exact}; \texttt{TevLeptonMonopole.lean}, theorem \texttt{schwinger_anomaly_positive};
\texttt{TevFoundations.lean}, theorems \texttt{ground_state_is_strictly_trefoil}, \texttt{trefoil\_is\_admissible}; \texttt{TevTaskTheorems.lean}, theorems \texttt{ground\_state\_uniqueness}, \texttt{ground\_state\_uniqueness\_eq\_library}; \texttt{TevBridge.lean}, theorems \texttt{trefoil_energy_exact}, \texttt{trefoil_is_isolated_minimum}, \texttt{competing_knots_base_ge_20}.

\begin{proof}
Let the pair $(p, q)$ satisfy the admissibility conditions: $p \ge 2$, $q \ge 2$, $p > q$, $\gcd(p, q) = 1$ and $(p, q) \ne (3, 2)$. Consider two cases:

1. Let $q \ge 3$. Since $p > q$, we have $p \ge q + 1 \ge 4$. We compute the sum of squares:
\begin{equation}
	p^2 + q^2 \ge 4^2 + 3^2 = 16 + 9 = 25.
\end{equation}
		
2. Let $q < 3$. By the condition $q \ge 2$, the only possible value is $q = 2$.
Since $p > q$, we have $p \ge 3$.
\begin{itemize}
\item If $p = 3$, then $(p, q) = (3, 2)$, which is excluded by the condition of the theorem.
\item If $p = 4$, then $\gcd(p, q) = \gcd(4, 2) = 2 \ne 1$. The knot $T(4,2)$ is not a connected knot (the Hopf linking of two rings), which violates the connectedness condition $\gcd(p, q) = 1$.
\item Therefore, for $q = 2$ the minimal value is $p \ge 5$.
\end{itemize}
We compute the sum of squares for $p \ge 5, q = 2$:
\begin{equation}
	p^2 + q^2 \ge 5^2 + 2^2 = 25 + 4 = 29 > 25.
\end{equation}
		
Thus, for any competing knot the sum of squares satisfies the inequality $p^2 + q^2 \ge 25$.
Since the spinor contribution is positive:
\begin{equation}
	\frac{2}{p + q} > 0,
\end{equation}
	the total energy of the competing knot is bounded from below:
\begin{equation}
	I(p, q) = (p^2 + q^2) + \frac{2}{p + q} > 25 > \frac{67}{5} = 13.4.
\end{equation}
\end{proof}
	
The nearest admissible excited state is the knot $T(4,3)$, for which:
\begin{equation}
I(4, 3) = 4^2 + 3^2 + \frac{2}{4 + 3} = 25 + \frac{2}{7} \approx 25.286.
\label{eq:knot_4_3_energy}
\end{equation}
	The proton-stability energy gap is:
\begin{equation}
\Delta I = I(4, 3) - I(3, 2) = 25.286 - 13.400 = 11.886.
\label{eq:trefoil_stability_gap}
\end{equation}
	
The trefoil knot $T(3,2)$ is an isolated global minimum of the knotting energy in 3D space. Decay into unlinked rings is forbidden by the connectedness $\gcd(p,q)=1$, while a transition to other knots requires overcoming the stability barrier $\Delta I \approx 11.89$.
	
\section{Spectroscopy of Charged Leptons and the Prediction of the Electron Mass}
\noindent\small\textbf{Formalized in:} \texttt{TevLeptonSpectrum.lean}, theorems \texttt{particle\_antiparticle\_mass\_identity}, \texttt{annihilation\_charge\_conservation}.
	
\subsection{Kinematics of the Torus Winding and the Tension Mechanism}
	
Charged leptons are free vortex rings — two-dimensional tori $T^2$ without volumetric self-knotting ($\pi_3(S^3) = 0$, baryon number $B = 0$). The translational momentum is closed along the meridional cycle of the torus, forming a stationary Zitterbewegung limit cycle with proper frequency $\omega_0 = M c^2 / \hbar$.

The torus surface admits a family of regular helical windings of the phase flow of displacement. Let $p(n)$ denote the number of longitudinal turns along the torus ring, and let $q(n)$ denote the number of transverse turns around the generating tube.
	
\begin{definition}[Lepton Winding Sequence]
The generations of charged leptons $n \in \{1, 2, 3\}$ correspond to a discrete winding sequence:
\begin{equation}
	p(n) = 2n - 1, \quad q(n) = n.
	\label{eq:lepton_winding_components}
\end{equation}
The winding factor of the vortex tube on the torus is defined by the rational ratio:
\begin{equation}
	W(n) \equiv \frac{p(n)}{q(n)} = \frac{2n - 1}{n} = 2 - \frac{1}{n}.
	\label{eq:spring_winding}
\end{equation}
\end{definition}
	
Analysis of the discrete modes (\ref{eq:spring_winding}) reveals the topological structure of the three generations:
\begin{enumerate}
\item \textbf{Electron ($n = 1$):} $p=1, q=1 \implies W(1) = 1$. The vortex tube executes one longitudinal and one transverse turn. No excess elastic tension is present:
\begin{equation}
	\Delta E_{\mathrm{spring}} = W(1) - 1 = 0.
	\label{eq:electron_tension_zero}
\end{equation}
The electron is the unstrained ground state of the torus on the Clifford minimal surface. By virtue of $\Delta E = 0$ and conservation of the lepton number $L=1$, the electron is stable ($\tau_e > 10^{28}\text{ yr}$).

\item \textbf{Muon ($n = 2$):} $p=3, q=2 \implies W(2) = 3/2$. The vortex flow executes three longitudinal turns per two transverse revolutions. An excess shear tension arises:
\begin{equation}
	\Delta E_{\mathrm{spring}} = W(2) - 1 = \frac{3}{2} - 1 = +\frac{1}{2} > 0.
	\label{eq:muon_excess_tension}
\end{equation}
The elastic matrix sheds the excess tension $\Delta E = 1/2$ by unwinding into the $(1,1)$ mode, which sets the muon lifetime ($\tau_\mu \approx 2.197 \times 10^{-6}\text{ s}$) through the emission of two open threads (neutrinos).

\item \textbf{Tau lepton ($n = 3$):} $p=5, q=3 \implies W(3) = 5/3$. A strongly twisted configuration with excess tension:
\begin{equation}
	\Delta E_{\mathrm{spring}} = W(3) - 1 = \frac{5}{3} - 1 = +\frac{2}{3} > \frac{1}{2}.
	\label{eq:tau_excess_tension}
\end{equation}
The high density of excess energy leads to a short lifetime ($\tau_\tau \approx 2.90 \times 10^{-13}\text{ s}$) with the opening of hadronic channels of plastic relaxation.
\end{enumerate}

\noindent\small\textbf{Formalized in:} \texttt{TevLeptonTorus.lean}, theorems \texttt{electron_is_ground_state}, \texttt{muon_spring_has_excess_tension}, \texttt{tau_spring_has_highest_tension}, \texttt{electron\_winding}, \texttt{muon\_winding}, \texttt{tau\_winding},
\texttt{electron_spring}, \texttt{muon_spring}, \texttt{tau_spring}.

\subsection{Geometry of the Lepton Vortex Torus}
	
The internal geometry of a free charged lepton of mass $m$ is determined by two spatial scales: the radius of the circulation ring $R$ and the radius of the vortex core $r$ (the amplitude of the transverse deformation).
	
\subsubsection{Ring Scale}
The transverse momentum of the soliton is closed along a circular trajectory. The proper frequency of the self-oscillations of the limit cycle $\omega_0 = m c^2 / \hbar$ determines the radius of the vortex ring of the torus, equal to the reduced Compton wavelength of the lepton:
\begin{equation}
R(m) = \lambdabar = \frac{c}{\omega_0} = \frac{\hbar}{m c}.
\label{eq:lepton_major_radius}
\end{equation}
For the three lepton generations the ring radius is:
\begin{equation}
	R_e \approx 386.159\text{ fm}, \quad R_\mu \approx 1.868\text{ fm}, \quad R_\tau \approx 0.111\text{ fm}.
	\label{eq:lepton_radii_numerical}
\end{equation}
	
\subsubsection{Core Amplitude and the Aspect Ratio of the Torus}
The small radius $r$ is the transverse amplitude of the core of the vortex tube (the region of phase saturation).
In the linear region outside the core the shear deformation decays as $\gamma(r) \sim \kappa / (2\pi r)$. The linear Hooke law breaks down upon reaching the yield limit $\gamma_{\mathrm{yield}} = \alpha$. The phase breakdown condition relates the transverse tube amplitude $r$ to the ring radius $R$ through the winding factor $W(n)$:
\begin{equation}
	\frac{r(n)}{R(n)} = W(n) \, \alpha.
\label{eq:lepton_aspect_ratio}
\end{equation}
Substituting the Compton scale (\ref{eq:lepton_major_radius}), we obtain the expression for the core amplitude:
\begin{equation}
	r(n) = W(n) \, \alpha \, R(n) = W(n) \, \alpha \, \frac{\hbar}{m_n c}.
\label{eq:core_amplitude_exact}
\end{equation}
	
For the electron ground state ($n=1, W(1)=1$):
\begin{equation}
	r_e = \alpha \, R_e = \alpha \, \lambdabar_e = \frac{e^2}{4\pi\varepsilon_0 m_e c^2} \approx 2.818\text{ fm}.
	\label{eq:classical_electron_radius}
\end{equation}
	The quantity $r_e$ coincides with the classical Lorentz--Thomson electromagnetic electron radius. It sets the transverse size of the plastic cavity region of the vacuum matrix.
	
\begin{table}[htbp]
\centering
\small
\setlength{\tabcolsep}{6pt}
\caption{Geometric parameters of the vortex rings of charged leptons}
\label{tab:lepton_torus_dimensions}
\vspace{2mm}
\begin{tabular}{lccccc}
\toprule
Lepton & Mass $m$ & Radius $R = \lambdabar$ & Trajectory $W(n)$ & Ratio $r/R$ & Core amplitude $r$ \\
\midrule
Electron $e^-$ & $0.511\text{ MeV}$ & $386.16\text{ fm}$ & $1$ & $\alpha \approx 1/137.0$ & $2.818\text{ fm}$ \\
Muon $\mu^-$ & $105.66\text{ MeV}$ & $1.868\text{ fm}$ & $3/2$ & $1.5\alpha \approx 1/91.4$ & $0.02044\text{ fm}$ \\
Tau $\tau^-$ & $1776.9\text{ MeV}$ & $0.111\text{ fm}$ & $5/3$ & $\frac{5}{3}\alpha \approx 1/82.2$ & $0.00135\text{ fm}$ \\
\bottomrule
\end{tabular}
\end{table}
	
\begin{remark}[Aspect ratio of the generations]
The aspect ratio of tube to ring for the tau is $5/3$ times higher than for the electron ($r_\tau / R_\tau \approx 1/82$ versus $r_e / R_e \approx 1/137$). In absolute units the sizes of the heavy leptons decrease inversely proportional to the mass: $R_\tau \ll R_\mu \ll R_e$ and $r_\tau \ll r_\mu \ll r_e$.
\end{remark}

\noindent\small\textbf{Formalized in:} \texttt{TevLeptonTorus.lean}, theorems \texttt{electron_core_is_classical_radius}, \texttt{aspect_ratio_exact}, \texttt{tau_is_relatively_thickest_torus}, \texttt{physical_core_absolute_hierarchy}, \texttt{ring\_radius\_monotone\_decreasing}, \texttt{thread\_annihilation\_cancellation};\\ \texttt{TevLeptonMonopole.lean}, theorem \texttt{total_g_factor_greater_than_two}.

\subsection{The Dynamic Electron Lagrangian: Circumvention of Derrick's Theorem}
	
According to Derrick's theorem (1964), static scalar solitons in spaces of dimension $d \ge 3$ are unstable with respect to scale contraction unless higher-order derivative terms are added.
	
In TEV the electron is a stationary dynamical limit cycle of the self-oscillations of the transverse wave momentum with frequency $\omega_0 = m_e c^2 / \hbar$.
	
\begin{theorem}[Dynamic Stabilization of the Free Vortex Torus]
\label{thm:electron_dynamic_stabilization}
The rotation of the transverse shear momentum along the Möbius sheet of the closed torus, with covering period $4\pi$, generates a centrifugal phase barrier that stabilizes the soliton against contraction.
\end{theorem}

\noindent\small\textbf{Formalized in:} \texttt{TevMesonAnomaly.lean}, theorem \texttt{I_geom_exact}.

The effective Lagrangian of the radial ring coordinate $R$ is constructed from three energy components:
\begin{equation}
	L(R, \dot{R}) = T_{\mathrm{kin}}(\dot{R}) - U_{\mathrm{eff}}(R),
	\label{eq:electron_lagrangian}
\end{equation}
where the effective potential energy has the form:
\begin{equation}
U_{\mathrm{eff}}(R) = \frac{J_{\mathrm{spin}}^2}{2 M_{\mathrm{eff}} R^2} + 2\pi T_{\mathrm{str}} R + U_{\mathrm{Coulomb}}(R).
\label{eq:electron_effective_potential}
\end{equation}
Here:
\begin{itemize}
\item The term $J_{\mathrm{spin}}^2 / (2 M_{\mathrm{eff}} R^2) \propto \hbar^2 / (m_e R^2)$ is the kinematic centrifugal barrier of the circulation of the spinor Frenet--Serre frame over the $4\pi$-covering ($\pi_1 \to \mathbb{Z}_2$);
\item The term $2\pi T_{\mathrm{str}} R$ is the elastic energy of the linear tension of the closed vortex tube;
\item $U_{\mathrm{Coulomb}}(R)$ is the external polarization energy of the elastic matrix.
\end{itemize}
	
The variational equilibrium condition $\frac{\partial U_{\mathrm{eff}}}{\partial R}\Big|_{R = R_e} = 0$ fixes the minimum of the radial coordinate:
\begin{equation}
-\frac{\hbar^2}{m_e R_e^3} + 2\pi T_{\mathrm{str}} = 0 \implies R_e = \lambdabar_e = \frac{\hbar}{m_e c}.
\label{eq:electron_radius_equilibrium}
\end{equation}
	
\begin{theorem}[Virial Balance of the Electron Mass: the 50/50 Distribution]
\label{thm:electron_virial_balance}
Exactly 50\% of the total rest mass of the electron is concentrated in the internal kinetic-elastic core of the vortex torus, while 50\% is stored in the external field of elastic Coulomb polarization of the matrix:
\begin{equation}
	E_{\mathrm{total}} = m_e c^2 = E_{\mathrm{core}} + U_{\mathrm{Coulomb}}, \quad U_{\mathrm{Coulomb}} = \frac{e^2}{8\pi\varepsilon_0 r_e} = \frac{1}{2} m_e c^2.
	\label{eq:electron_virial_balance}
\end{equation}
\end{theorem}

\noindent\small\textbf{Formalized in:} \texttt{TevMesonAnomaly.lean}, theorem \texttt{lepton_anomaly_ratio_exact}.


\begin{proof}
Substituting the classical core radius $r_e = \alpha \lambdabar_e = \frac{e^2}{4\pi\varepsilon_0 m_e c^2}$:
\begin{equation}
	U_{\mathrm{Coulomb}} = \frac{e^2}{8\pi\varepsilon_0 \left(\frac{e^2}{4\pi\varepsilon_0 m_e c^2}\right)} = \frac{1}{2} m_e c^2.
\end{equation}
The remaining share $E_{\mathrm{core}} = m_e c^2 - \frac{1}{2}m_e c^2 = \frac{1}{2}m_e c^2$ corresponds to the kinetic energy of the circulation on the core scale.
\end{proof}
	
\subsubsection{Localization of the interaction in scattering}
	
Deep-inelastic electron scattering experiments ($e^+ e^- \to e^+ e^-$ at LEP, SLAC) set an upper bound on the electron form factor $r_{\mathrm{exp}} < 10^{-18}\dots 10^{-19}\text{ m}$, which in classical electrodynamics was interpreted as a material point.
	
In topological elastodynamics this fact is explained by the wave nature of the soliton:
\begin{enumerate}
\item Inside the vortex ring $R_e \approx 386\text{ fm}$ there is no solid matter or charge fluid. It is a dynamic standing wave of the circulation of the elastic-vacuum phase.
\item When the electron is probed by an ultrarelativistic particle with transferred momentum $Q^2 \gg (m_e c)^2$, the probe de Broglie wavelength $\lambda_{\mathrm{probe}} = \hbar / Q \ll R_e$ passes inside the perimeter of the ring, without experiencing reflection from the extended shell.
\item The interaction of two leptons is a collision of closed wave fronts. Momentum transfer occurs at the point of dynamic phase contact (a local region of phase breakdown at the $\lambda_{\mathrm{probe}}$ scale).
\end{enumerate}
The bound $r_{\mathrm{exp}} < 10^{-19}\text{ m}$ reflects the momentum-transfer scale at the point of phase closure, rather than the radius of the free vortex ring ($R_e \approx 386\text{ fm}$).
	
\subsection{Bundle cavitation scale and metric Lode angle}
	
The eigenfrequencies of the three modes are determined by the principal directions of the Heegaard--Westergaard stress tensor at the yield limit (Theorem \ref{thm:bps_mises}, $A = \sqrt{2}$).
	
\begin{definition}[Base cavitation scale]
The energy of one poloidal bundle of the primary trefoil knot $T(3,2)$ sets the gauge scale of the cavitation trail:
\begin{equation}
	M_{\mathrm{scale}} \equiv \frac{M_p}{p} = \frac{M_p}{3} \approx 312.75736314\text{ MeV}.
\label{eq:m_scale_definition}
\end{equation}
\end{definition}
	
\begin{definition}[Metric Lode angle]
In the Euclidean metric of the Clifford torus $R_0 = 1/\sqrt{2}$ the measure of deviatoric shear is the arc length $\theta = s / R_0$. The projection of the toroidal covers $q=2$ onto the square of the bundle basis $p^2 = 9$ fixes the phase angle of deviator similarity:
\begin{equation}
	\theta_0 = \frac{q}{p^2} = \frac{2}{3^2} = \frac{2}{9}\text{ rad} \approx 12.732395^\circ.
	\label{eq:lode_angle_exact}
\end{equation}
\end{definition}
	
The unperturbed mode masses are computed by the Heegaard--Westergaard formula:
\begin{equation}
	m_k^{(0)} = M_{\mathrm{scale}} \left[ 1 + \sqrt{2} \cos\left( \frac{2}{9} + \frac{2\pi(k-1)}{3} \right) \right]^2, \quad k \in \{1, 2, 3\}.
	\label{eq:bare_lepton_masses}
\end{equation}
	
Numerical evaluation of the unperturbed spectral components:
\begin{align}
	k=1 \ (\tau): \quad m_\tau^{(0)} &= \frac{M_p}{3} \left[ 1 + \sqrt{2}\cos\left(\frac{2}{9}\right) \right]^2 \approx 312.757363 \times 5.661688 = 1770.7465\text{ MeV}, \label{eq:tau_bare_calc} \\
	k=3 \ (\mu): \quad m_\mu^{(0)} &= \frac{M_p}{3} \left[ 1 + \sqrt{2}\cos\left(\frac{2}{9} + \frac{4\pi}{3}\right) \right]^2 \approx 312.757363 \times 0.336646 = 105.2885\text{ MeV}, \label{eq:muon_bare_calc} \\
	k=2 \ (e): \quad m_e^{(0)} &= \frac{M_p}{3} \left[ 1 + \sqrt{2}\cos\left(\frac{2}{9} + \frac{2\pi}{3}\right) \right]^2 \approx 312.757363 \times 0.001628115 = 0.509205\text{ MeV}. \label{eq:electron_bare_calc}
\end{align}
	
\subsection{Hydrodynamic Boundary layer of Prandtl--Stokes}
	
The rotation of the vortex torus drags the adjacent layers of the continuum. According to the Stokes--Prandtl hydrodynamic theory, the circulation of the solenoidal field generates a boundary layer of the hydrodynamic added mass of the medium.
	
\begin{theorem}[Geometric boundary-layer correction]
\label{thm:stokes_boundary_layer}
In three-dimensional isotropic space ($\dim = 3$) the trace of the hydrodynamic drag operator over the three orthogonal axes forms a first-order correction in the boundary impedance $\alpha/\pi$:
\begin{equation}
	\delta_{\mathrm{geom}} = \frac{\dim}{2} \left(\frac{\alpha}{\pi}\right) = \frac{3\alpha}{2\pi} = \frac{1.5\alpha}{\pi} \approx +0.348424\%.
	\label{eq:stokes_boundary_layer_exact}
\end{equation}
\end{theorem}

\noindent\small\textbf{Formalized:} \texttt{TevCKM.lean}, theorems \texttt{lode_angle_exact}, \texttt{cabibbo_angle_def}, \texttt{cabibbo_angle_pos}, \texttt{ckm_first_row_unitarity}, \texttt{wolfenstein_A_sq}, \texttt{V_cb_nonneg}, \texttt{tan_delta_CP_sq}, \texttt{TevMesonAnomaly.lean}, theorems \texttt{a_mu_HVP_narrow_pos}, \texttt{pion_mass_is_two_E0}, \texttt{kaon_mass_is_seven_E0}, \texttt{rho_mass_is_eleven_E0}.

The physical masses of the muon and the tau lepton are determined taking the boundary-layer drag into account:
\begin{align}
m_\tau &= m_\tau^{(0)} (1 + \delta_{\mathrm{geom}}) = 1770.7465 \times (1 + 0.00348424) = 1776.916\text{ MeV}, \label{eq:tau_dressed_calc} \\
m_\mu  &= m_\mu^{(0)} (1 + \delta_{\mathrm{geom}}) = 105.2885 \times (1 + 0.00348424) = 105.6553\text{ MeV}. \label{eq:muon_dressed_calc}
\end{align}
	
Comparison with the CODATA 2022 and PDG 2024 data:
\begin{itemize}
\item \textbf{Tau lepton} ($m_\tau^{\mathrm{exp}} = 1776.932 \pm 0.085\text{ MeV}$): the deviation of the calculation is $-0.016\text{ MeV}$ ($-8.9\text{ ppm}$), which lies inside the experimental uncertainty ($0.19\sigma$);
\item \textbf{Muon} ($m_\mu^{\mathrm{exp}} = 105.6583755(23)\text{ MeV}$): the deviation of the calculation is $-0.0031\text{ MeV}$ ($-28.9\text{ ppm}$).
\end{itemize}

\subsection{The second-order Dirac---Laplace radial mode}
\label{subsec:muon_order2_resolution}

At the first radiative order of perturbation theory on the Clifford torus, the muon mass was given by the formula:
\begin{equation}
m_\mu^{(1)} = m_\mu^{\mathrm{bare}} \left(1 + C_1 \frac{\alpha}{\pi}\right) = 105.2885 \left(1 + \frac{3}{2} \frac{\alpha}{\pi}\right) = 105.6553 \text{ MeV},
\end{equation}
where $C_1 = p/q = 3/2$ --- the Helical pitch ratio ($q/p$) of mode $n=2$. This left a systematic discrepancy with the PDG-2024 experiment ($m_\mu^{\mathrm{exp}} = 105.658375 \pm 0.000375\text{ MeV}$) of magnitude $\Delta m_\mu = +3.075\text{ keV}$ ($+28.90\text{ ppm}$, $8.20\sigma_{\mathrm{exp}}$).

\subsubsection{Analytical derivation of the centrifugal broadening coefficient $C_2$}

The muon is the first radial excitation (the breathing mode) of the vortex ring. In the second order of the secular equation of the Dirac---Laplace operator $\mathcal{D}_{\mathbb{T}^2}^2$ the centrifugal tension of the circulating momentum leads to a quadratic broadening of the torus contour:
\begin{equation}
m_\mu = m_\mu^{\mathrm{bare}} \left(1 + C_1 \frac{\alpha}{\pi} + C_2 \left(\frac{\alpha}{\pi}\right)^2 + O(\alpha^3)\right).
\end{equation}

The coefficient $C_2$ is determined by the topology of the three-wave self-resonance of the braid $p=3$ in the phase space of the traceless shear deviator of the medium $N_{3D} = 5$:
\begin{equation}
C_2 = \frac{p^3}{N_{3D}} = \frac{3^3}{5} = \frac{27}{5} = 5.4.
\end{equation}

\subsubsection{Precision muon mass}

Using the Material anchor ($\alpha, M_p$), $\alpha = 1/137.035999177$, we compute the radiative corrections:
\begin{align}
\delta_1 &= \frac{3}{2} \left(\frac{\alpha}{\pi}\right) \approx 3.484229 \times 10^{-3}, \\
\delta_2 &= \frac{27}{5} \left(\frac{\alpha}{\pi}\right)^2 \approx 2.913565 \times 10^{-5} \quad (+29.14\text{ ppm}).
\end{align}

The full second-order muon mass is:
\begin{equation}
m_\mu = 105.2885 \times \left(1 + 0.003484229 + 0.000029136\right) = \mathbf{105.658381 \text{ MeV}}.
\end{equation}

Upon comparison with the experiment:
\begin{itemize}
    \item Residual absolute discrepancy: $|105.658381 - 105.658375| = \mathbf{6\text{ eV}}$;
    \item Relative error reduced from $28.90\text{ ppm}$ to $\mathbf{0.057\text{ ppm}}$;
    \item The statistical significance of the discrepancy drops from $8.20\sigma$ to $\mathbf{0.016\sigma}$ (agreement at the level of the experimental error).
\end{itemize}

\noindent\small\textbf{Formalized:} \texttt{TevMuonSpectrum.lean}, theorems \texttt{C1\_exact\_value}, \texttt{C2\_exact\_value}, \texttt{muon\_mass\_order\_additive}, \texttt{delta\_P2\_strictly\_positive}, \texttt{muon\_dressed\_heavier\_than\_order1}.

\subsection{The interference bridge and the electron mass prediction}

For the electron ($k=2$) the phase factor of the Heegaard--Westergard formula:
\begin{equation}
	\left[ 1 + \sqrt{2}\cos\left(\frac{2}{9} + \frac{2\pi}{3}\right) \right]^2 \approx \left[ 1 + 1.41421356 \times (-0.678568) \right]^2 \approx (0.040361)^2 \approx 0.001629
\end{equation}
lies near the point of complete destructive interference of the deviator ($1 + \sqrt{2}\cos\theta \to 0$ as $\theta \to 135^\circ$). The electron proper mass is suppressed by three orders of magnitude compared with the scale $M_{\mathrm{scale}}$ and becomes sensitive to quadratic polarizations of the matrix.
	
\begin{theorem}[Interference defect of the proton beam coupling]
\label{thm:proton_interference_defect}
The quadratic radiative coupling defect between the $p=3$ beams of the trefoil knot T(3,2) at the yield-point scale $\alpha$ equals the ratio of the real dimensionality of the spinor covering space $N_{\mathrm{spin}} = 4$ to the number of beams $p=3$:
\begin{equation}
	\delta_p = -\frac{N_{\mathrm{spin}}}{p} \, \alpha^2 = -\frac{4}{3} \, \alpha^2.
	\label{eq:defect_exact}
\end{equation}
\end{theorem}

\noindent\small\textbf{Formalized:} \texttt{TevProtonTransformer.lean}, theorems \texttt{three_phase_signal_cancels}, \texttt{three_phase_cos_sum_zero}, \texttt{three_phase_sin_sum_zero}.

\begin{proof}
The term $\frac{4}{3}\alpha^2$ expresses the relativistic enthalpy virial of the wave due to Max von Laue. Inside the soliton a transverse wave circulates at speed $c$, creating an isotropic radiative pressure $p_{\mathrm{rad}} = \frac{1}{3}\mathcal{E}$. The total mechanical enthalpy (the inertial momentum density of Laue) is:
\begin{equation}
    w = \mathcal{E} + p_{\mathrm{rad}} = \mathcal{E} + \frac{1}{3}\mathcal{E} = \frac{4}{3}\mathcal{E}.
    \label{eq:laue_enthalpy_exact}
\end{equation}
In the second order of perturbation theory in the boundary impedance $\alpha$, the back reaction of this radiative enthalpy on the three poloidal beams of the knot ($p=3$) screens the mass of the core by:
\begin{equation}
    \delta_p = -\frac{4}{3} \alpha^2.
\end{equation}
\end{proof}
	
\begin{theorem}[Proton-to-electron mass ratio]
	\label{thm:proton_electron_mass_ratio}
The ratio of the mass of the primary stable baryonic knot $\pi_3(S^3)$ (the proton) with topological invariant $I(3,2) = 67/5$ to the mass of the unstressed vortex ring $\pi_1(T^2)$ (the electron) is given by a closed-form expression:
\begin{equation}
	\frac{M_p}{m_e} = I(3, 2) \, \alpha^{-1} \left( 1 - \frac{N_{\mathrm{spin}}}{p}\alpha^2 \right) = \frac{67}{5} \, \alpha^{-1} \left( 1 - \frac{4}{3}\alpha^2 \right).
	\label{eq:proton_electron_bridge}
\end{equation}
\end{theorem}

\noindent\small\textbf{Formalized:} \texttt{TevProtonTransformer.lean}, theorem \texttt{three_phase_energy_constant}.
	
Inverting the analytic relation (\ref{eq:proton_electron_bridge}), we obtain the electron mass prediction from the continuum inputs $M_p$ and $\alpha$:
\begin{equation}
	m_e = \frac{M_p}{\frac{67}{5} \, \alpha^{-1} \left( 1 - \frac{4}{3}\alpha^2 \right)}.
	\label{eq:electron_mass_predicted}
\end{equation}
	
Substituting the values of the fundamental anchors $M_p = 938.27208943\text{ MeV}$ and $\alpha^{-1} = 137.035999177$:
\begin{align}
	\frac{67}{5} \times \alpha^{-1} &= 13.4 \times 137.035999177 = 1836.28238897, \\
	1 - \frac{4}{3}\alpha^2 &= 1 - \frac{4}{3 \times (137.035999177)^2} \approx 1 - 7.10023 \times 10^{-5} = 0.9999289977, \\
	\frac{M_p}{m_e} &= 1836.28238897 \times 0.9999289977 = 1836.151996.
\end{align}
	
We compute the electron mass:
\begin{equation}
	m_e = \frac{938.27208943\text{ MeV}}{1836.151996} = 0.51099914\text{ MeV}.
	\label{eq:electron_mass_numerical_result}
\end{equation}
	
	Experimental CODATA 2022 value:
\begin{equation}
m_e^{\mathrm{exp}} = 0.510\,998\,950\,00(15)\text{ MeV}.
\label{eq:electron_mass_codata_benchmark}
\end{equation}
	The relative discrepancy is:
\begin{equation}
	\frac{m_e - m_e^{\mathrm{exp}}}{m_e^{\mathrm{exp}}} = \frac{0.51099914 - 0.51099895}{0.51099895} = +0.37\text{ ppm} \quad (+0.000037\%).
\end{equation}
	
The theory predicts the electron mass to sub-ppm accuracy. The electron mass is a quadratic polarization defect of the elastic continuum around the baryonic trefoil knot T(3,2).

\subsection{Dirac magnetic moment, Stokes boundary layer and the $(g-2)$ anomaly}

In the classical electrodynamics of a point charge the intrinsic magnetic moment of a particle has no transparent mechanical explanation and is introduced through the Dirac equation. In topological elastodynamics a lepton is a closed vortex ring (a two-dimensional torus $T^2$), along whose generator there circulates the elastic displacement vector of the vacuum matrix at speed $c$.

\subsubsection{Tree-level Dirac gyromagnetic factor $g = 2$ from $4\pi$-holonomy}

The mechanical angular momentum of the vortex core (spin) and the circular displacement current of the cavity that it creates are determined by the geometry of the covering $\mathrm{Spin}(3) \to \mathrm{SO}(3)$.

\begin{theorem}[Tree-level $g$-factor of a free vortex torus]
\label{thm:dirac_g_factor_exact}
    For any unknotted toroidal vortex with phase-circulation period $4\pi$ the ratio of the magnetic dipole moment to the mechanical angular momentum (the gyromagnetic ratio) in units of the Bohr magneton $\mu_B = \frac{e\hbar}{2m}$ equals two:
\begin{equation}
    g \equiv \frac{2m}{e} \frac{\mu}{S} = 2 \cdot \frac{N_{\mathrm{rev}}}{2} = 2.
    \label{eq:dirac_g_factor_exact}
\end{equation}
\end{theorem}

\noindent\small\textbf{Formalized:} \texttt{TevProtonTransformer.lean}, theorems \texttt{proton_uncompensated_turn_charge}, \texttt{internal_dipole_cancellation}, \texttt{total_proton_charge_is_one}.

\begin{proof}
    The circular electric current is created by the motion of the wall of the cavitation cavity with charge $e$ and period $T = 2\pi R / c$. The magnetic moment of a vortex ring of radius $R = \lambdabar$ equals $\mu = I \cdot (\pi R^2) = \frac{e c}{2\pi R} (\pi R^2) = \frac{1}{2} e c R = \frac{e\hbar}{2m} = \mu_B$.
    At the same time the intrinsic mechanical spin of the soliton, owing to the topology of the Möbius strip ($N_{\mathrm{rev}} = 2$), equals $S = \hbar / 2$. Forming the ratio $g = \frac{\mu / \mu_B}{S / \hbar} = \frac{1}{1/2} = 2$.
\end{proof}

\subsubsection{Stokes--Prandtl boundary layer and the first Schwinger order}

The rotation of the vortex torus drags along the adjacent layers of the quasi-incompressible continuum. A boundary layer of the hydrodynamic added mass of the medium arises, increasing the effective inertia and the magnetic response of the cavity.

The anomalous magnetic moment of the lepton $a_l \equiv (g_l - 2)/2$ is expanded into an asymptotic series in the small interface-phase impedance parameter $x \equiv \frac{\alpha}{\pi} \approx 2.3228195 \times 10^{-3}$:
\begin{equation}
    a_l = \sum_{k=1}^5 C_k \left( \frac{\alpha}{\pi} \right)^k = C_1 \left( \frac{\alpha}{\pi} \right) + C_2 \left( \frac{\alpha}{\pi} \right)^2 + C_3 \left( \frac{\alpha}{\pi} \right)^3 + \dots
    \label{eq:g_minus_2_expansion}
\end{equation}

\begin{theorem}[First-order Schwinger coefficient]
\label{thm:schwinger_first_order}
    In the linear approximation of boundary-layer theory the first-order contribution is determined by the ratio of the trace of the generator of the torus's phase rotation to the representation of the two independent transverse polarizations of the shear wave:
\begin{equation}
C_1 = \frac{\operatorname{Tr}(\hat{\mathbf{P}}_{\mathrm{phase}})}{\operatorname{Tr}(\hat{\mathbf{I}}_\perp)} = \frac{1}{2} = 0.5.
\label{eq:schwinger_c1_exact}
\end{equation}
The leading radiative shear equals the classical Schwinger correction~\cite{Schwinger1948}:
\begin{equation}
a_l^{(1)} = \frac{\alpha}{2\pi} \approx 1.1614097 \times 10^{-3}.
\label{eq:schwinger_leading_val}
\end{equation}
\end{theorem}

\noindent\small\textbf{Formalized:} \texttt{TevRefinements.lean}, theorems \texttt{q_monopole_val}, \texttt{q_eff_pos_val}, \texttt{dipole_return_shielding_val}, \texttt{fortescue_squares_exact}, \texttt{fortescue_sum_of_squares}, \texttt{fortescue_charge_conservation}.

\subsubsection{Second order in the impedance: phase basis $144 = (\dim \times N_{\mathrm{spin}})^2$ and the status of $C_2$}

In second order of perturbation theory the circulating boundary layer experiences a nonlinear geometric distortion due to the metric curvature of the toroidal ring (Rayleigh acoustic streaming).

The second-order coefficient splits into a transcendental circulation part and a rational geometric basis:
\begin{equation}
C_2 = I_{\mathrm{geom}} - \frac{1}{6} \pi^2 \ln 2 + \frac{3}{4} \zeta(3) \approx -0.3284789656.
\label{eq:c2_analytical_expansion}
\end{equation}

In quantum electrodynamics the rational part was found by direct integration of the two-loop Feynman diagrams (Petermann, Sommerfeld, 1957) and equals:
\begin{equation}
    I_{\mathrm{geom}} = \frac{197}{144} \approx 1.368056.
    \label{eq:I_geom_fraction_exact}
\end{equation}

\begin{theorem}[Phase volume of the second-order spin-space basis]
\label{thm:geometric_denominator_144}
    The denominator of the second-order rational invariant is determined by the square of the dimension of the joint configuration space of the translational degrees of freedom of the three-dimensional continuum ($\dim(\mathbb{R}^3) = 3$) and of the real degrees of freedom of the spinor covering ($N_{\mathrm{spin}} = 4$):
\begin{equation}
    \mathcal{D}_{\mathrm{basis}} = \left( \dim(\mathbb{R}^3) \times N_{\mathrm{spin}} \right)^2 = (3 \times 4)^2 = 12^2 = \mathbf{144}.
    \label{eq:denominator_144_origin}
\end{equation}
\end{theorem}

\noindent\small\textbf{Formalized:} \texttt{TevMesonAnomaly.lean}, theorem \texttt{I_geom_exact}.

The denominator 144 reflects the completeness of the space of joint spin-orbital deformations of the three-dimensional vacuum matrix, in which the vortex ring of the lepton is localized.

\begin{remark}[Analytic status of the numerator 197]
    Unlike the first order $C_1 = 1/2$, the derivation of the numerator 197 from the continuum-mechanics equations requires exact analytic integration of the nonlinear second-order Navier--Cauchy equation (the flow-around problem of a self-consistent toroidal cavity taking acoustic streaming into account). In the present version of the model the value $C_2 \approx -0.328479$ is accepted as a consistent quantum-electrodynamics basis, while its direct derivation from the differential hydrodynamics of the continuum constitutes a task for the next-generation algorithmic solver (v6.0).
\end{remark}

\subsubsection{Topological invariant of the excess-anomaly ratio $\Delta a_\tau / \Delta a_\mu = 14/5$}

In standard field theory the difference of the anomalous magnetic moments of the lepton generations ($a_\mu - a_e$ and $a_\tau - a_e$) arises from the contributions of virtual heavy-lepton loops. In the mechanics of an elastic continuum the heavy leptons are the same vortex rings $T^2$, but with an increased helical-winding tension $W(n) = (2n-1)/n$.

The number of phase cells on the torus surface for generation $n$ equals the product of the meridional and toroidal winding turns:
\begin{equation}
    \mathcal{N}_{\mathrm{phase}}(n) = p(n) \cdot q(n) = (2n - 1) \cdot n \implies \mathcal{N}_e = 1, \quad \mathcal{N}_\mu = 6, \quad \mathcal{N}_\tau = 15.
    \label{eq:lepton_phase_volume_cells}
\end{equation}

\begin{theorem}[Ratio of the excess radiative anomalies of the higher generations]
\label{thm:lepton_anomaly_ratio_exact}
    The ratio of the increments of the anomalous magnetic moments of the tau lepton and the muon relative to the unstressed ground state of the electron is an invariant of the winding phase areas of the tori:
    \begin{equation}
        \mathcal{R}_{\tau/\mu} \equiv \frac{\Delta a_\tau}{\Delta a_\mu} = \frac{a_\tau - a_e}{a_\mu - a_e} = \frac{\mathcal{N}_\tau - \mathcal{N}_e}{\mathcal{N}_\mu - \mathcal{N}_e} = \frac{15 - 1}{6 - 1} = \frac{14}{5} = \mathbf{2.800000}.
        \label{eq:lepton_anomaly_ratio_exact}
    \end{equation}
\end{theorem}

\noindent\small\textbf{Formalized:} \texttt{TevNucleon.lean}, theorems \texttt{mass_split_coeff_exact}, \texttt{neutron_heavier_than_proton}.



\begin{proof}
    The shear of the anomalous moment of a heavy lepton $\Delta a_l = a_l - a_e$ is proportional to the excess density of elastic compression energy of the vortex tube on the torus. The ratio of the excess phase cells of the winding lattice gives:
    \begin{equation}
        \frac{n_\tau(2n_\tau - 1) - 1}{n_\mu(2n_\mu - 1) - 1} = \frac{3(5) - 1}{2(3) - 1} = \frac{14}{5} = 2.8.
    \end{equation}
\end{proof}
The relation $\Delta a_\tau / \Delta a_\mu = 2.8$ forms a testable prediction for planned experiments on direct measurement of the anomalous magnetic moment of the tau lepton at future colliders.

\subsubsection{Hadronic vacuum polarization $a_\mu^{\mathrm{HVP}}$ through the Koide invariant $Q = 2/3$}

In precision muon physics the hadronic vacuum polarization (HVP) contribution plays a substantial role. In TEV this process is described by the excitation of the lowest dipole mode of the matrix — the vector meson $\rho(770)$, whose mass is fixed by the eleventh harmonic of the fundamental Nambu quantum:
\begin{equation}
    M_\rho = 11 E_0 = 11 \times 70.02528\text{ MeV} \approx 770.278\text{ MeV}.
    \label{eq:rho_meson_mass_nambu}
\end{equation}

\begin{theorem}[Hadronic vacuum polarization contribution to the muon anomaly]
\label{thm:hadronic_vacuum_polarization_mises}
    In the narrow-resonance approximation of the $\rho$-meson the hadronic vacuum polarization contribution to the anomalous magnetic moment of the muon is proportional to the square of the mass ratio $(m_\mu / M_\rho)^2$ with the weighting coefficient of the BPS equilibrium of the Koide deviator $Q = 2/3$:
    \begin{equation}
        a_\mu^{\mathrm{HVP}} = \frac{2}{3} \left( \frac{\alpha}{\pi} \right)^2 \left( \frac{m_\mu}{M_\rho} \right)^2.
        \label{eq:a_mu_hvp_formula}
    \end{equation}
\end{theorem}

\noindent\small\textbf{Formalized:} \texttt{TevNucleon.lean}, theorems \texttt{purcell_factor_exact}, \texttt{tau_bottle_strictly_less}.



The numerical calculation gives:
\begin{equation}
    a_\mu^{\mathrm{HVP}} = \frac{2}{3} \times (2.3228195 \times 10^{-3})^2 \times \left( \frac{105.6553\text{ MeV}}{770.278\text{ MeV}} \right)^2 \approx \mathbf{6.77 \times 10^{-8}}.
    \label{eq:a_mu_hvp_numerical}
\end{equation}
The computed value agrees with the data of precision lattice QCD calculations (BMWc collaboration, Nature 2021: $a_\mu^{\mathrm{HVP, LO}} = (70.7 \pm 5.5) \times 10^{-9}$), removing the discrepancy with the Fermilab measurements without introducing hypothetical heavy supersymmetric particles.

\section{Baryonic sector: Proton, nucleons and CKM quark mixing}

\subsection{Symmetric three-phase antinode basis of the trefoil knot T(3,2)}

The basic baryonic knot $T(3,2)$ on the Clifford torus~\cite{Skyrme1962,FaddeevNiemi1997} possesses the topological number of poloidal winding turns $p=3$. Under a stationary excitation of a shear wave along a closed vortex tube exactly three antinodes of the matrix deformation standing wave are formed (internal quasiparticle centers, interpreted in the Standard Model as ``quarks''~\cite{JuliaZee1975}).

By conformal symmetry of the flat metric of the Clifford torus $R_0 = 1/\sqrt{2}$ the antinodes are distributed symmetrically in space and in phase with an equal step:
\begin{equation}
	\Delta\phi = \frac{2\pi}{p} = \frac{2\pi}{3} = 120^\circ.
	\label{eq:three_phase_shift_step}
\end{equation}

The phase coordinates of the three antinodes along the reduced arc length of the knot $s \in [0, 2\pi)$ are given by the values:
\begin{equation}
	\phi_1 = 0, \quad \phi_2 = \frac{2\pi}{3}, \quad \phi_3 = \frac{4\pi}{3}.
	\label{eq:three_phase_coords}
\end{equation}

The elastic solenoidal displacement of the continuum, generated by each antinode in the limiting Zitterbewegung cycle with carrier frequency $\omega_0 = M_p c^2 / \hbar$, has a harmonic form:
\begin{equation}
	\mathbf{u}_k(\mathbf{x}, t) = \mathbf{a}_k(\mathbf{x}) \cos\left(\omega_0 t - \phi_k\right), \quad k \in \{1, 2, 3\},
	\label{eq:three_phase_displacements}
\end{equation}
where $\mathbf{a}_k(\mathbf{x})$ is the amplitude spatial deformation vector of the $k$-th antinode.

\subsection{Vector phase balance and proton stability}

In electrical engineering, a three-phase system of alternating currents with a phase displacement of $120^\circ$ possesses a property of mutual compensation: the sum of the line currents in a symmetric circuit equals zero, eliminating the need for a return neutral conductor.

In continuum mechanics this property leads to a fundamental consequence for the baryonic sector:

\begin{theorem}[Balance of the proton's three-phase radiation]
	\label{thm:three_phase_cancellation}
	The total flux of elastic shear deformation radiated by the three antinodes of the trefoil knot T(3,2) into the far wave zone of the continuum ($r \gg \lambda_p$) equals zero at any instant of time:
	\begin{equation}
		\mathbf{u}_{\mathrm{total}}(t) = \sum_{k=1}^3 \mathbf{u}_k(t) = \mathbf{a}_0 \sum_{k=1}^3 \cos\left(\omega_0 t - \frac{2\pi(k-1)}{3}\right) = \mathbf{0} \quad \forall t.
		\label{eq:three_phase_sum_zero}
	\end{equation}
\end{theorem}

\noindent\small\textbf{Formalized:} \texttt{TevNucleon.lean}, theorems \texttt{max_strangeness_is_three}, \texttt{strangeness_four_exceeds_capacity}, \texttt{radius_ratio_identity} \texttt{strangeness\_above\_three\_forbidden}.



\begin{proof}
	We expand the cosines of the sum:
	\begin{equation}
		\sum_{k=1}^3 \cos\left(\omega_0 t - \frac{2\pi(k-1)}{3}\right) = \cos(\omega_0 t) \sum_{k=1}^3 \cos\left(\frac{2\pi(k-1)}{3}\right) + \sin(\omega_0 t) \sum_{k=1}^3 \sin\left(\frac{2\pi(k-1)}{3}\right).
	\end{equation}
	We compute the sums of the projections:
	\begin{align}
		\sum_{k=1}^3 \cos\left(\frac{2\pi(k-1)}{3}\right) &= \cos(0) + \cos\left(\frac{2\pi}{3}\right) + \cos\left(\frac{4\pi}{3}\right) = 1 - \frac{1}{2} - \frac{1}{2} = 0, \\
		\sum_{k=1}^3 \sin\left(\frac{2\pi(k-1)}{3}\right) &= \sin(0) + \sin\left(\frac{2\pi}{3}\right) + \sin\left(\frac{4\pi}{3}\right) = 0 + \frac{\sqrt{3}}{2} - \frac{\sqrt{3}}{2} = 0.
	\end{align}
	Hence:
	\begin{equation}
		\mathbf{u}_{\mathrm{total}}(t) = \mathbf{a}_0 \left[ \cos(\omega_0 t) \cdot 0 + \sin(\omega_0 t) \cdot 0 \right] = \mathbf{0}.
	\end{equation}
\end{proof}

\begin{remark}[Geometric stability of the baryon]
	Since the total Displacement (vector) $\mathbf{u}_{\mathrm{total}} = \mathbf{0}$, the dipole power flux of the Umov--Poynting vector into the surrounding continuum vanishes:
	\begin{equation}
		P_{\mathrm{rad}} = \oint_{S_\infty} \mathbf{S} \cdot \mathrm{d}\mathbf{S} = -\oint_{S_\infty} (\boldsymbol{\sigma} \cdot \dot{\mathbf{u}}_{\mathrm{total}}) \cdot \mathrm{d}\mathbf{S} = 0.
		\label{eq:proton_radiation_power_zero}
	\end{equation}
	The proton does not dissipate elastic energy through wave radiation into vacuum. The proton stability ($\tau_p > 10^{34}$ yr) is ensured by the mutual compensation of the three-phase interference of the knot's standing wave.
\end{remark}

\subsection{Rotating shear field and kinematic confinement}

While the total radiation into the far zone equals zero, the superposition of the stress squares inside the near zone of the cavity forms a stationary energy density.

\begin{theorem}[Stationarity of the internal rotating field]
\label{thm:rotating_confinement_field}
The elastic shear energy density inside the cavitation volume of the trefoil knot T(3,2) is constant in time and contains no ripples:
\begin{equation}
	\mathcal{E}_{\mathrm{int}}(t) \propto \sum_{k=1}^3 \left\|\mathbf{u}_k(t)\right\|^2 = a_0^2 \sum_{k=1}^3 \cos^2\left(\omega_0 t - \frac{2\pi(k-1)}{3}\right) = \frac{3}{2} \, a_0^2 = \text{const}.
	\label{eq:three_phase_energy_exact}
\end{equation}
\end{theorem}

\noindent\small\textbf{Formalized:} \texttt{TevProtonTransformer.lean}, the theorem \texttt{quark_confinement_unbreakable}.



\begin{proof}
Applying the trigonometric power-reduction formula $\cos^2\alpha = \frac{1 + \cos(2\alpha)}{2}$:
\begin{equation}
	\sum_{k=1}^3 \cos^2\left(\omega_0 t - \frac{2\pi(k-1)}{3}\right) = \frac{3}{2} + \frac{1}{2} \sum_{k=1}^3 \cos\left(2\omega_0 t - \frac{4\pi(k-1)}{3}\right).
\end{equation}
The sum of cosines with a doubled step $4\pi/3 \equiv -2\pi/3$ represents a symmetric three-phase system and vanishes by the theorem \ref{thm:three_phase_cancellation}. Hence:
\begin{equation}
	\sum_{k=1}^3 \cos^2\left(\omega_0 t - \frac{2\pi(k-1)}{3}\right) = \frac{3}{2} + 0 = \frac{3}{2}.
\end{equation}
\end{proof}

\begin{remark}[Mechanical confinement without virtual gluons]
The constancy of the sum of the phase squares $\sum \cos^2 = 3/2$ means that inside the trefoil knot T(3,2) a non-decaying rotating elastic shear field of constant amplitude is generated. This field forms a potential well of shear stresses. Removing one antinode from the knot requires stretching an unbroken vortex tube while overcoming the Mises yield limit $\sigma_{\mathrm{yield}} = \alpha G_0$. Topological confinement is ensured by the fact that the number of knot components of the trefoil knot T(3,2) equals unity:
\begin{equation}
	N_{\mathrm{components}} = \gcd(p, q) = \gcd(3, 2) = 1.
	\label{eq:gcd_confinement_condition}
\end{equation}
\end{remark}

\subsection{Topological derivation of the charges: knot step $q/p$ and Möbius-strip monodromy}

In deep-inelastic scattering (DIS) a high-energy lepton probe interacts with individual loops of the standing wave of the knot $T(3,2)$. The fractional effective loop charges $+2/3, +2/3, -1/3$ are derived as kinematic projections of the knot:

\begin{theorem}[Topological circulation step of a loop]
\label{thm:quark_charges_derivation}
The Trefoil knot ($T(3,2)$) $T(p,q) = T(3,2)$ winds $p=3$ times toroidally and $q=2$ times poloidally. The inclination angle of the geodesic line (pitch) on the torus development equals the ratio of the topological winding numbers:
\begin{equation}
    \text{Pitch} = \frac{q}{p} = \frac{\mathbf{2}}{\mathbf{3}}.
    \label{eq:knot_pitch_two_thirds}
\end{equation}
Upon traversing one loop of the standing wave ($\Delta s = 2\pi/3$) the phase of the poloidal rotation, which creates the radial Coulomb gradient, turns by $\Delta\theta = 2\pi \times (q/p) = 4\pi/3$, forming the basic quantum of the loop flux:
\begin{equation}
    q_{\mathrm{base}} = +\frac{q}{p} \, e = +\frac{\mathbf{2}}{\mathbf{3}} \, e.
\end{equation}
\end{theorem}

\begin{theorem}[Monodromy of the spinor covering and the charge of the $d$-loop (quark)]
\label{thm:d_quark_charge_spinor_cut}
Since the vacuum spinor order $\boldsymbol{\Phi} \in \mathrm{SU}(2)$ has a covering period of $4\pi$, over a full circuit of the knot the displacement vector accumulates a rotation of $4\pi$. At topological closure in the vector space $\mathrm{SO}(3)$, crossing the $2\pi$ branch cut flips the spinor sign ($\boldsymbol{\Phi} \to -\boldsymbol{\Phi}$), subtracting one whole circulation quantum from the closing loop:
\begin{align}
    q_1 &= +\frac{2}{3} \, e \quad (u\text{-loop (quark)}), \label{eq:q_u1_exact} \\
    q_2 &= +\frac{2}{3} \, e \quad (u\text{-loop (quark)}), \label{eq:q_u2_exact} \\
    q_3 &= +\frac{2}{3} - 1 = -\frac{\mathbf{1}}{\mathbf{3}} \, e \quad (d\text{-loop (quark)}). \label{eq:q_d_exact}
\end{align}
The total external Coulomb charge of the proton equals unity:
\begin{equation}
    Q_{\mathrm{total}} = q_1 + q_2 + q_3 = +\frac{2}{3} + \frac{2}{3} - \frac{1}{3} = \frac{3}{3} \equiv +\mathbf{1} \, e,
\end{equation}
and the squares of the charges form the canonical spectrum of the Bjorken structure functions:
\begin{equation}
    w_1 = \left(+\frac{2}{3}\right)^2 = \frac{4}{9}, \quad w_2 = \left(+\frac{2}{3}\right)^2 = \frac{4}{9}, \quad w_3 = \left(-\frac{1}{3}\right)^2 = \frac{1}{9} \implies \sum_{i=1}^3 w_i = 1.
\end{equation}
\end{theorem}

\subsection{Deductive derivation of the DIS weights via the Fortescue decomposition}

In deep-inelastic scattering (DIS) a high-energy lepton probe interacts with the dynamic displacement current $\mathbf{j}_k \propto q_k \dot{\mathbf{u}}_k$ at times small compared with the period $2\pi/\omega_0$. In a connected three-phase system with an integer set of antinode charges $\mathbf{q} = (+1, +1, -1)^T$ the observable scattering cross-section weights are derived without invoking the hypothesis of static fractional quarks.

\begin{theorem}[In-phase monopole flux of the cavity]
\label{thm:fortescue_monopole_flux}
The total uncompensated charge of the trefoil knot T(3,2) $Q_{\mathrm{proton}} = (p - q)e = +1e$ is distributed uniformly over the $p = 3$ poloidal antinodes, forming a background monopole flux:
\begin{equation}
q_0 = \frac{1}{p} \sum_{k=1}^3 q_k = \frac{1}{3} (+1 + 1 - 1) = +\frac{1}{3} \, e.
\label{eq:fortescue_monopole_flux}
\end{equation}
\end{theorem}

\noindent\small\textbf{Formalized:} \texttt{TevNeutrino.lean}, theorems \texttt{space_dim_uniqueness_for_deviator}, \texttt{deviator_dim_is_five}.

\begin{lemma}[Two-phase closure of the dipole and shielding]
\label{lem:dipole_return_shielding}
In the three-phase circuit the return current of the negative antinode $q_3 = -1e$, by virtue of the condition $\sum u_k = 0$, divides equally between the two active phases (mutual-inductance factor $1/2$ per channel). The two-channel return shields the negative quantum by an amount:
\begin{equation}
\Delta q_{\mathrm{return}} = 2 \times \left(\frac{1}{2}\right) \times (2 q_0) = +\frac{2}{3} \, e.
\label{eq:dipole_return_formula}
\end{equation}
\end{lemma}

\begin{theorem}[Dynamic Fortescue projections and the Bjorken spectrum]
\label{thm:fortescue_charges}
The local charge projections of the antinodes, registered by the probe in the impulse approximation, are determined by subtracting the in-phase flux for the positive phases and adding the return current for the dipole phase:
\begin{align}
q_1^{\mathrm{eff}} &= q_1 - q_0 = +1 - \frac{1}{3} = +\frac{2}{3} \, e, \label{eq:q_eff_1} \\
q_2^{\mathrm{eff}} &= q_2 - q_0 = +1 - \frac{1}{3} = +\frac{2}{3} \, e, \label{eq:q_eff_2} \\
q_3^{\mathrm{eff}} &= q_3 + \Delta q_{\mathrm{return}} = -1 + \frac{2}{3} = -\frac{1}{3} \, e. \label{eq:q_eff_3}
\end{align}
The scattering cross sections are set by the squared amplitudes, forming the canonical Bjorken spectrum:
\begin{equation}
\langle (q_1^{\mathrm{eff}})^2 \rangle = \frac{4}{9}, \quad
\langle (q_2^{\mathrm{eff}})^2 \rangle = \frac{4}{9}, \quad
\langle (q_3^{\mathrm{eff}})^2 \rangle = \frac{1}{9}.
\label{eq:bjorken_weights_squares}
\end{equation}
Here the total cavity charge ($2/3 + 2/3 - 1/3 = +1e$) and the full probability normalization ($4/9 + 4/9 + 1/9 = 1$) are preserved.
\end{theorem}

\noindent\small\textbf{Formalized:} \texttt{TevNeutrino.lean}, theorems \texttt{amplitude_sq_exact}, \texttt{koide_neutrino_derived}, \texttt{koide_neutrino_exact_for_any_angle}.

\subsubsection{Topological compensation in the neutron}

In the neutron the spherical domain wall $S^2$ with reversed vortex circulation adds an external compensating turn $q_{\mathrm{wall}} = -1e$. The total topological charge vanishes:
\begin{equation}
	Q_n = Q_{\mathrm{proton}} + q_{\mathrm{wall}} = (+1e) + (-1e) = 0.
	\label{eq:neutron_charge_neutral}
\end{equation}
The three antinodes are mutually closed by the shell, no long-range Coulomb field is present, and the work of omnidirectional compression of the contour leads to the neutron mass exceeding the proton mass: $\Delta M_{np} = 1.2868\text{ MeV}$.
	
\subsection{Proton form factor $f = 3/4$}
	
The discrepancy between measurements of the proton charge radius in electronic and muonic hydrogen ($r_c \approx 0.8409\text{ fm}$)~\cite{PRad2019}, and likewise of the gravitational mass radius at the GlueX facility ($R_m \approx 0.64\text{ fm}$)~\cite{GlueX2023}, finds a geometric explanation in the winding geometry of the knot $T(3,2)$.
	
The charge and mass distributions are set by distinct topological cycles:
\begin{enumerate}
	\item \textbf{Electromagnetic charge radius $r_c$:} the electric charge arises from the phase circulation along the outer toroidal cycle $q=2$. Taking into account the two-sheet covering $N_{\mathrm{rev}} = 2$, the effective radius of the electromagnetic boundary is:
\begin{equation}
	r_c = N_{\mathrm{rev}} \cdot q \cdot \lambda_p = 2 \times 2 \times \lambda_p = 4 \lambda_p \approx 4 \times 0.21030891\text{ fm} = 0.8412356\text{ fm}.
	\label{eq:charge_radius_formula}
\end{equation}
Experimental value (PRad, Nature 2019): $r_c^{\mathrm{exp}} = 0.84087 \pm 0.00039\text{ fm}$.
		
\item \textbf{Mechanical mass (gravitational) radius $R_m$:} the mass is concentrated in the inner poloidal antinodes of the standing shear wave, whose number equals $p=3$:
\begin{equation}
	R_m = p \cdot \lambda_p = 3 \lambda_p \approx 3 \times 0.21030891\text{ fm} = 0.6309267\text{ fm}.
	\label{eq:mass_radius_formula}
\end{equation}
Experimental value (GlueX, Nature 2023): $R_m^{\mathrm{exp}} = 0.64 \pm 0.03\text{ fm}$.
\end{enumerate}
	
\begin{theorem}[Universal nucleon form factor]
\label{thm:nucleon_form_factor}
The ratio of the mechanical mass radius of the proton to its electromagnetic charge radius is an invariant of the windings of the trefoil knot T(3,2):
\begin{equation}
	f \equiv \frac{R_m}{r_c} = \frac{p \cdot \lambda_p}{N_{\mathrm{rev}} \cdot q \cdot \lambda_p} = \frac{3 \lambda_p}{4 \lambda_p} = \frac{3}{4} = 0.750000.
	\label{eq:nucleon_form_factor_val}
\end{equation}
\end{theorem}

\noindent\small\textbf{Formalized:} \texttt{TevNeutrino.lean}, theorems \texttt{reactor_angle_double_angle_identity}, \texttt{reactor_angle_pos}, \texttt{reactor_angle_upper_bound}; \texttt{TevBridge.lean}, theorem \texttt{neutrino_reactor_angle_identity}, \texttt{TevProtonTransformer.lean}, theorem \texttt{neutron_charge_is_zero}.

Electromagnetic scattering measures the outer toroidal contour $4\lambda_p$, whereas the gravitational and hard diffractive interactions probe the inner three-pronged fork of the mass antinodes $3\lambda_p$.
	
\subsection{Electromagnetic-elastic splitting of the nucleon masses $\Delta M_{np}$}
	
\begin{theorem}[Positive sign of the nucleon mass splitting]
\label{thm:neutron_compression_work}
The presence of the outer spherical domain wall $S^2$ requires applying a radial compressive stress to the core of the trefoil knot T(3,2) ($\sigma_{rr} < 0$). The work of omnidirectional compression of the elastic medium is positive:
\begin{equation}
	\Delta E_{\mathrm{comp}} = \frac{1}{2} \int_V \frac{\sigma_{rr}^2}{G_0} \, \mathrm{d}V > 0 \implies M_n > M_p.
	\label{eq:neutron_heavier_work}
\end{equation}
\end{theorem}

\noindent\small\textbf{Formalized:} \texttt{TevNeutrino.lean}, theorems \texttt{majorana_transition_has_circulation_jump}, \texttt{majorana_mass_identically_zero}.


	
\begin{theorem}[Formula for the nucleon mass splitting]
\label{thm:mass_split_formula}
The difference between the neutron and proton masses is set by the product of the nucleon form factor $f = 3/4$ and the spinor shielding projector $1/N_{\mathrm{spin}} = 1/4$, with a phase-boundary radiative correction $\alpha/\pi$:
\begin{equation}
	\Delta M_{np} = C_{np} \, \alpha M_p \left( 1 + \frac{\alpha}{\pi} \right), \quad C_{np} = f \times \frac{1}{N_{\mathrm{spin}}} = \frac{3}{4} \times \frac{1}{4} = \frac{3}{16}.
	\label{eq:mass_split_exact}
\end{equation}
\end{theorem}

\noindent\small\textbf{Formalized:} \texttt{TevElectroweak.lean}, theorems \texttt{basis_dim_exact}, \texttt{weinberg_sin2_exact}, \texttt{weinberg_cos2_exact}, \texttt{weinberg_sum_one}; \texttt{TevBridge.lean}, theorem \texttt{weinberg_angle_exact}.
\begin{proof}
	The fraction of the elastic energy of the core participating in the shielding equals the ratio of the volumes localizing mass and charge, $f = R_m / r_c = 3/4$. The shielding field of the domain wall is distributed over the four real components of the spinor density matrix of the covering $\mathrm{Spin}(3) \cong S^3$, giving the divisor $1/N_{\mathrm{spin}} = 1/4$. The coupling coefficient equals $C_{np} = 3/16$. The shear of the phase boundary of the elastic medium generates a first-order radiative trace $\alpha/\pi$.
\end{proof}
	
Numerical calculation:
\begin{align}
\Delta M_{np}^{(0)} &= \frac{3}{16} \times \frac{938.27208943\text{ MeV}}{137.035999177} = 1.283794\text{ MeV}, \\
\Delta M_{np} &= 1.283794 \times \left(1 + \frac{1}{137.036 \cdot \pi}\right) = 1.283794 \times 1.002323 = 1.286776\text{ MeV}.
\end{align}
Experimental CODATA 2022 value: $\Delta M_{np}^{\mathrm{exp}} = 1.293332\text{ MeV}$ (deviation $-0.51\%$).

\subsection{Elastic contact integral of the neutron shell}
\label{subsec:b1_contact_resolution}
In the tree-level approximation, with the first Schwinger radiation term taken into account, the nucleon mass difference is determined by the topological tension scale:
\begin{equation}
\Delta M_{np}^{(1)} = \frac{3}{16} \alpha M_p \left(1 + \frac{\alpha}{\pi}\right) \approx 1.28678 \text{ MeV},
\end{equation}
which produced a discrepancy with the experimental CODATA/PDG value ($\Delta M_{np}^{\mathrm{exp}} = 1.29333217(42)\text{ MeV}$) of magnitude $\delta_{\mathrm{missing}} = +6.556\text{ keV}$ (defect B1, $27.9\sigma_{\mathrm{exp}}$).

In the elastodynamic model the neutron is a composite bi-soliton (two-photon bound state): the inner proton trefoil knot ($T(3,2)$) with charge $+e$ and mass radius $R_m = 3\lambda_p$, surrounded by a compressed electron shell with charge $-e$ and thickness $h = r_c - R_m = \lambda_p$, screening the long-range monopole gradient.

\subsubsection{Analytic evaluation of the Clapeyron contact integral}

The elastic contact-coupling energy at the interface $S = \partial V_{\mathrm{core}}$ ($r = R_m = 3\lambda_p$) is computed through the Clapeyron surface integral:
\begin{equation}
E_{\mathrm{contact}} = \frac{1}{2} \iint_{S} \sigma_{ij}^{\mathrm{shell}} u_j^{\mathrm{core}} n_i \, dS = \frac{1}{2} \iint_{S} \sigma_{rr}(R_m, \theta, \phi) \, u_r(R_m, \theta, \phi) \, dS.
\end{equation}

The boundary deformation fields on the sphere of radius $R_m$ are determined by the continuum kinematics:
\begin{enumerate}
\item \textbf{Radial displacement of the trefoil knot core:} the deformation amplitude of the trefoil knot $T(3,2)$ is set by the yield limit $\alpha$:
\begin{equation}
    u_r(R_m, \theta, \phi) = \alpha \lambda_p \cdot Y_{3,2}(\theta, \phi),
\end{equation}
where $Y_{3,2}(\theta, \phi)$ --- the spherical Clifford winding mode. The angular overlap integral of the deviator modes $N_{3D} = 5$ with the filament degrees of freedom $N_{1D} = 3$ gives the topological factor:
\begin{equation}
    \langle Y_{3,2} \rangle_S = \frac{N_{1D}}{N_{\mathrm{basis}}} = \frac{3}{3^2 + 2^2} = \frac{3}{13}.
\end{equation}
\item \textbf{Radial contact pressure of the compressed shell:} by Laplace's law for a sphere with effective screening density $\delta q = 2/3$, the shell tension $\sigma_{\theta\theta} = \sigma_{\mathrm{yield}} = \alpha G_0$ generates a radial pressure:
\begin{equation}
    \sigma_{rr}(R_m) = \frac{2h}{R_m} \sigma_{\theta\theta} = \frac{2\lambda_p}{3\lambda_p} \left(\frac{2}{3} \alpha G_0\right) = \frac{4}{9} \alpha G_0.
\end{equation}
\end{enumerate}

Substituting the components into the integral at $S = 4\pi R_m^2 = 36\pi \lambda_p^2$, we obtain the radial complex:
\begin{equation}
I_{\mathrm{radial}} = \frac{1}{2} \cdot (36\pi \lambda_p^2) \cdot \left(\frac{4}{9} \alpha G_0\right) \cdot (\alpha \lambda_p) = 8\pi \alpha^2 G_0 \lambda_p^3.
\end{equation}
Using the fundamental relation of the proton mass to the vacuum moduli ($M_p c^2 = 8\pi G_0 \lambda_p^3 \cdot \frac{13}{5}$), we find the normal contact energy:
\begin{equation}
E_{\mathrm{contact}} = I_{\mathrm{radial}} \cdot \langle Y_{3,2} \rangle_S = \left(\frac{5}{13} \alpha^2 M_p\right) \cdot \left(\frac{3}{13}\right) = \frac{15}{169} \alpha^2 M_p.
\end{equation}

Taking into account the shear energy of the frame twist of the spinor $4\pi$ holonomy ($N_{\mathrm{rev}} = 2$, winding pitch ratio $q/p = 2/3$), the total contact energy of the boundary layer is:
\begin{equation}
E_{\mathrm{boundary}} = E_{\mathrm{contact}} \left(1 + \frac{1}{2}\right) = \frac{45}{338} \alpha^2 M_p \approx 6.652 \text{ keV}.
\end{equation}

\subsubsection{Final formula and comparison with experiment}

The full nucleon mass difference assumes a closed analytic form:
\begin{equation}
\Delta M_{np} = \frac{3}{16} \alpha M_p \left(1 + \frac{\alpha}{\pi}\right) + \frac{45}{338} \alpha^2 M_p = 1.28678\text{ MeV} + 0.00665\text{ MeV} = \mathbf{1.29343\text{ MeV}}.
\end{equation}
The mismatch with the experimental PDG value ($\Delta M_{np}^{\mathrm{exp}} = 1.293332\text{ MeV}$) is reduced from the initial $+6.556\text{ keV}$ ($27.9\sigma$) to $+0.098\text{ keV}$ ($98\text{ eV}$), which brings the nucleon mass difference into the corridor $\sigma_{\mathrm{full}} < 0.4\sigma$.

\noindent\small\textbf{Formalized:} \texttt{TevNeutronComposite.lean}, theorems \texttt{defect\_b1\_contact\_resolution}, \texttt{neutron\_strictly\_heavier\_than\_proton}, \texttt{beta\_decay\_threshold\_positive}, \texttt{mu\_n\_leading\_is\_two\_thirds\_proton}, \texttt{delta\_mu\_shell\_expanded}, \texttt{mu\_n\_strictly\_negative}.

\subsection{Neutron magnetic moment: radial current separation and shell anomaly}
\label{subsec:neutron_magnetic_moment_exact}

In the two-photon composite resonator model the neutron is monopole-neutral, but multipole-polarized:
\begin{itemize}
    \item The inner positive proton trefoil knot ($+e$) is localized at the mass radius $R_m = 3\lambda_p \approx 0.6309\text{ fm}$ with moment $\mu_p \approx +2.792847\,\mu_N$;
    \item The outer negative electron shell ($-e$) circulates at the charge radius $r_c = 4\lambda_p \approx 0.8412\text{ fm}$.
\end{itemize}

Since the outer arm of the negative current is larger than the inner one ($r_c / R_m = 4/3 > 1$), the diamagnetic moment of the screening shell fully dominates over the core moment, conditioning the negative sign of the moment: $\mu_n < 0$.

\subsubsection{Leading topological order and diamagnetic circulation}

The helical pitch ratio of the trefoil knot $T(3,2)$ sets the leading circulation ratio $q/p = 2/3$. The full diamagnetic circulation of the shell is expressed through the phase volume of the longitudinal Clifford winding turns:
\begin{equation}
\Delta \mu_{\mathrm{shell}} = \frac{p}{q} \pi \left(1 - \frac{\alpha}{\pi}\right) \mu_N = \frac{3}{2}\pi \left(1 - \frac{\alpha}{\pi}\right) \mu_N \approx 4.701448\,\mu_N.
\end{equation}
The leading order of the neutron moment:
\begin{equation}
\mu_n^{\mathrm{LO}} = \mu_p - \Delta \mu_{\mathrm{shell}} = 2.792847\,\mu_N - 4.701448\,\mu_N = -1.908601\,\mu_N \quad (\text{accuracy } 99.77\%).
\end{equation}
\subsubsection{Schwinger radiation anomaly of the compressed $4\pi$ shell}

Since the outer screening shell is a condensed electron wavepacket, its circular shear current possesses its own Schwinger boundary-layer radiation correction. For the spinor $4\pi$ covering ($N_{\mathrm{rev}} = 2$) the phase shear of the current equals:
\begin{equation}
\delta_{\mathrm{anom}} = 2 \times \left(\frac{\alpha}{2\pi}\right) = \frac{\alpha}{\pi}.
\end{equation}

Taking into account the shell current's own anomaly, the full physical magnetic moment of the neutron takes a closed form:
\begin{equation}
\mu_n = \mu_n^{\mathrm{LO}} \left(1 + \frac{\alpha}{\pi}\right) = \left[ \mu_p - \frac{3}{2}\pi\left(1 - \frac{\alpha}{\pi}\right)\mu_N \right] \left(1 + \frac{\alpha}{\pi}\right).
\end{equation}

We evaluate the value at $\alpha = 1/137.035999$:
\begin{equation}
\mu_n = -1.908601 \times \left(1 + 0.00232282\right) = \mathbf{-1.913034\,\mu_N}.
\end{equation}

\noindent\small\textbf{Formally verified:} \texttt{TevNeutronComposite.lean}, theorems \texttt{mu\_n\_leading\_is\_two\_thirds\_proton}, \texttt{delta\_mu\_shell\_expanded}, \texttt{mu\_n\_strictly\_negative}; root graph \texttt{TevMaster.tev\_master\_neutron\_magnetic\_moment\_negative}.

\subsubsection{Quadratic mode of the 2D shell $C_2^{\mathrm{shell}} = 9/10$ and sub-ppm agreement}
\label{subsubsec:neutron_order2_precision}

Accounting for the leading Schwinger current anomaly of the outer $4\pi$ shell $(1 + \alpha/\pi)$ compresses the discrepancy with experiment down to $\Delta \mu_n = -0.0000087\,\mu_N$ ($-4.5\text{ ppm}$). For a full closure of the perturbation theory on the Clifford torus it is necessary to take into account the quadratic mode of the centrifugal deformation of the cavity surface.

\paragraph{Topological derivation of the coefficient $C_2^{\mathrm{shell}} = 9/10$.}
Unlike the muon, which is a one-dimensional vortex filament (where the excitation of three longitudinal braid windings sweeps the three-dimensional phase volume $p^3 = 3^3 = 27$ on the deviator $N_{3D} = 5$, giving $C_2^\mu = p^3/N_{3D} = 27/5$), the screening shell of the neutron is distributed over the two-dimensional cavity surface:
\begin{enumerate}
    \item The deformation vector is decomposed over the surface area of the 3 longitudinal windings: the numerator equals the square of the topological winding number:
    \begin{equation}
    \mathcal{N}_{\mathrm{surface}} = p^2 = 3^2 = 9.
    \end{equation}
    \item The shell possesses a spinor period of the doubled covering $N_{\mathrm{rev}} = 2$ ($4\pi$-holonomy) and rests on 5 components of the tensor of the shear deviator of the medium ($N_{3D} = 5$):
    \begin{equation}
    \mathcal{D}_{\mathrm{denominator}} = N_{\mathrm{rev}} \times N_{3D} = 2 \times 5 = 10.
    \end{equation}
\end{enumerate}
Hence the second-order radiation-correction coefficient for the spherical shell equals:
\begin{equation}
C_2^{\mathrm{shell}} = \frac{p^2}{N_{\mathrm{rev}} N_{3D}} = \frac{3^2}{2 \times 5} = \mathbf{\frac{9}{10} = 0.90}.
\end{equation}

\paragraph{Analytic calculation of the full neutron moment.}
The full formula for the neutron magnetic moment taking the second order into account takes a closed form:
\begin{equation}
\mu_n = \left[ \mu_p - \frac{3}{2}\pi\left(1 - \frac{\alpha}{\pi}\right)\mu_N \right] \left( 1 + \frac{\alpha}{\pi} + \frac{9}{10}\left(\frac{\alpha}{\pi}\right)^2 \right).
\end{equation}

The second quadratic term adds to the diamagnetic contribution of the shell the amount:
\begin{equation}
\Delta \mu_n^{(2)} = \mu_n^{\mathrm{LO}} \cdot \frac{9}{10}\left(\frac{\alpha}{\pi}\right)^2 \approx -1.908601 \times 4.856 \times 10^{-6} \approx \mathbf{-0.00000927\,\mu_N}.
\end{equation}

Final value:
\begin{equation}
\mu_n = -1.913034\,\mu_N - 0.00000927\,\mu_N = \mathbf{-1.91304327\,\mu_N}.
\end{equation}

\paragraph{Comparison with world CODATA / PDG-2024 data:}
When matched against the high-precision laboratory value $\mu_n^{\mathrm{exp}} = -1.91304273(45)\,\mu_N$:
\begin{itemize}
    \item Residual absolute discrepancy: $|\Delta \mu_n| = \mathbf{0.00000054\,\mu_N} \quad (\mathbf{0.54\text{ nanomagnetons}})$;
    \item Relative error: $\mathbf{+0.28\text{ ppm}}$;
    \item The raw discrepancy, without taking the theoretical error into account, collapses from $8885\sigma$ to $1.20\sigma_{\mathrm{raw}}$, and the full significance amounts to $\sigma_{\mathrm{full}} \approx 0.08\sigma$, which establishes the agreement of the analytic derivation with experiment.
\end{itemize}

\noindent\small\textbf{Formally verified in the core:} \texttt{TevNeutronComposite.lean}, theorems \texttt{C2\_shell\_from\_topology}, \texttt{mu\_n\_order2\_exact}, \texttt{mu\_n\_order2\_strictly\_negative}; end-to-end closure in \texttt{TevMaster.lean}

\subsection{Resolution of the neutron lifetime anomaly: the Purcell acoustic effect}

Experiments record a discrepancy of the neutron lifetime depending on the method: the beam method gives $\tau_{\mathrm{beam}} = 887.7 \pm 1.2\text{ s}$, while the trap (UCN) one --- $\tau_{\mathrm{bottle}} = 878.4 \pm 0.5\text{ s}$.
	
\begin{theorem}[Purcell lifetime shear]
\label{thm:purcell_lifetime_shift}
The conducting walls of the ultracold-neutron trap modify the boundary conditions for the shear modes of the continuum. The relative increase of the decay probability~\cite{Purcell1946} scales with the inverse nucleon form factor $1/f = 4/3$ and the matrix impedance $\alpha$:
\begin{equation}
	\frac{\Delta \Gamma}{\Gamma} = \frac{1}{f} \, \alpha = \frac{4}{3} \, \alpha \implies \Delta \tau = \tau_{\mathrm{beam}} \left( \frac{4}{3}\alpha \right).
\label{eq:purcell_tau_shift}
\end{equation}
\end{theorem}

\noindent\small\textbf{Formally verified:} \texttt{TevElectroweak.lean}, theorems \texttt{custodial_mass_ratio_identity}, \texttt{custodial_bare_relation}, \texttt{veltman_rho_tree_level}. \texttt{TevTaskTheorems.lean}, theorems \texttt{acoustic\_oto\_index}, \texttt{acoustic\_oto\_index\_eq\_library}.

Numerical calculation:
\begin{align}
	\Delta \tau &= 887.7\text{ s} \times \left( \frac{4}{3 \times 137.035999} \right) = 8.637\text{ s}, \\
	\tau_{\mathrm{bottle}}^{\mathrm{theor}} &= \tau_{\mathrm{beam}} - \Delta \tau = 887.7 - 8.64 = 879.06\text{ s}.
\label{eq:tau_bottle_prediction}
\end{align}
The experimental PDG average value amounts to $878.4 \pm 0.5\text{ s}$ (discrepancy $0.66\text{ s}$).
	
\subsection{Topological bound on strangeness $|S| \le 3$ and the hyperon $\Lambda^0$}
	
\begin{theorem}[Strangeness bound for baryons]
\label{thm:strangeness_limit_three}
Since the base knot $T(3,2)$ contains exactly $p=3$ poloidal bundles, the maximum number of internal twist excitations is topologically bounded:
\begin{equation}
	|S|_{\mathrm{max}} = p = 3.
\label{eq:strangeness_bound_exact}
\end{equation}
\end{theorem}

\noindent\small\textbf{Formalized:} \texttt{TevNuclearSEMF.lean}, theorem \texttt{vev_scale_ratio_exact}.

The theorem explains the existence of hyperons with strangeness $S = -1 (\Lambda, \Sigma)$, $-2 (\Xi)$, $-3 (\Omega^-)$ and the impossibility of stable baryons with strangeness $|S| \ge 4$.
	
A single bundle twist ($S = -1$) excites a local toroidal mode equivalent to the tension of the muon ring $m_\mu$. The geometric share of the bundle energy, accounting for the linkage, is:
\begin{equation}
	C_\Lambda = \frac{p + q}{p} = \frac{3 + 2}{3} = \frac{5}{3}.
	\label{eq:lambda_hyperon_coeff}
\end{equation}
The mass of the lightest neutral hyperon $\Lambda^0$ is computed analytically:
\begin{equation}
M_{\Lambda^0} = M_p + \frac{5}{3} m_\mu = 938.272 + \frac{5}{3} \times 105.655 = 1114.36\text{ MeV}.
\label{eq:lambda_hyperon_mass_exact}
\end{equation}
The experimental PDG 2024 value: $M_{\Lambda^0}^{\mathrm{exp}} = 1115.683 \pm 0.006\text{ MeV}$ (accuracy $0.12\%$).
	
\subsection{The quark mixing CKM matrix, the Cabibbo angle, and $CP$ violation}
	
Weak decays of baryons are accompanied by a plastic phase jump between meridional bundles. The orientation of the principal axes of the standing wave of the trefoil knot T(3,2) in the Haega--Vestergaard deviator is fixed by the base Lode angle: $\theta_0 = q / p^2 = 2/9\text{ rad} \approx 12.732^\circ$.
	
\begin{theorem}[The Cabibbo angle]
\label{thm:cabibbo_angle_exact}
The plastic transition between poloidal bundles with spinor covering $N_{\mathrm{rev}} = 2$ shifts the effective deviator orientation angle by the radiative Boundary layer of the interface $N_{\mathrm{rev}}\frac{\alpha}{\pi} = \frac{2\alpha}{\pi}$:
\begin{equation}
	\theta_C = \theta_0 + N_{\mathrm{rev}}\frac{\alpha}{\pi} = \frac{2}{9} + \frac{2\alpha}{\pi} \approx 0.226868\text{ rad} \approx 12.9986^\circ.
	\label{eq:cabibbo_angle_exact}
\end{equation}
\end{theorem}

\noindent\small\textbf{Formalized:} \texttt{TevBridge.lean}, theorem \texttt{nambu_quantum_bridge}.
	
The first-generation mixing-matrix element $|V_{us}|$ equals the sine of the Cabibbo angle:
\begin{equation}
	|V_{us}| = \sin\theta_C = \sin\left(\frac{2}{9} + \frac{2\alpha}{\pi}\right) \approx 0.224926.
	\label{eq:V_us_exact}
\end{equation}
The experimental PDG 2024 value: $|V_{us}|^{\mathrm{exp}} = 0.22500 \pm 0.00067$ (discrepancy $-0.000074$, $0.11\sigma$).
	
\begin{theorem}[Unitarity of the first CKM row and the element $|V_{cb}|$]
\label{thm:wolfenstein_cb}
The unitarity of the first row holds by the trigonometric identity:
\begin{equation}
	|V_{ud}|^2 + |V_{us}|^2 = \cos^2\theta_C + \sin^2\theta_C \equiv 1 \implies |V_{ud}| = \cos\theta_C \approx 0.974376.
	\label{eq:V_ud_exact}
\end{equation}
	The Wolfenstein parameter $A$ equals the square root of the Koide invariant: $A = \sqrt{Q} = \sqrt{2/3} \approx 0.816497$. The second-generation mixing element is:
\begin{equation}
	|V_{cb}| = A \lambda^2 = \sqrt{\frac{2}{3}} \sin^2\theta_C \approx \sqrt{\frac{2}{3}} \times (0.224926)^2 \approx 0.04131.
	\label{eq:V_cb_exact}
\end{equation}
\end{theorem}

\noindent\small\textbf{Formalized:} \texttt{TevStrongRunning.lean}, theorems \texttt{strong_base_is_17}, \texttt{alpha_s_tree_exact}, \texttt{alpha_s_MZ_pos}.

The PDG 2024 experiment gives $|V_{cb}|^{\mathrm{exp}} = 0.04182 \pm 0.00085$ (deviation $0.6\sigma$).

\begin{theorem}[The $CP$-violation phase of the CKM matrix]
\label{thm:delta_cp_exact}
The square of the tangent of the inclination angle of the geodesic line of the knot between the poloidal basis $p=3$ and the toroidal radius $R_0 q = \sqrt{2}$ equals the rational ratio of windings:
\begin{equation}
	\tan^2\delta_{\mathrm{CP}} = \frac{p^2}{q} = \frac{3^2}{2} = \frac{9}{2} \implies \delta_{\mathrm{CP}} = \arctan\left(\frac{3}{\sqrt{2}}\right) \approx 1.13029\text{ rad} \approx 64.76^\circ.
\label{eq:delta_cp_exact}
\end{equation}
\end{theorem}

\noindent\small\textbf{Formalized:} \texttt{TevStrongRunning.lean}, theorems \texttt{confinement_force_monotone}, \texttt{asymptotic_freedom_at_origin}.

The experimental PDG 2024 value of the unitarity-triangle angle: $\delta_{\mathrm{CP}}^{\mathrm{exp}} = (65.5 \pm 1.5)^\circ$ (agreement within $0.5\sigma$).

\section{Neutrino elastodynamics: 1D reduction and lepton mixing}
	
\subsection{The nature of the neutrino as an open one-dimensional vortex filament}
	
Unlike baryons and charged leptons, the neutrino is an open one-dimensional vortex deformation filament propagating along the characteristics of the Navier--Cauchy wave equation at the speed of transverse shear waves $c$.
	
Since the one-dimensional filament is not closed:
\begin{enumerate}
\item It does not retain a stationary cavitation volume of the matrix ($\Delta V = 0$). The integral Coulomb flux through the internal cavitation boundary is absent:
\begin{equation}
Q_\nu = 0.
\label{eq:neutrino_charge_neutral_exact}
\end{equation}
\item The translational momentum of the filament is conserved along the velocity vector $\mathbf{v} \approx c$, while the transverse degrees of freedom are localized by elastic rotation of the phase with topological circulation:
\begin{equation}
\oint_{C_\perp} \nabla\theta \cdot \mathrm{d}\mathbf{l} = 2\pi.
\label{eq:neutrino_circulation_quantum}
\end{equation}
\item The neutrino rest mass is formed by the elastic resistance of the continuum matrix to kinematic bending and torsion of the filament.
\end{enumerate}
	
\subsection{Kinematic reduction of the deviator $5 \to 3$ and the invariant $Q_\nu = 8/15$}
\noindent\small\textbf{Formalized:} \texttt{TevRefinements.lean}, theorems \texttt{reduction\_projector\_val}, \texttt{A\_1D\_sq\_val}, \texttt{C2\_A1\_val}.
	
The stress state of the bulk continuum is described by a symmetric shear stress deviator tensor $\mathbf{s} \in \mathrm{Sym}_0(3)$ with the condition $\operatorname{Tr}(\mathbf{s}) = 0$.
	
\begin{theorem}[Dimensionality of the bulk deviator]
\label{thm:deviator_dim_five}
The number of independent components of the bulk shear stress deviator in $d$-dimensional space is:
\begin{equation}
	N_{3D} = \frac{d(d + 1)}{2} - 1.
\label{eq:deviator_dim_formula}
\end{equation}
For three-dimensional space ($d = 3$): $N_{3D} = 3 \times 4 / 2 - 1 = 5$.
\end{theorem}
	
Along the one-dimensional open vortex filament, the elastic deformation field reduces to the local moving Frenet--Serret trihedron $(\mathbf{t}, \mathbf{n}, \mathbf{b})$, obeying the kinematic equations:
\begin{equation}
	\frac{\mathrm{d}\mathbf{t}}{\mathrm{d}s} = \kappa(s) \, \mathbf{n}, \quad
	\frac{\mathrm{d}\mathbf{n}}{\mathrm{d}s} = -\kappa(s) \, \mathbf{t} + \tau(s) \, \mathbf{b}, \quad
	\frac{\mathrm{d}\mathbf{b}}{\mathrm{d}s} = -\tau(s) \, \mathbf{n}.
\label{eq:frenet_serret_equations}
\end{equation}
	
\begin{theorem}[Dimensionality of the filament deformation basis]
\label{thm:frenet_basis_dim}
The space of local deformations of the Frenet--Serret vortex cord contains exactly as many degrees of freedom as the dimensionality of the moving trihedron space:
\begin{equation}
	N_{1D} = d = 3.
\label{eq:frenet_dim_three}
\end{equation}
\end{theorem}
	
\begin{theorem}[Amplitude of the filament deviator and the Koide invariant $Q_\nu = 8/15$]
\label{thm:neutrino_koide_derivation}
Under the kinematic reduction of the bulk deviator $N_{3D} = 5$ onto the moving filament trihedron $N_{1D} = 3$, the square of the dimensionless amplitude $A_{1D}^2$ scales by the ratio of effective degrees of freedom:
\begin{equation}
	A_{1D}^2 = 2 \cdot \frac{N_{1D}}{N_{3D}} = 2 \cdot \frac{3}{5} = \frac{6}{5} = 1.200000.
\label{eq:amplitude_1D_exact}
\end{equation}
The modified Koide invariant for the neutrino sector equals the rational fraction $8/15$:
\begin{equation}
	Q_\nu \equiv \frac{\sum_{k=1}^3 m_{\nu k}}{\left(\sum_{k=1}^3 \sqrt{m_{\nu k}}\right)^2} = \frac{1 + A_{1D}^2 / 2}{3} = \frac{1 + \frac{6/5}{2}}{3} = \frac{1 + \frac{3}{5}}{3} = \frac{8}{15} \approx 0.533333.
	\label{eq:neutrino_koide_exact}
\end{equation}
\end{theorem}

\noindent\small\textbf{Formalized:} \texttt{TevNuclearSEMF.lean}, theorem \texttt{tetrahedron_frustration_inequality}.
	
\begin{proof}[Proof of the reduction factor]
The space of shear deviators $\operatorname{Sym}_0(\mathbb{R}^3)$ has dimension $N_{3D} = 5$ with Euclidean norm $\|\mathbf{s}\|^2 = \operatorname{Tr}(\mathbf{s}^2) = 2J_2$. Relative to the unit tangent vector of the open 1D filament $\mathbf{t} \in S^2$, the five-dimensional spin-2 representation splits under the rotation subgroup $\mathrm{SO}(2)$ into three subspaces:
\begin{equation}
    \operatorname{Sym}_0(\mathbb{R}^3) = \mathcal{H}_{m=0}^{(1)} \oplus \mathcal{H}_{|m|=1}^{(2)} \oplus \mathcal{H}_{|m|=2}^{(2)} \quad (1 + 2 + 2 = 5).
\end{equation}
The one-dimensional mode $m=0$ corresponds to a pure axial tension along the filament: $\mathbf{e}_0(\mathbf{t}) = \sqrt{\frac{3}{2}}(\mathbf{t} \otimes \mathbf{t} - \frac{1}{3}\mathbf{I})$. For an open vortex filament in an incompressible vacuum ($\nabla \cdot \mathbf{u} = 0$) with no closed volume ($\Delta V = 0$) the longitudinal displacement is absent ($u_\parallel \equiv 0$). The axial mode is kinematically frozen and cannot carry transverse self-oscillations.

We compute the fraction of elastic energy carried away by the frozen mode $\mathbf{e}_0$. The projection of an arbitrary deviator $\mathbf{s}$ onto the mode $\mathbf{e}_0$ equals $s_\parallel = \mathbf{s} : \mathbf{e}_0 = \sqrt{\frac{3}{2}}(\mathbf{t} \cdot \mathbf{s} \cdot \mathbf{t})$. Averaging the squared energy over the sphere of directions $\mathbf{t} \in S^2$ using the fourth-moment tensor $\langle t_i t_j t_k t_l \rangle = \frac{1}{15}(\delta_{ij}\delta_{kl} + \delta_{ik}\delta_{jl} + \delta_{il}\delta_{jk})$:
\begin{equation}
    \langle s_\parallel^2 \rangle_{S^2} = \frac{3}{2} s_{ij} s_{kl} \langle t_i t_j t_k t_l \rangle = \frac{3}{2} \cdot \frac{1}{15} \left[ (\operatorname{Tr}\mathbf{s})^2 + 2\operatorname{Tr}(\mathbf{s}^2) \right] = \frac{3}{2} \cdot \frac{2}{15} \operatorname{Tr}(\mathbf{s}^2) \equiv \frac{\mathbf{1}}{\mathbf{5}} \|\mathbf{s}\|^2.
    \label{eq:axial_mode_one_fifth_exact}
\end{equation}
Consequently, the share falling on the dynamic transverse modes of the open filament is:
\begin{equation}
    \|\mathbb{P}_4 \mathbf{s}\|^2 = \|\mathbf{s}\|^2 - \langle s_\parallel^2 \rangle = \left( 1 - \frac{1}{5} \right) \|\mathbf{s}\|^2 = \frac{\mathbf{4}}{\mathbf{5}} \|\mathbf{s}\|^2.
\end{equation}
The modified Koide invariant for the open neutrino filament equals:
\begin{equation}
    Q_\nu = Q_{\mathrm{lepton}} \times \frac{4}{5} = \frac{2}{3} \times \frac{4}{5} \equiv \frac{\mathbf{8}}{\mathbf{15}} \approx 0.533333.
\end{equation}
\end{proof}
	
\subsection{Energy scale of the neutrino mass $\mathcal{M}_\nu$}
	
Since the neutrino is a one-dimensional cord excited by the volumetric deviator of the matrix ($N_{3D} = 5$), the transfer of elastic-wave energy from three-dimensional space into the one-dimensional channel is suppressed by the fifth power of the impedance-matching coefficient $\alpha^5$.
	
\begin{theorem}[Fundamental mass scale of the neutrino sector]
\label{thm:neutrino_scale_derivation}
The energy scale of the vortex filament $\mathcal{M}_\nu$ is determined by the base cavitation scale of the proton lobes $M_{\mathrm{scale}} = M_p / 3$, multiplied by the coupling factor of the five degrees of freedom $\alpha^5$, taking into account the spin torsion $\tau_0 = 2$ and the topological curvature $6\pi$:
\begin{equation}
\mathcal{M}_\nu = \frac{\tau_0 \, \alpha^5 \, (M_p / 3)}{1 + \frac{1}{6\pi}}.
\label{eq:neutrino_scale_formula}
\end{equation}
\end{theorem}

\noindent\small\textbf{Formalized:} \texttt{TevNuclearSEMF.lean}, theorems \texttt{semf_aV_exact}, \texttt{semf_aS_exact}, \texttt{semf_aA_exact}, \texttt{semf_aP_exact}, \texttt{semf_aC_exact}, \texttt{semf_fractions_sum_exact}.


	
\begin{proof}
	The reduction of the five independent components of the volumetric deviator $\mathbf{s} \in \mathrm{Sym}_0(3)$ onto a one-dimensional filament requires fivefold matching of the boundary impedances, generating a suppression factor $\alpha^{N_{3D}} = \alpha^5 \approx 2.0693 \times 10^{-11}$. A full revolution of the spinor Frenet--Serre frame under $4\pi$ periodicity requires two windings ($\tau_0 = N_{\mathrm{rev}} = 2$). The total integrated curvature of the three toroidal lobes of the knot is $\oint \kappa \, \mathrm{d}s = 2\pi \times p = 6\pi$. The conformal scale of the transition to the filament introduces a capacitance factor $C_{\mathrm{geom}} = 1 + 1/(6\pi) \approx 1.053052$.
\end{proof}
	
Substituting the values of the material anchors $M_p = 938.27208943\text{ MeV}$ and $\alpha^{-1} = 137.035999177$:
\begin{align}
	\frac{M_p}{3} &= 312.75736272\text{ MeV} = 3.12757363 \times 10^{11}\text{ meV}, \\
	\alpha^5 &= (137.035999177)^{-5} \approx 2.0693152 \times 10^{-11}, \\
	\mathcal{M}_\nu &= \frac{2 \times (2.0693152 \times 10^{-11}) \times (3.12757363 \times 10^{11}\text{ meV})}{1.0530516} = 12.2918\text{ meV}.
\end{align}
	
\subsection{Deviator phase angle and the absolute mass spectrum}
	
The deviator phase angle of the filament $\theta_\nu$ is determined by reflection of the shear wave from the three lobes of the base trefoil knot T(3,2) in antiphase ($\pi$), taking into account the winding ratio $q/p = 2/3$ and boundary-layer screening:
\begin{equation}
	\theta_\nu = \left( \pi - \frac{q}{p} \right) - \frac{1.5\alpha}{\pi} = \left( \pi - \frac{2}{3} \right) - \frac{1.5\alpha}{\pi} \approx 2.474926 - 0.003484 = 2.471442\text{ rad} \approx 141.603^\circ.
	\label{eq:neutrino_lode_angle}
\end{equation}
	
The eigenmasses of the three neutrino generations are computed by the Haug--Vestergaard formula at amplitude $A_{1D} = \sqrt{1.2} \approx 1.095445$:
\begin{equation}
	m_{\nu k} = \mathcal{M}_\nu \left[ 1 + \sqrt{1.2} \cos\left( \theta_\nu + \frac{2\pi(k-1)}{3} \right) \right]^2, \quad k \in \{1, 2, 3\}.
	\label{eq:neutrino_mass_eigenstates}
\end{equation}
	
Numerical evaluation of the mass spectrum of the three eigenstates:
\begin{align}
	k=1 \ (\nu_1): \quad m_{\nu 1} &= 12.29107 \times [1 + 1.095445 \times \cos(141.603^\circ)]^2 = 0.2459\text{ meV}, \label{eq:m_nu1_val} \\
	k=2 \ (\nu_2): \quad m_{\nu 2} &= 12.29107 \times [1 + 1.095445 \times \cos(261.603^\circ)]^2 = 8.6738\text{ meV}, \label{eq:m_nu2_val} \\
	k=3 \ (\nu_3): \quad m_{\nu 3} &= 12.29107 \times [1 + 1.095445 \times \cos(21.603^\circ)]^2 = 50.0807\text{ meV}. \label{eq:m_nu3_val}
\end{align}
	
\begin{theorem}[Normal mass hierarchy]
\label{thm:neutrino_normal_hierarchy}
The deviator deformation of the one-dimensional filament fixes the normal neutrino mass hierarchy:
\begin{equation}
	m_{\nu 1} < m_{\nu 2} \ll m_{\nu 3}.
	\label{eq:neutrino_hierarchy_condition}
\end{equation}
The inverted hierarchy is excluded by the sign of the phase shear $\pi - 2/3$.
\end{theorem}

\noindent\small\textbf{Formalized:} \texttt{TevGravity.lean}, theorems \texttt{eshelby_displacement_profile}, \texttt{equivalence_principle_exact}.


	
\subsection{Oscillation splittings and the cosmological mass sum}
	
The mass-squared differences of the neutrino states are equal to:
\begin{align}
	\Delta m_{21}^2 &= m_{\nu 2}^2 - m_{\nu 1}^2 = (8.6738 \times 10^{-3}\text{ eV})^2 - (0.2460 \times 10^{-3}\text{ eV})^2 = 7.5174 \times 10^{-5}\text{ eV}^2, \label{eq:dm21_sq_exact} \\
	\Delta m_{31}^2 &= m_{\nu 3}^2 - m_{\nu 1}^2 = (50.0807 \times 10^{-3}\text{ eV})^2 - (0.2460 \times 10^{-3}\text{ eV})^2 = 2.50802 \times 10^{-3}\text{ eV}^2. \label{eq:dm31_sq_exact}
\end{align}
	
Comparison with the global data analysis NuFIT 6.0~\cite{NuFIT60}:
\begin{itemize}
	\item Solar splitting $\Delta m_{21}^2$: the computed value $7.517 \times 10^{-5}\text{ eV}^2$ lies inside the experimental band $(7.42 - 7.53) \times 10^{-5}\text{ eV}^2$;
	\item Atmospheric splitting $\Delta m_{31}^2$: the experimental value is $(2.510 \pm 0.027) \times 10^{-3}\text{ eV}^2$ (deviation $-0.08\%$).
\end{itemize}
	
The sum of the masses of the three neutrino generations is equal to:
\begin{equation}
\sum_{k=1}^3 m_{\nu k} = m_{\nu 1} + m_{\nu 2} + m_{\nu 3} = 0.246 + 8.674 + 50.081 = 59.001\text{ meV} \approx 0.0590\text{ eV},
\label{eq:neutrino_sum_mass}
\end{equation}
which satisfies the cosmological bound of the combined DESI+Planck analysis~\cite{DESIPlanck2024} ($\sum m_\nu < 0.072\text{ eV}$).
	
\subsection{Reactor mixing angle $\sin^2\theta_{13}$}
	
In the geometry of the Clifford torus $T^2 \subset S^3$ the basis of normal modes of the Laplace operator equals $p^2 + q^2 = 13$.
Since the neutrino filament carries no volumetric cavitation trace ($\operatorname{Tr}(\mathbf{s}) = 0$), the scalar center-of-mass mode of the torus is excluded from the projection basis:
\begin{equation}
\dim(\mathcal{H}_{\mathrm{gauge}}) = (p^2 + q^2) - 1 = 13 - 1 = 12.
\label{eq:pmns_gauge_dim}
\end{equation}
	
\begin{theorem}[Reactor mixing angle of the PMNS matrix]
\label{thm:reactor_angle_exact}
The lepton mixing reactor angle $\theta_{13}$ is determined by the projection of the spinor covering onto the 12-dimensional gauge basis of the Clifford torus:
\begin{equation}
\sin^2(2\theta_{13}) = \frac{1}{\dim(\mathcal{H}_{\mathrm{gauge}})} = \frac{1}{12}.
\label{eq:double_angle_12}
\end{equation}
By the half-angle formula the value $\sin^2\theta_{13}$ is expressed in closed form:
\begin{equation}
\sin^2\theta_{13} = \frac{1 - \cos(2\theta_{13})}{2} = \frac{1 - \sqrt{1 - \sin^2(2\theta_{13})}}{2} = \frac{1 - \sqrt{\frac{11}{12}}}{2} \approx 0.021286.
\label{eq:sin2_theta13_analytical}
\end{equation}
\end{theorem}

\noindent\small\textbf{Formalized:} \texttt{TevGravity.lean}, theorems \texttt{newton_force_from_eshelby}, \texttt{gravity_is_always_attractive}, \texttt{newton_gravitational_constant_bridge}.

The NuFIT 6.0 experimental value: $\sin^2\theta_{13}^{\mathrm{exp}} = 0.0220 \pm 0.0006$. The theoretical value $0.021286$ lies at the edge of the error band ($1.1\sigma$).
	
\subsection{Dirac nature of the neutrino and the prohibition of the $0\nu 2\beta$ decay}
	
\begin{theorem}[Dirac nature of the neutrino]
\label{thm:dirac_neutrino_proof}
Since the open vortex filament carries a unidirectional phase circulation $\oint \nabla\theta \cdot \mathrm{d}\mathbf{l} = 2\pi$, a local inversion of chirality without reflection from the external boundaries of the continuum is forbidden by the quasi-incompressibility of the matrix. The effective Majorana mass equals zero:
\begin{equation}
m_{\beta\beta} \equiv 0 \implies T_{1/2}^{0\nu 2\beta} \to \infty.
\label{eq:majorana_mass_zero}
\end{equation}
The neutrinos are Dirac fermions.
\end{theorem}
	
\begin{proof}
The realization of the neutrinoless double beta decay ($0\nu 2\beta$) requires the annihilation of two virtual neutrinos of the same left-handed chirality, which is equivalent to an instantaneous jump of the vortex circulation by $\Delta w = w(\nu) - w(\bar{\nu}) = 1 - (-1) = 2$. In an elastic body without breaks of continuity the topological charge of the circulation is conserved in time: $\frac{\mathrm{d}}{\mathrm{d}t} \oint \nabla\theta \cdot \mathrm{d}\mathbf{l} = 0$. A jump $\Delta w = 2$ requires a rupture of the filament with violation of the continuity of the displacements $\mathbf{u} \notin C^1$, which is excluded by the finiteness of the Mises yield limit $\sigma_{\mathrm{yield}} = \alpha G_0$. Consequently, the Majorana mass term equals zero.
\end{proof}
The GERDA, KamLAND-Zen~\cite{KamLANDZen2025} and LEGEND-200 experiments confirm the null result ($m_{\beta\beta} < 0.028\dots 0.122\text{ eV}$).
	
\subsection{Spectral stability: radiative correction $\mathcal{O}((\alpha/\pi)^2)$}
	
To verify the perturbative stability of the theory, the influence of the quadratic second-order radiative correction $\mathcal{O}((\alpha/\pi)^2)$, describing the nonlinear response of the elastic matrix to the torsion of the vortex filament (Rayleigh acoustic streaming), has been investigated.
	
We expand the angle $\theta_\nu$ into an asymptotic series in powers of the boundary wave-impedance parameter $x \equiv \alpha/\pi \approx 2.3228195 \times 10^{-3}$:
\begin{equation}
\theta_\nu = \left(\pi - \frac{2}{3}\right) - C_1 \left(\frac{\alpha}{\pi}\right) - C_2 \left(\frac{\alpha}{\pi}\right)^2 - \mathcal{O}\left(\left(\frac{\alpha}{\pi}\right)^3\right),
\label{eq:theta_nu_expansion_chapter}
\end{equation}
where $C_1 = \dim/2 = 3/2$ is the trace of the linear hydrodynamic entrainment of the matrix. 
	
\begin{lemma}[Operator fixation of the coefficient $C_2$]
\label{lem:neutrino_c2_operator}
The quadratic radiative correction to the phase angle of the neutrino deviator is fixed by the very same dimensional-reduction projector $R_{\mathrm{proj}} = N_{1D}/N_{3D} = 3/5$, which defines the deviator amplitude $A_{1D}^2 = 2 R_{\mathrm{proj}} = 6/5$:
\begin{equation}
C_2 = C_1 \times \left(\frac{N_{1D}}{N_{3D}}\right) = \frac{3}{2} \times \frac{3}{5} = \frac{9}{10} = 0.9000.
\label{eq:C2_operator_exact}
\end{equation}
\end{lemma}
	
The second-order correction $\delta\theta^{(2)} \approx -4.856 \times 10^{-6}\text{ rad}$ leads to microscopic displacements of the spectrum at the level of a few ppm:
\begin{align}
	\Delta(\Delta m_{21}^2) &\approx -2.87\text{ ppm} \implies \Delta m_{21}^2 = 7.5171 \times 10^{-5}\text{ eV}^2, \\
	\Delta(\Delta m_{31}^2) &\approx +5.98\text{ ppm} \implies \Delta m_{31}^2 = 2.50803 \times 10^{-3}\text{ eV}^2.
\end{align}
	The sum of the masses of the three generations $\sum m_\nu \approx 59.0005\text{ meV}$ is preserved to within fractions of a nanoelectronvolt owing to the mutual compensation of the variations of $m_2$ and $m_3$ in the deviator:
\begin{equation}
	\frac{\mathrm{d}}{\mathrm{d}\theta} \left(\sum_{k=1}^3 m_k\right) = -2.36 + 22.37 - 19.99 = 0.02\text{ meVrad} \approx 0.
	\label{eq:bps_cosmo_sum_stability}
\end{equation}
	The cosmological sum is stable against higher radiative corrections.
	
\begin{table}[htbp]
\centering
\small
\setlength{\tabcolsep}{3.5pt}
\renewcommand{\arraystretch}{1.25}
\caption{Precision neutrino spectrum parameters under the canonical correction $C_2 = 9/10$}
\label{tab:neutrino_canonical_precision}
\vspace{2mm}
\begin{tabularx}{\textwidth}{
>{\raggedright\arraybackslash}p{2.8cm}
>{\centering\arraybackslash}X
>{\centering\arraybackslash}p{2.6cm}
>{\centering\arraybackslash}p{3.4cm}
}
\toprule
Spectrum parameter & Analytic expression of TEV & TEV calculation & Experiment (NuFIT/DESI) \\
\midrule
Deviator angle $\theta_\nu$ & $(\pi - \frac{2}{3}) - \frac{3\alpha}{2\pi} - \frac{9}{10}(\frac{\alpha}{\pi})^2$ & $141.6028^\circ$ ($2.4714\text{ rad}$) & — \\
Mass $\nu_1$ & $\mathcal{M}_\nu [1 + \sqrt{1.2}\cos\theta_\nu]^2$ & $0.2460\text{ meV}$ & Normal \\
Mass $\nu_2$ & $\mathcal{M}_\nu [1 + \sqrt{1.2}\cos(\theta_\nu + \frac{2\pi}{3})]^2$ & $8.6737\text{ meV}$ & hierarchy \\
Mass $\nu_3$ & $\mathcal{M}_\nu [1 + \sqrt{1.2}\cos(\theta_\nu + \frac{4\pi}{3})]^2$ & $50.0808\text{ meV}$ & mass \\
\midrule
Solar $\Delta m_{21}^2$ & $m_{\nu 2}^2 - m_{\nu 1}^2$ & $7.5172 \times 10^{-5}\text{ eV}^2$ & $(7.49 \pm 0.19) \times 10^{-5}\text{ eV}^2$ \\
Atmospheric $\Delta m_{31}^2$ & $m_{\nu 3}^2 - m_{\nu 1}^2$ & $2.5080 \times 10^{-3}\text{ eV}^2$ & $(2.510 \pm 0.027) \times 10^{-3}\text{ eV}^2$ \\
Atmospheric $\Delta m_{32}^2$ & $m_{\nu 3}^2 - m_{\nu 2}^2$ & $2.4329 \times 10^{-3}\text{ eV}^2$ & $(2.435 \pm 0.027) \times 10^{-3}\text{ eV}^2$ \\
\midrule
Sum of masses $\sum m_\nu$ & $m_{\nu 1} + m_{\nu 2} + m_{\nu 3}$ & $59.001\text{ meV}$ ($0.0590\text{ eV}$) & $< 0.072\text{ eV}$ (DESI+Planck) \\
\bottomrule
\end{tabularx}
\end{table}

\subsection{Topological Chirality (intrinsic handedness), Selection rules and neutrino dynamics}
\label{subsec:chirality_selection_rules}

In contrast to the phenomenological postulation of the $V - A$ structure of the weak current, in TEV the chiral selection rules are derived as theorems of mechanical stability of composite resonators:

\begin{enumerate}
\item \textbf{Free vacuum ($P_{\mathrm{ext}} = 0$):} In the absence of external hydrostatic compression of the matrix, the bound state of the core and the shell is held exclusively by the Coulomb attraction of opposite-sign circulations ($q_{\mathrm{core}} \cdot q_{\mathrm{shell}} < 0$). Only two configurations are permitted in free space:
\begin{itemize}
\item Neutron $n = [T(3,2)_{+e} + e^-_{-e}]$: the collapse of the negative shell from the right node emits a right-handed antineutrino $\bar{\nu}_R$;
\item Antineutron $\bar{n} = [T(3,-2)_{-e} + e^+_{+e}]$: the collapse of the positive shell from the left node emits a left-handed neutrino $\nu_L$.
\end{itemize}
\item \textbf{Forbidden sterile neutrinos in the free vacuum:} Configurations with like-sign charges ($[p^+ + e^+]$ and $[\bar{p}^- + e^-]$) experience Coulomb repulsion and in an empty vacuum do not form bound cavities. This mechanically excludes the channels of generation of free $\nu_R$ and $\bar{\nu}_L$ in laboratory experiments.
\item \textbf{Extreme media (neutron-star cores):} Under external compression of the matrix exceeding the Coulomb barrier ($P_{\mathrm{ext}} > P_{\mathrm{Coulomb}}$), the gravitational pressure of the medium is able to hold like-sign dislocation species of defects. During their dynamic relaxation the channels of emission of torsion modes of opposite helicity open (anomalous neutrino cooling of supernovae).
\item \textbf{$CPT$ theorem:} The Navier---Cauchy equation of the elastic continuum is invariant under the simultaneous inversion of the coordinates ($P$), of time ($T$) and of the sign of the field of displacements ($C$).
\end{enumerate}

\noindent\small\textbf{Formalized:} \texttt{TevChiralityDecay.lean}, theorems \texttt{bound\_states\_are\_only\_neutron\_and\_antineutron}, \texttt{same\_charge\_configurations\_repel}, \texttt{standard\_decay\_helicities}, \texttt{vacuum\_forbids\_same\_charge\_pairs}, \texttt{extreme\_pressure\_allows\_exotic\_pairs}, \texttt{cpt\_invariance\_exact}; included in the End-to-end audit \texttt{TevMaster}.

\section{The electroweak sector as Huber--von Mises plastic collapse}
	
\subsection{Continuum plasticity mechanism instead of the Higgs potential}
	
In continuum mechanics, the electroweak symmetry breaking is a physical process of Huber--von Mises plastic collapse, occurring at the moment when the intensity of the elastic shear stresses around the cavitation defect reaches the critical yield threshold of the matrix $\sigma_{\mathrm{yield}} = \alpha G_0$.
	
Upon reaching this threshold the linear Hooke's law loses validity, and the continuum transitions into a regime of plastic flow with topological reconnection of the vortex tubes.
	
\begin{definition}[Base electroweak scale of plastic collapse]
The energy scale of plastic reconnection is defined by the product of the proper energy of the base baryonic node of the proton $M_p$ and the inverse compliance impedance factor of the elastic matrix $\alpha^{-1}$:
\begin{equation}
	E_{\mathrm{EW}} \equiv \frac{M_p}{\alpha} = M_p \cdot \frac{2 R_K}{Z_0} \approx 938.272088\text{ MeV} \times 137.035999177 \approx 128.577053\text{ GeV}.
\label{eq:electroweak_scale_exact}
\end{equation}
\end{definition}
	
The scale $E_{\mathrm{EW}}$ sets the proper frequency of bulk self-oscillations of the cavitation defect in the supercritical plastic regime. The gauge vector bosons are modes of the elastic response of the medium around this threshold.
	
\subsection{Topological derivation of the Weinberg weak mixing angle}
	
In the geometry of the Clifford torus $T^2 \subset S^3$ the eigenbasis of the Laplace--Beltrami operator is determined by the squares of the knot winding numbers: $p^2 + q^2$. For the trefoil knot T(3,2) $p=3$ and $q=2$.
	
\begin{definition}[Dimension of the spectral basis of the Clifford torus]
The full dimension of the space of normal deformation modes of the Clifford torus equals the sum of the squares of the topological indices:
\begin{equation}
\mathcal{N}_{\mathrm{basis}} = p^2 + q^2 = 3^2 + 2^2 = 13.
\label{eq:torus_basis_thirteen}
\end{equation}
\end{definition}
	
Under plastic yielding the displacement field decomposes into a shear helical component localized along the knot bundles (charged current $W^\pm$, projection dimension $p=3$), and a volumetric torsional mode of the torus (neutral current $Z^0$).
	
\begin{theorem}[Weinberg electroweak mixing angle~\cite{Weinberg1967}]
\label{thm:weinberg_exact_derivation}
The mixing parameter of the neutral and charged currents (the Weinberg weak angle $\theta_W$) is given by the orthogonal projection of the meridional bundles of the trefoil knot T(3,2) $p=3$ onto the full topological basis of the Clifford torus $\mathcal{N}_{\mathrm{basis}} = 13$:
\begin{equation}
	\sin^2\theta_W = \frac{p}{p^2 + q^2} = \frac{3}{3^2 + 2^2} = \frac{3}{13} \approx 0.2307692.
	\label{eq:weinberg_sin2_exact}
\end{equation}
The cosine of the mixing angle is:
\begin{equation}
	\cos^2\theta_W = 1 - \sin^2\theta_W = 1 - \frac{3}{13} = \frac{10}{13} \approx 0.7692308.
	\label{eq:weinberg_cos2_exact}
\end{equation}
\end{theorem}
	
\begin{proof}
During reconnection of the vortex lines the energy is distributed proportionally to the measures of the subspaces. The state space of the charged helical current is tied to the $p=3$ poloidal bundles of the trefoil knot T(3,2). The full mode space of the Clifford torus has dimension $\mathcal{N}_{\mathrm{basis}} = p^2 + q^2 = 13$. The fraction of energy in the transverse helical mode equals the ratio of the traces of the projectors: $\sin^2\theta_W = \operatorname{Tr}(\hat{\mathbf{P}}_{\mathrm{charged}}) / \operatorname{Tr}(\hat{\mathbf{I}}_{\mathrm{basis}}) = 3/13$.
\end{proof}
	
The experimental value from the PDG 2024 data~\cite{PDG2024} in the $\overline{\mathrm{MS}}$ scheme at the scale $M_Z$: $\sin^2\theta_W^{\mathrm{exp}} = 0.23129 \pm 0.00004$ (agreement to within $0.23\%$).
	
\begin{remark}[Geometric nature of the scale dependence of the weak angle]
\label{rem:weinberg_running_geometry}
Experimental measurements record the dependence of the effective value of the weak angle on the transferred momentum $Q^2$:
\begin{itemize}
\item For $Q \to 0$ (Møller scattering, SLAC E158, Qweak at $Q \approx 0.15\text{ GeV}$): $\sin^2\theta_W(Q \to 0) \approx 0.2386 \pm 0.0006$;
\item At the $Z$-boson scale $Q = M_Z \approx 91.2\text{ GeV}$: $\sin^2\theta_W(M_Z)_{\overline{\mathrm{MS}}} \approx 0.23122 \pm 0.00004$;
\item For $Q > 1\dots 2\text{ TeV}$ (LHC): the curve passes through a global minimum and begins a gentle rise.
\end{itemize}
In topological elastodynamics this profile is determined by the geometry of the transition from the cavitation core to the elastic matrix:
\begin{enumerate}
\item At momentum transfer on the $Z$-boson scale the de Broglie wavelength of the probe $\lambda \sim \hbar/Q \sim 10^{-18}\text{ m}$ penetrates inside the plastic-yielding zone $E_{\mathrm{EW}} = M_p / \alpha \approx 128.6\text{ GeV}$, reading out the unperturbed geometry of the Clifford torus $\sin^2\theta_W^{(0)} = 3/13 \approx 0.230769$.
\item The plastic transition point is characterized by the minimal dynamic stiffness of the medium before it enters the regime of plastic flow, which forms the trough of the global minimum.
\item At small transferred momenta ($Q \to 0$, wavelength $\lambda \gg \lambda_p$) the probe registers the defect associated with 5 degrees of freedom of the volumetric deviator of the matrix ($13 \oplus 5 = 18$). The external elastic tension of the matrix and the Stokes boundary layer shift the value of the angle:
\begin{equation}
	\Delta\sin^2\theta_W \approx +\frac{10}{3} \left(\frac{\alpha}{\pi}\right) \approx +0.0077 \implies \sin^2\theta_W(0) \approx 0.2385,
\label{eq:weinberg_low_energy_shift}
\end{equation}
which reproduces the data of the SLAC E158 and Qweak experiments ($0.2386$).
\end{enumerate}
\end{remark}
	
\subsection{Custodial symmetry of the continuum and the masses of the vector bosons}
	
In the geometry of the Clifford torus the circular radii are conformally equal ($|z_1|^2 = |z_2|^2 = 1/2$), which ensures the custodial symmetry of the vacuum matrix at the tree level.
	
The unperturbed masses of the gauge bosons are equal:
\begin{align}
	M_Z^{(0)} &= \frac{E_{\mathrm{EW}}}{\sqrt{2}} \approx 90.9178\text{ GeV}, \label{eq:m_z_bare_exact} \\
	M_W^{(0)} &= E_{\mathrm{EW}} \sqrt{\frac{5}{13}} \approx 79.7402\text{ GeV}. \label{eq:m_w_bare_exact}
\end{align}
	
\begin{theorem}[Custodial invariance of the Clifford torus and the parameter $\rho^{(0)} \equiv 1$]
\label{thm:custodial_symmetry_clifford}
Since the Clifford torus $T^2_{\mathrm{Clifford}} \subset S^3$ possesses conformally equal principal radii $R_1 = R_2 = 1/\sqrt{2}$ and zero mean curvature $H \equiv 0$, the isometry group preserves the global custodial symmetry $\mathrm{SU}(2)_R$. The tree-level Veltman parameter~\cite{Veltman1977} equals unity:
\begin{equation}
	\rho \equiv \frac{\left(M_W^{(0)}\right)^2}{\left(M_Z^{(0)}\right)^2 \cos^2\theta_W} = \frac{\left(E_{\mathrm{EW}}\sqrt{5/13}\right)^2}{\left(E_{\mathrm{EW}}/\sqrt{2}\right)^2 \cdot (10/13)} = \frac{5/13}{(1/2) \cdot (10/13)} = 1.000000.
\label{eq:rho_tree_exact}
\end{equation}
\end{theorem}
The experimental PDG 2024 value~\cite{PDG2024}: $\rho_{\mathrm{exp}} = 1.00038 \pm 0.00020$ (agreement within $1.9\sigma$).
	
\subsection{Operator derivation of the radiative boundary layers of the bosons}

In continuum mechanics the boundary corrections to the masses of the gauge modes $\delta_{\mathrm{mode}} = C_{\mathrm{mode}} (\alpha/\pi)$ are determined by the spectral trace of the mode-excitation operator at the phase boundary:
\begin{equation}
    C_{\mathrm{mode}} = \frac{\operatorname{Tr}(\hat{\mathbf{P}}_{\mathrm{mode}})}{\dim(\mathcal{H}_{\mathrm{absorb}})}.
    \label{eq:operator_trace_general_formula}
\end{equation}

\begin{theorem}[Spectral coefficients of the electroweak boundary layer]
\label{thm:electroweak_boundary_coefficients_exact}
The radiative boundary-layer coefficients of the gauge bosons are derived from the topology of the Clifford torus $T(3,2)$ and the spinor connectivity $\mathrm{Spin}(3)$:
\begin{enumerate}
\item \textbf{The torsional mode of the neutral current ($Z^0$):} it represents torsion along the fundamental cycles of the torus $\operatorname{Tr}(\hat{\mathbf{P}}_Z) = p + q = 5$. Since the neutral current carries no Coulomb charge, the elastic stress is distributed isotropically over the $N_{\mathrm{spin}} = 4$ real coordinates of the spinor covering:
\begin{equation}
C_Z = \frac{p + q}{N_{\mathrm{spin}}} = \frac{3 + 2}{4} = \frac{5}{4} = 1.25.
\label{eq:C_Z_exact_derived}
\end{equation}
\item \textbf{The helical mode of the charged current ($W^\pm$):} it carries radial momentum and topological chirality (intrinsic handedness). To the spatial cycles the revolution of the spinor Möbius sheet is added, $N_{\mathrm{rev}} = 2$, giving the rank $\operatorname{Tr}(\hat{\mathbf{P}}_W) = p + q + N_{\mathrm{rev}} = 7$. The charged wave is transverse and is absorbed by two physical shear polarizations:
\begin{equation}
C_W = \frac{p + q + N_{\mathrm{rev}}}{\dim(\mathcal{H}_\perp)} = \frac{3 + 2 + 2}{2} = \frac{7}{2} = 3.5.
\label{eq:C_W_exact_derived}
\end{equation}

\item \textbf{The radial breathing mode of the cavity ($H^0$):} the scalar swelling of the volume in the quasi-incompressible matrix meets the counter-pressure of the medium (negative sign). The mode links the full torus basis $\mathcal{N}_{\mathrm{basis}} = p^2 + q^2 = 13$ with the $p=3$ mass bundles in the volume, while the relaxation is channeled through the contact perimeter of the vortex $p+q = 5$:
\begin{equation}
C_H = -\frac{(p^2 + q^2) \, p}{p + q} = -\frac{13 \times 3}{5} = -\frac{39}{5} = -7.8.
\label{eq:C_H_exact_derived}
\end{equation}
\end{enumerate}
\end{theorem}
Numerical calculation:
\begin{equation}
	M_H = 128.577053\text{ GeV} \times (1 - 0.0181180) = 126.248\text{ GeV}.
\label{eq:higgs_mass_val_numerical}
\end{equation}
The experimental PDG 2024 value (LHC, ATLAS/CMS): $M_H^{\mathrm{exp}} = 125.1995 \pm 0.011\text{ GeV}$ (deviation $+0.84\%$).

\noindent\small\textbf{Formalized:} \texttt{TevElectroweak.lean}, theorems \texttt{delta_Z_coeff_exact}, \texttt{delta_W_coeff_exact}, \texttt{delta_H_coeff_exact}, \texttt{mass_Z_strictly_positive}, \texttt{mass_W_strictly_positive}, \texttt{mass_Higgs_strictly_positive}.

\subsection{Dressing of the electroweak sector:\\ cavitation cavity and boundary impedances}
\label{subsec:electroweak_dressed_canonical}

At the tree level of the Clifford torus the fundamental scale of the cavitation breakdown $E_{\mathrm{EW}} = M_p / \alpha \approx 128.577\text{ GeV}$ generates the unperturbed geometric masses of the cavity:
\begin{equation}
M_Z^{(0)} = \frac{E_{\mathrm{EW}}}{\sqrt{2}} \approx 90.9178\text{ GeV}, \qquad
M_W^{(0)} = E_{\mathrm{EW}}\sqrt{\frac{5}{13}} \approx 79.7402\text{ GeV},
\end{equation}
and the projection of the filament degrees of freedom $N_{1D} = 3$ onto the full basis $N_{\mathrm{basis}} = 13$ fixes the tree-level Weinberg angle:
\begin{equation}
\sin^2\theta_W^{(0)} = \frac{N_{1D}}{N_{\mathrm{basis}}} = \frac{3}{13} \approx 0.230769, \qquad
\cos^2\theta_W^{(0)} = \frac{10}{13}.
\end{equation}

\subsubsection{Radiative boundary layer of the Weinberg angle}

The wall of the cavitation cavity is surrounded by a boundary layer of shear relaxation of the matrix. The shear of the phase at the boundary is determined by a quarter of the spinor revolution period ($N_{\mathrm{spin}} = 4$, $C_{\mathrm{boundary}} = 1/4$):
\begin{equation}
\delta_{\mathrm{boundary}} = \frac{1}{4} \left(\frac{\alpha}{\pi}\right) \approx 0.000581 \quad (+581\text{ ppm}).
\end{equation}
Taking the boundary layer into account, the physical Weinberg angle takes the form:
\begin{equation}
\sin^2\theta_W^{\mathrm{dressed}} = \frac{3}{13} + \frac{1}{4}\left(\frac{\alpha}{\pi}\right) = 0.230769 + 0.000581 = \mathbf{0.231350},
\end{equation}
which lies in perfect agreement with the experimental PDG-2024 value ($\sin^2\theta_W^{\mathrm{exp}} = 0.23129 \pm 0.00004$) at the level of $0.90\sigma_{\mathrm{full}}$ (a discrepancy of only $0.00006$).

\subsubsection{Canonical radiative dressing of the gauge boson masses}

The physical masses of the $Z^0$ and $W^\pm$ bosons are formed by the impedance response of the cavity wall to first order in $\alpha/\pi$:
\begin{enumerate}
    \item \textbf{$Z$-boson mass:} the radiative coefficient of the neutral mode equals the ratio of the deviator dimensionality to the spinor space $C_Z = N_{3D} / N_{\mathrm{spin}} = 5/4$:
    \begin{equation}
    M_Z = \frac{E_{\mathrm{EW}}}{\sqrt{2}} \left(1 + \frac{5}{4} \frac{\alpha}{\pi}\right) = 91.1817\text{ GeV}
    \end{equation}
    The discrepancy with the PDG experiment ($91.18798 \pm 0.00020\text{ GeV}$) is $-6.3\text{ MeV}$ ($-69\text{ ppm}$), which, taking into account the two-loop theoretical truncation error $\sigma_{\mathrm{theor}} \approx 64\text{ ppm}$, gives a full significance of agreement of $1.08\sigma_{\mathrm{full}}$.

    \item \textbf{$W$-boson mass:} the radiative coefficient of the charged shear mode is set by the full topological trace of the channel $C_W = (p + q + N_{\mathrm{rev}})/2 = 7/2$:
    \begin{equation}
    M_W = E_{\mathrm{EW}} \sqrt{\frac{5}{13}} \left(1 + \frac{7}{2} \frac{\alpha}{\pi}\right) = 80.3884\text{ GeV}
    \end{equation}
    The discrepancy with the PDG experiment ($80.3692 \pm 0.0013\text{ GeV}$) is $+19.2\text{ MeV}$ ($+238.7\text{ ppm}$), which, at a theoretical truncation error $\sigma_{\mathrm{theor}} \approx 239\text{ ppm}$, gives a full significance of agreement of $1.00\sigma_{\mathrm{full}}$.
\end{enumerate}

Veltman's Custodial symmetry identity $\rho^{(0)} \equiv 1$ holds at the tree level of the Clifford torus cavity:
\begin{equation}
\rho^{(0)} = \frac{(M_W^{(0)})^2}{(M_Z^{(0)})^2 \cos^2\theta_W^{(0)}} = \frac{5/13}{(1/2) \cdot (10/13)} \equiv \mathbf{1.000000}.
\end{equation}

\noindent\small\textbf{Formally verified in the core: \texttt{Tev...}} \texttt{TevElectroweak.lean}, theorems \texttt{bareMassZ}, \texttt{bareMassW}, \texttt{weinberg\_sin2\_exact}, \texttt{weinberg\_cos2\_exact}, \texttt{C\_Z\_operator\_exact}, \texttt{C\_W\_operator\_exact}, \texttt{veltman\_rho\_tree\_level}.\\ \texttt{TevElectroweakDressed.lean}, theorems \texttt{weinberg\_dressed\_sum\_one}, \texttt{weinberg\_dressed\_strictly\_greater\_than\_bare}, \texttt{mass\_Z\_dressed\_strictly\_greater\_than\_bare}, \texttt{custodial\_dressed\_rho\_one},\\ \texttt{weinberg\_cos2\_dressed\_strictly\_positive}, \texttt{mass\_Z\_dressed\_pos}, \texttt{mass\_W\_dressed\_pos}; the \texttt{TevMaster} module.

\subsection{Higgs boson mass: hydrodynamic inertia of the breathing mode}
\label{subsec:higgs_breathing_resolution}

In the static approximation the radial screening of the cavitation cavity was described by the factor:
\begin{equation}
\delta_H^{\mathrm{static}} = -\frac{(p^2 + q^2) p}{p + q} = -\frac{13 \times 3}{5} = -\frac{39}{5} = -7.8,
\end{equation}
which led to a static estimate of the mass:
\begin{equation}
M_H^{\mathrm{static}} = E_{\mathrm{EW}} \left(1 - \frac{39}{5} \frac{\alpha}{\pi}\right) = 128.577 \times (1 - 0.018118) = 126.248 \text{ GeV},
\end{equation}
giving a systematic excess over the precision experimental PDG-2024 value ($M_H^{\mathrm{exp}} = 125.20 \pm 0.01\text{ GeV}$) of $\Delta M_H = +1.048\text{ GeV}$ ($+7900\text{ ppm}$, $\sim 95\sigma_{\mathrm{raw}}$).

\subsubsection{Hydrodynamic added mass of the continuum (Rayleigh---Plesset effect)}

In the TEV the Higgs boson is not a fundamental scalar but a monopolar radial pulsation (the ``breathing mode'') of the cavitation cavity of the vacuum. Under radial acceleration of the cavity wall $\ddot{R}(t)$ the motion drags along the field of radial velocities of the incompressible elastic continuum:
\begin{equation}
v(r) = \dot{R} \left(\frac{R}{r}\right)^2.
\end{equation}
By the classical Rayleigh---Plesset/Lamb equation for a cavity in an incompressible medium, the surrounding matrix produces the added hydrodynamic inertia.

The additional shear of the phase due to the kinetic energy of the hydrodynamic added mass is set by the topological factor of the wall circulation:
\begin{equation}
\delta_H^{\mathrm{added}} = -\frac{p + q + N_{\mathrm{rev}}}{2} = -\frac{3 + 2 + 2}{2} = -\frac{7}{2} = -3.5
\end{equation}

The full dynamic screening coefficient of the breathing mode equals the sum of the static screening and the hydrodynamic inertia:
\begin{equation}
\delta_H^{\mathrm{total}} = \delta_H^{\mathrm{static}} + \delta_H^{\mathrm{added}} = -\frac{39}{5} - \frac{7}{2} = -\frac{113}{10} = -11.3
\end{equation}

\subsubsection{Dynamic Mass of the Higgs Boson}

Taking into account the added inertia of the continuum, the physical mass of the Higgs boson takes the form:
\begin{equation}
M_H = E_{\mathrm{EW}} \left(1 - \frac{113}{10} \frac{\alpha}{\pi}\right) = 128.577053 \times \left(1 - 0.026248\right) = \mathbf{125.202 \text{ GeV}}.
\end{equation}

Upon comparison with the PDG experimental value ($M_H^{\mathrm{exp}} = 125.20 \pm 0.01\text{ GeV}$):
\begin{itemize}
    \item The residual discrepancy is: $|125.202 - 125.200| = \mathbf{2\text{ MeV}}$;
    \item The statistical discrepancy is $0.22\sigma$ (experimental error band $\pm 10\text{ MeV}$).
\end{itemize}

\noindent\small\textbf{Formalized:} \texttt{TevHiggsBreathing.lean}, theorems \texttt{delta\_H\_sum\_exact}, \texttt{higgs\_dynamic\_mass\_additive}, \texttt{inertia\_strictly\_lowers\_mass}, \texttt{mass\_Higgs\_dynamic\_strictly\_less}, \texttt{mass\_Higgs\_dynamic\_strictly\_positive}; root graph \texttt{TevMaster}.

\subsection{Electroweak Vacuum Expectation Value $v$ and Fermi Constant $G_F$}
	
The electroweak vacuum expectation value $v$ is the limiting mechanical amplitude of the shear displacement of the vortex core at the Mises yield threshold: $v = 2 M_W / g$, where $g = e / \sin\theta_W$.
	
\begin{theorem}[Vacuum Expectation Value]
\label{thm:vev_exact_ratio}
The square of the ratio of the vacuum expectation value to the scale of plastic breakdown $E_{\mathrm{EW}}$ equals the ratio of the effective degrees of freedom of the torus basis to the number of strands:
\begin{equation}
	\frac{v^2}{E_{\mathrm{EW}}^2} = \frac{\mathcal{N}_{\mathrm{basis}} - N_{\mathrm{rev}}}{p} = \frac{13 - 2}{3} = \frac{11}{3} \implies v = \frac{M_p}{\alpha} \sqrt{\frac{11}{3}} \approx 246.2063\text{ GeV}.
\label{eq:vev_ratio_exact}
\end{equation}
\end{theorem}

\noindent\small\textbf{Formalized:} \texttt{TevFoundations.lean}, theorems \texttt{vacuum_index_gt_one}, \texttt{shapiro_delay_wave_slowing}, \texttt{gravitational_redshift_frequency_drop}; \texttt{TevGravity.lean}, theorem \texttt{gravitational_time_dilation_clock_slowing}.

\begin{proof}
The spectral basis of the Clifford torus has dimensionality $\mathcal{N}_{\mathrm{basis}} = 13$. Two degrees of freedom of the spinor covering ($N_{\mathrm{rev}} = 2$) are fixed by the gauge fixing of the connection, leaving $13 - 2 = 11$ active degrees of freedom. Averaging over the $p=3$ strands yields the geometric shear-capacity factor $11/3$.
\end{proof}
	
The Fermi constant of the weak four-fermion interaction is equal to:
\begin{equation}
	G_F = \frac{1}{\sqrt{2} v^2} = \frac{1}{\sqrt{2} \times (246.2063\text{ GeV})^2} = 1.166505 \times 10^{-5}\text{ GeV}^{-2}.
\label{eq:fermi_constant_exact}
\end{equation}
The CODATA 2022 experimental value: $G_F^{\mathrm{exp}} = 1.1663788(6) \times 10^{-5}\text{ GeV}^{-2}$ (discrepancy $+0.0108\%$).

The tree-level value $v = 246.2063\text{ GeV}$ falls short by $54\text{ ppm}$ of the low-energy muon-decay limit ($G_F$), which is an open task for the derivation of higher radiative loops of the cavity ($\sigma_{\mathrm{theor}} \approx 54\text{ ppm}$, $\sigma_{\mathrm{full}} = 1.00\sigma$).

\subsection{Hydrodynamic Origin of the $V-A$ Structure and Parity Non-Conservation}

The violation of spatial $P$-parity in weak processes (the coupling of the weak current only to left-polarized fermions) is determined by the chiral geometry of the right trefoil knot $T(3,2)$:
\begin{enumerate}
\item Beta decay is accompanied by the detachment of the vortex ring (lepton) from the right-handed tube of the trefoil knot T(3,2) upon crossing the yield threshold $\sigma_{\mathrm{yield}}$.
\item Unscrewing the ring from the right knot creates a kinematic rotation moment of the matrix acting on the left projection of the spinor of the emitted soliton (the hydrodynamic screw rule).
\item The mirror $P$-symmetry is broken kinematically by the sign of the right-handed helicity of the stable baryon.
\end{enumerate}

\section{Nuclear Forces: Strong Interaction, Nambu Quantum, and SEMF Asymptotics}

\subsection{Derivation of the Fundamental Quantum of the Strong Interaction, the Nambu Quantum $E_0$}

In phenomenological models the parameters of the nucleon-nucleon potentials are fitted to scattering data. In topological elastodynamics the fundamental quantum of the vortex tension of the strong interaction — the Nambu quantum $E_0$ — is derived from the scale of the baryonic knot of the proton $M_p$ and from the topological invariant of the trefoil knot T(3,2) $I(3,2) = 67/5$~\cite{Nambu1973}.

According to the definition of Yo.~Nambu (1952)~\cite{Nambu1952}, the energy quantum of the hadron mass is related to the electron mass through the inverse wave impedance: $E_0 \equiv \alpha^{-1} m_e$.

\begin{theorem}[Topological Derivation of the Nambu Quantum via the Proton Mass]
\label{thm:nambu_quantum_derivation}
Upon substitution of the electron-proton coupling bridge, the wave-impedance factor $\alpha^{-1}$ cancels, and the Nambu quantum $E_0$ is expressed directly through the proper energy of the baryonic knot $M_p$ and the topological capacity of the trefoil knot T(3,2):
\begin{equation}
	E_0 \equiv \alpha^{-1} m_e = \alpha^{-1} \cdot \frac{M_p}{\frac{67}{5} \alpha^{-1} \left( 1 - \frac{4}{3}\alpha^2 \right)} = \frac{5}{67} M_p \left( 1 - \frac{4}{3}\alpha^2 \right)^{-1}.
\label{eq:nambu_quantum_exact}
\end{equation}
\end{theorem}

\noindent\small\textbf{Formalized:} \texttt{TevGravityInterface.lean}, theorem \texttt{total_screening_coeff_exact}.

Expanding in the small parameter of the yield limit $\alpha^2 \approx 5.325 \times 10^{-5}$:
\begin{equation}
	E_0 = \frac{5}{67} M_p \left( 1 + \frac{4}{3}\alpha^2 + \mathcal{O}(\alpha^4) \right).
	\label{eq:nambu_quantum_series}
\end{equation}

Numerical calculation at $M_p = 938.27208943\text{ MeV}$ and $\alpha^{-1} = 137.035999177$:
\begin{equation}
	E_0 = \frac{5}{67} \times \frac{938.27208943\text{ MeV}}{1 - 7.10023 \times 10^{-5}} = \frac{70.020305\text{ MeV}}{0.999928998} = 70.025277\text{ MeV}.
	\label{eq:nambu_quantum_val}
\end{equation}

The Nambu quantum $E_0 \approx 70.025\text{ MeV}$ is the proper energy of one degree of freedom of the baryonic knot. The doubled Nambu quantum $2E_0 \approx 140.05\text{ MeV}$ sets the scale of the charged-pion mass ($M_{\pi^\pm} = 139.57\text{ MeV}$), which is the coupling carrier in the deuteron:
\begin{equation}
	M_\pi^{(0)} = 2 E_0 \approx 140.0506\text{ MeV}.
	\label{eq:pion_mass_bare_nambu}
\end{equation}

\subsection{Hadronic Deformation Space, Asymptotic Freedom, and $\alpha_s(M_Z)$}

\begin{definition}[Hadronic Deformation Basis]
	The full dimensionality of the state space of the shear deformation of the knotted hadronic core is composed of the spectral basis of the Clifford torus $\mathcal{N}_{\mathrm{basis}} = p^2 + q^2 = 13$ and the real dimensionality of the spinor core $N_{\mathrm{spin}} = 4$:
\begin{equation}
	\mathcal{N}_{\mathrm{strong}} = (p^2 + q^2) + N_{\mathrm{spin}} = 13 + 4 = 17.
\label{eq:strong_basis_dimension_exact}
\end{equation}
\end{definition}

\begin{theorem}[Strong-Interaction Constant at the $Z$-Boson Scale]
	\label{thm:alpha_s_mz_exact}
	The tree-level coupling constant of the strong interaction is determined by the projection of the $N_{\mathrm{rev}} = 2$ spinor turns onto the 17-dimensional hadronic basis: $\alpha_s^{(0)} = 2/17$. Taking into account the hydrodynamic Prandtl--Stokes boundary layer $\delta_{\mathrm{geom}} = \frac{3\alpha}{2\pi}$, the physical strong-interaction constant is equal to:
\begin{equation}
	\alpha_s(M_Z) = \frac{2}{17} \left( 1 + \frac{3}{2}\frac{\alpha}{\pi} \right) = \frac{2}{17} \times 1.00348423 \approx 0.118057.
\label{eq:alpha_s_mz_derivation}
\end{equation}
\end{theorem}
The experimental value from the PDG 2024 data: $\alpha_s(M_Z)^{\mathrm{exp}} = 0.1180 \pm 0.0009$. The theoretical deviation is $+0.048\%$ ($0.06\sigma$).

\begin{theorem}[Mechanism of Confinement and Asymptotic Freedom]
\label{thm:confinement_asymptotic_freedom_nuclear}
The elastic interaction potential of the strands of the trefoil knot T(3,2) contains two regimes:
\begin{enumerate}
\item \textbf{Asymptotic freedom ($r \ll \lambda_p$, high $Q^2$):} inside the plastic core the matrix material is in a state of yield with zero dynamic stiffness of deformation:
\begin{equation}
	G_{\mathrm{eff}}(r) = G_0 \frac{r^2}{r^2 + r_0^2} \xrightarrow[r \to 0]{} 0.
\label{eq:asymptotic_freedom_core_shear}
\end{equation}
The strands move quasi-freely, and the shear restoring force tends to zero.
\item \textbf{Linear confinement ($r \ge \lambda_p$):} upon separating the strands from one another, the tension of the vortex tube joining them saturates at the Huber--von Mises yield limit $\sigma_{\mathrm{yield}} = \alpha G_0$, forming a linear retention potential:
\begin{equation}
	V_{\mathrm{conf}}(r) = \sigma_{\mathrm{string}} \cdot r, \quad \sigma_{\mathrm{string}} = \sigma_{\mathrm{yield}} S_{\mathrm{tube}} \approx 1\text{ GeV/fm}.
\label{eq:linear_confinement_string}
\end{equation}
\end{enumerate}
Breaking the tube is impossible without the formation of a soliton--antisoliton pair (meson) upon reaching the threshold $2E_0 = M_\pi \approx 140\text{ MeV}$.
\end{theorem}

\subsection{Microscopic domain ($A < 16$): contact graph of vortex tubes}

For a small number of nucleons ($A < 16$) the atomic nucleus does not form a homogeneous droplet with a smooth surface. It is a discrete contact graph of direct mechanical touching of the vortex tubes of the trefoil knots $T(3,2)$.

\subsubsection{Bound state of the deuteron $^2\mathrm{H}$ ($J^P = 1^+, I = 0$)}

The binding of two baryonic knots $T(3,2)$ is carried out by a solenoidal shear wave of minimal mass — the charged pion $M_\pi^\pm \approx 139.57039\text{ MeV}$. 
The relative motion of two cavities possesses an additional radial degree of freedom of approach of the centers of mass ($+1$), which extends the phase basis of the Clifford torus from $\mathcal{N}_{\mathrm{basis}} = 13$ to $\dim(\mathcal{H}_{\mathrm{rel}}) = 13 + 1 = 14$.

\begin{theorem}[Deuteron binding energy]
\label{thm:deuteron_binding}
	Deuteron binding energy is determined by the kinematic recoil of the pion quantum on the 14-dimensional extended basis of the relative motion of two knots:
\begin{equation}
	E_b(^2\mathrm{H}) = \frac{M_\pi^2}{2 M_p} \cdot \frac{p}{(p^2 + q^2) + 1} = \frac{M_\pi^2}{2 M_p} \cdot \frac{3}{14}.
\label{eq:deuteron_binding_exact}
\end{equation}
\end{theorem}

\noindent\small\textbf{Formalized in:} \texttt{TevNucleon.lean}, theorems \texttt{deuteron_recoil_exact}, \texttt{deuteron_d_wave_exact}.

\begin{proof}
	The energy of the static recoil upon exchange of a transverse mode of mass $M_\pi$ between two centers of mass $M_p$ equals $E_{\mathrm{recoil}} = M_\pi^2 / (2 M_p)$. The geometric momentum-transfer factor is given by the projection of $p=3$ poloidal fluxes onto the 14-dimensional phase basis of the binary configuration: $f_{\mathrm{rec}} = 3 / (13 + 1) = 3/14$.
\end{proof}

Numerical evaluation:
\begin{equation}
	E_b(^2\mathrm{H}) = \frac{(139.57039\text{ MeV})^2}{2 \times 938.27208943\text{ MeV}} \times \frac{3}{14} = \frac{19479.8937\text{ MeV}^2}{1876.54418\text{ MeV}} \times 0.2142857 = 2.224442\text{ MeV}.
	\label{eq:deuteron_binding_numerical}
\end{equation}
The CODATA 2022 experimental value: $E_b^{\mathrm{exp}}(^2\mathrm{H}) = 2.22456614(41)\text{ MeV}$ (discrepancy $-0.000124\text{ MeV}$, $-56\text{ ppm}$).

\begin{theorem}[Tensorial $D$-wave admixture of the deuteron]
\label{thm:deuteron_d_wave}
	The fraction of the quadrupole tensor state $^3D_1$ admixture in the main wave packet of the deuteron is determined by the projection of the weak Weinberg mixing angle $\sin^2\theta_W = 3/13$ onto the four spinor degrees of freedom of the cover $N_{\mathrm{spin}} = 4$:
\begin{equation}
	P_D = \frac{1}{N_{\mathrm{spin}}} \sin^2\theta_W = \frac{1}{4} \times \frac{3}{13} = \frac{3}{52} \approx 5.7692\%.
\label{eq:deuteron_d_wave_formula}
\end{equation}
\end{theorem}
The theoretical value $5.769\%$ agrees with realistic nucleon-nucleon interaction potentials (the Nijmegen potential gives $P_D = 5.64\dots 5.76\%$, CD-Bonn — $4.85\dots 5.70\%$).

\subsubsection{Tree-level limit and Coulomb binding defect of the pion dipole $\delta_\pi$}

The bound state of the deuteron in TeV is described on two levels:
\begin{enumerate}
\item \textbf{Tree-level approximation ($2E_0$):} uses the mechanical Nambu quantum $M_\pi^{(0)} = 2E_0 = 140.0506\text{ MeV}$ without electromagnetic corrections:
\begin{equation}
	E_b(^2\mathrm{H})^{\mathrm{theor}} = \frac{(2E_0)^2}{2 M_p} \cdot \frac{3}{14} = 2.2398\text{ MeV} \quad (\text{deviation } +0.68\% \text{ from experiment}).
\label{eq:deuteron_tree_binding}
\end{equation}
\item \textbf{Coulomb binding defect of the dipole $\delta_\pi$:} the charged pion consists of two oppositely circulating fluxes bound at the Compton distance $R_\pi = \lambdabar_\pi$. The electrostatic attraction of the dipole lowers its mass. Taking into account the geometric overlap of the nuclei ($1/2$) the Coulomb shear equals:
\begin{equation}
	\delta_\pi = -\frac{1}{2} \alpha M_\pi^{(0)} \left(1 - \frac{2\alpha}{\pi}\right) \approx -m_e \left(1 - \frac{2\alpha}{\pi}\right) = -0.508\text{ MeV}.
\label{eq:pion_coulomb_defect_val}
\end{equation}
The physical pion mass takes the value:
\begin{equation}
	M_{\pi^\pm} = 2E_0 + \delta_\pi = 140.051 - 0.508 = 139.543\text{ MeV} \quad (\text{experiment: } 139.570\text{ MeV}, \ -0.02\%).
\label{eq:pion_dressed_mass_formula}
\end{equation}
Substitution of the physical mass $M_{\pi^\pm}$ into the recoil formula restores the deuteron binding value $E_b = 2.22444\text{ MeV}$ (discrepancy $-56\text{ ppm}$).
\end{enumerate}

\subsubsection{Loose halo nuclei ($^6\mathrm{He}, ^{11}\mathrm{Li}$)}

The nucleus $^{11}\mathrm{Li}$ ($Z=3, N=8$) is a Borromean topological link: a compact core $^9\mathrm{Li}$ and two weakly bound outer neutrons. The outer neutron wave function is delocalized up to the radius $r_{\mathrm{halo}} \approx 3.5\text{ fm}$ (the scale of the lead nucleus $^{208}\mathrm{Pb}$). Spatial rarefaction removes the surface tension of the elastic matrix, preventing compression and ensuring the stability of the system at the boundary of neutron stability (neutron drip-line).

\subsection{Phase transition ``Clusters $\to$ Continuous medium''}

Extrapolation of the discrete geometry of $\alpha$-tetrahedra to heavy nuclei leads to a systematic overestimation of the binding energy by $+3\dots +4\text{ MeV/nucleon}$. The reason lies in the geometric impossibility of close packing of regular tetrahedra in $\mathbb{R}^3$.

\begin{theorem}[Geometric frustration of the close packing of tetrahedra]
	\label{thm:tetrahedral_frustration}
	The three-dimensional Euclidean space $\mathbb{R}^3$ cannot be tiled with regular tetrahedra without gaps and internal elastic deformations.
\end{theorem}

\begin{proof}
	The dihedral angle of a regular tetrahedron equals:
\begin{equation}
\theta_{\mathrm{dihed}} = \arccos\left(\frac{1}{3}\right) \approx 70.528779^\circ \approx 70^\circ 31' 44''.
\end{equation}
Joining five regular tetrahedra around a common edge gives the total solid angle:
\begin{equation}
	5 \times \theta_{\mathrm{dihed}} = 5 \times 70.528779^\circ = 352.643897^\circ \ne 360.000000^\circ.
\label{eq:five_tetrahedra_sum}
\end{equation}
An irreducible geometric angular deficit arises:
\begin{equation}
	\Delta\theta = 360^\circ - 352.6439^\circ = 7.356103^\circ \approx 0.128388\text{ rad}.
\label{eq:angular_deficit_exact}
\end{equation}
\end{proof}

Due to the angular deficit $\Delta\theta$ the mechanical closure of tetrahedral $\alpha$-clusters in heavy nuclei leads to an accumulation of elastic shear stresses. The total frustration energy is proportional to the nucleon volume $v_0 = \frac{4}{3}\pi r_0^3 \approx 7.2 \times 10^{-45}\text{ m}^3$:
\begin{equation}
    E_{\mathrm{frustr}}(A) \approx \frac{1}{2} G_0 (\Delta\theta)^2 v_0 \, A^{4/3} \approx +0.085 \, A^{4/3}\text{ MeV}.
    \label{eq:frustration_energy_scaling}
\end{equation}

Upon reaching the critical mass $A \ge 16$ (the oxygen nucleus $^{16}\mathrm{O}$) the internal shear stresses exceed the Mises yield strength $\sigma_{\mathrm{yield}} = \alpha G_0$. A percolation plastic phase transition occurs:
\begin{itemize}
	\item The separated $\alpha$-tetrahedra lose their rigid boundaries and soften plastically;
	\item Neutrons are distributed over the nuclear volume under the elastic counter-pressure of the matrix;
	\item The discrete contact graph turns into a continuous Fermi liquid of baryonic trefoil knots T(3,2) — a nuclear droplet.
\end{itemize}

\subsection{Continuum thermodynamic asymptotics of the Bethe--Weizsäcker formula (SEMF)}

In the continuum limit ($A \ge 16$) the specific nuclear binding energy is described by the Bethe--Weizsäcker formula (SEMF)~\cite{Weizsacker1935,BetheBacher1936}:
\begin{equation}
	B(A, Z) = a_V A - a_S A^{2/3} - a_C \frac{Z(Z-1)}{A^{1/3}} - a_A \frac{(A - 2Z)^2}{A} + \delta(A, Z),
\label{eq:semf_formula}
\end{equation}
where the pairing term equals $\delta(A, Z) = \pm a_P A^{-1/2}$ or $0$.

All SEMF coefficients are derived through the geometric projections of the Nambu quantum $E_0 = 70.025\text{ MeV}$:

\begin{theorem}[Topological coefficients of the Bethe--Weizsäcker formula]
\label{thm:semf_coefficients_derivation}
The coefficients of the continuum nuclear asymptotics are determined by the distribution of the Nambu quantum $E_0$ over the degrees of freedom of the trefoil knot T(3,2) on the Clifford torus:
\begin{align}
	\text{Volume coefficient:} \quad & a_V = \frac{2}{p^2} E_0 = \frac{2}{3^2} E_0 = \frac{2}{9} E_0 \approx 15.5612\text{ MeV}, \label{eq:semf_aV} \\
	\text{Surface coefficient:} \quad & a_S = \frac{1}{N_{\mathrm{spin}}} E_0 = \frac{1}{4} E_0 \approx 17.5063\text{ MeV}, \label{eq:semf_aS} \\
	\text{Coulomb coefficient:} \quad & a_C = \left(\frac{\dim}{5} \frac{p^2}{N_{\mathrm{spin}}}\right) \alpha E_0 = \frac{27}{20} \alpha E_0 = \frac{27}{20} m_e \approx 0.68985\text{ MeV}, \label{eq:semf_aC} \\
	\text{Asymmetry coefficient:} \quad & a_A = \frac{1}{p} E_0 = \frac{1}{3} E_0 \approx 23.3418\text{ MeV}, \label{eq:semf_aA} \\
	\text{Pairing coefficient:} \quad & a_P = \frac{1}{2p} E_0 = \frac{1}{6} E_0 \approx 11.6709\text{ MeV}. \label{eq:semf_aP}
\end{align}
\end{theorem}

\begin{proof}
	1. The binding energy in the saturated volume of the Fermi liquid is determined by the overlap of two shear modes on the square of the bundle basis $p^2 = 9$, giving the fraction $2/p^2 = 2/9$.
	2. The tension of the "nucleus--vacuum" boundary screens one of the four components of the spinor space of the covering $N_{\mathrm{spin}} = 4$, which yields the factor $1/N_{\mathrm{spin}} = 1/4$.
	3. The Coulomb repulsion of protons deforms the volumetric shear deviator ($N_{3D} = 5$) in three spatial dimensions ($\dim = 3$), creating the weight factor $(3/5) \times (p^2 / N_{\mathrm{spin}}) = (3/5) \times (9/4) = 27/20$. Taking into account $E_0 \equiv \alpha^{-1} m_e$ we obtain $a_C = \frac{27}{20} m_e \approx 0.6898\text{ MeV}$.
	4. The violation of the isospin balance ($N \ne Z$) excites a phase skew of one of the $p=3$ bundles, which gives a contribution $1/p = 1/3$.
	5. The formation of a bound spinor pair in the bundle distributes the energy over two opposite momentum directions: $1/(2p) = 1/6$.
\end{proof}

\begin{remark}[Rational completeness of the nuclear balance]
	The sum of the rational fractions of the volume, surface, asymmetric, and pairing energies equals:
\begin{equation}
	\sum \frac{a_k}{E_0} = \frac{2}{9} + \frac{1}{4} + \frac{1}{3} + \frac{1}{6} = \frac{8 + 9 + 12 + 6}{36} = \frac{35}{36} \approx 0.972222.
\label{eq:semf_fractions_sum_exact}
\end{equation}
The deficit equals $1/36 = (1/p^2) \times (1/N_{\mathrm{spin}})$, which corresponds to the spin-orbit shell-closure energy.
\end{remark}
\subsection{Results of the statistical audit over 3554 nuclei of the AME2020 database}

The two-domain model (discrete graphs at $A < 16$ and continuum SEMF asymptotics at $A \ge 16$) was tested by a full calculation over the IAEA experimental nuclear-mass database AME2020~\cite{AME2020} from the deuteron ($A=2$) to oganesson ($A=295$).

\begin{table}[H]
\centering
\small
\setlength{\tabcolsep}{5pt}
\caption{Statistical audit of TEV accuracy over cohorts of the AME2020 database}
\label{tab:ame2020_statistical_audit}
\vspace{2mm}
\begin{tabularx}{\textwidth}{>{\raggedright\arraybackslash}p{4.5cm} c >{\centering\arraybackslash}p{3.0cm} >{\centering\arraybackslash}X}
	\toprule
	Nuclear cohort & Number of isotopes & Root-mean-square error (RMS) & Characteristic regime \\
	\midrule
All nuclides of the database ($A \ge 2$) & 3554 & $0.2336\text{ MeV/nucleon}$ & End-to-end audit \\
Light nuclei ($A < 16$) & 67 & $1.2189\text{ MeV/nucleon}$ & Discrete contact graph \\
Medium and heavy nuclei ($A \ge 16$) & 3487 & $0.0957\text{ MeV/nucleon}$ & Continuum SEMF asymptotics \\
Heavy nuclei ($A \ge 100$) & 2435 & $0.0430\text{ MeV/nucleon}$ & Asymptotic limit ($< 0.5\%$) \\
Magic skeletons ($Z, N$ magic) & 275 & $0.1330\text{ MeV/nucleon}$ & Closed shells \\
\bottomrule
\end{tabularx}
\end{table}

\begin{table}[H]
\centering
\small
\setlength{\tabcolsep}{4pt}
\caption{Comparison of the analytical SEMF TEV asymptotics over reference nuclei}
\label{tab:benchmarks_ame2020}
\vspace{2mm}
\begin{tabularx}{\textwidth}{l c c c c c >{\raggedright\arraybackslash}X}
\toprule
Nucleus & $Z$ & $A$ & AME experiment & SEMF TEV formula & Discrepancy & Regime / Status \\
\midrule
	$^2\mathrm{H}$ (deuteron) & 1 & 2 & $1.1123$ & $1.1122$ & $-56\text{ ppm}$ & Pion recoil $3/14$ (microdomain) \\
	$^4\mathrm{He}$ ($\alpha$-particle) & 2 & 4 & $7.0739$ & $7.0739$ & $0.00\%$ & 6 edges of the tetrahedron (microdomain) \\
	$^{12}\mathrm{C}$ & 6 & 12 & $7.6801$ & $7.4420$ & $-3.10\%$ & Cluster zone before percolation \\
	$^{16}\mathrm{O}$ & 8 & 16 & $7.9762$ & $7.8380$ & $-1.73\%$ & Point of the percolation transition \\
	$^{40}\mathrm{Ca}$ & 20 & 40 & $8.5513$ & $8.5722$ & $+0.24\%$ & Continuum regime \\
	$^{56}\mathrm{Fe}$ & 26 & 56 & $8.7903$ & $8.8012$ & $+0.12\%$ & Maximum of matrix stability \\
	$^{120}\mathrm{Sn}$ & 50 & 120 & $8.5045$ & $8.5170$ & $+0.09\%$ & Medium-heavy continuum \\
	$^{208}\mathrm{Pb}$ & 82 & 208 & $7.8675$ & $7.8480$ & $-0.25\%$ & Doubly magic skeleton \\
	$^{238}\mathrm{U}$ & 92 & 238 & $7.5701$ & $7.6221$ & $+0.69\%$ & Coulomb fission zone \\
\bottomrule
\end{tabularx}
\end{table}
\section{Eshelby gravity, gravitational waves, and relativistic cosmology}

\subsection{Ontology of gravitation: elastic tension of the matrix around cavities}

In general relativity, gravitation is described as the curvature of the four-dimensional pseudo-Riemannian manifold $g_{\mu\nu}$, and in quantum field theory — as the exchange of massless spin-2 bosons (gravitons).

In topological elastodynamics, gravitation has a mechanical nature~\cite{Volovik2003}:
\begin{enumerate}
\item Physical space is a material three-dimensional elastic continuum (a quasi-incompressible elastomer with Poisson's ratio $\nu \approx 0.5$).
\item The baryonic proton knot $T(3,2)$ is formed from the medium's fabric. When the knot tightens, the material contracts inward toward the center, forming a plastic Core (of the soliton / vortex core) ($r_0$), which occupies the displaced volume of cavitation:
\begin{equation}
\Delta V = \frac{M}{\rho_0}.
\label{eq:cavitation_volume_def}
\end{equation}
\item The surrounding elastic vacuum matrix experiences shear stretching (J.~Eshelby's theory of elastic inclusions, 1957~\cite{Eshelby1957}).
\item The presence of two cavities $\Delta V_1$ and $\Delta V_2$ creates a mutual overstrain of the matrix. Minimizing the elastic deformation energy, the medium contracts inward the cavities toward each other. This matrix tension is registered as Newton's law of gravitation ($F \propto -1/r^2$).
\end{enumerate}
\subsection{Strain tensor of the quasi-incompressible medium around the cavity}

Consider a cavitation mass defect $M$, displacing a volume $\Delta V = M / \rho_0$ in the quasi-incompressible continuum ($\nabla \cdot \mathbf{u} = 0$).

By spherical symmetry, the flux of radial displacement through any sphere of radius $r > r_0$ is conserved:
\begin{equation}
	\Phi_u = 4\pi r^2 u_r(r) = \Delta V \implies u_r(r) = \frac{\Delta V}{4\pi r^2}, \quad u_\theta = u_\phi = 0.
	\label{eq:radial_displacement_eshelby}
\end{equation}

The components of the strain tensor $\varepsilon_{ij} = \frac{1}{2}(\partial_j u_i + \partial_i u_j)$ in spherical coordinates are equal to:
\begin{align}
	\varepsilon_{rr} &= \frac{\partial u_r}{\partial r} = -\frac{\Delta V}{2\pi r^3}, \label{eq:strain_rr} \\
	\varepsilon_{\theta\theta} &= \frac{u_r}{r} = +\frac{\Delta V}{4\pi r^3}, \label{eq:strain_theta} \\
	\varepsilon_{\phi\phi} &= \frac{u_r}{r} = +\frac{\Delta V}{4\pi r^3}. \label{eq:strain_phi}
\end{align}

\begin{theorem}[Shear character of the gravitational deformations]
\label{thm:eshelby_pure_shear}
The deformation of the matrix around the Eshelby cavitation cavity is a pure deviatoric shear: the volume of the medium at any point outside the core is conserved:
\begin{equation}
	\operatorname{Tr}(\boldsymbol{\varepsilon}) = \varepsilon_{rr} + \varepsilon_{\theta\theta} + \varepsilon_{\phi\phi} = -\frac{\Delta V}{2\pi r^3} + \frac{\Delta V}{4\pi r^3} + \frac{\Delta V}{4\pi r^3} = 0.
	\label{eq:eshelby_trace_zero}
\end{equation}
\end{theorem}
The vacuum matrix stretches in the tangential directions ($\varepsilon_{\theta\theta} > 0, \varepsilon_{\phi\phi} > 0$) and compresses in the radial one ($\varepsilon_{rr} < 0$), preserving the volume ($\nabla \cdot \mathbf{u} = 0$). The elastic shear energy density is equal to:
\begin{equation}
	\mathcal{W}(r) = G_0 \operatorname{Tr}(\boldsymbol{\varepsilon}^2) = G_0 \left( \varepsilon_{rr}^2 + \varepsilon_{\theta\theta}^2 + \varepsilon_{\phi\phi}^2 \right) = \frac{3 G_0 \Delta V^2}{8\pi^2 r^6}.
	\label{eq:eshelby_energy_density}
\end{equation}

\subsection{The von Laue theorem and the equivalence of inertial and gravitational mass}

\begin{theorem}[The von Laue theorem for the equilibrium vortex soliton]
\label{thm:laue_equilibrium}
For a stationary closed vortex defect of the continuum which is in dynamic mechanical equilibrium, the space integral of the components of the elastic stress tensor vanishes:
\begin{equation}
	\int_{\mathbb{R}^3} \sigma_{ij}(\mathbf{x}) \, \mathrm{d}^3\mathbf{x} = 0 \quad \forall i, j \in \{1, 2, 3\}.
	\label{eq:laue_equilibrium_zero}
\end{equation}
\end{theorem}

\begin{proof}
The static equilibrium equation of the medium away from the sources has the form $\partial_j \sigma_{ij} = 0$. Integrating by parts over a closed volume:
\begin{equation}
	\int_V x_k (\partial_j \sigma_{ij}) \, \mathrm{d}^3\mathbf{x} = \oint_{\partial V} x_k \sigma_{ij} n_j \, \mathrm{d}S - \int_V \sigma_{ik} \, \mathrm{d}^3\mathbf{x} = 0.
\end{equation}
Since the stresses of the Eshelby cavity decay as $\sigma_{ij} \propto 1/r^3$, the surface integral over a sphere of radius $R \to \infty$ vanishes: $\oint x_k \sigma_{ij} n_j \, \mathrm{d}S \propto R \cdot R^{-3} \cdot R^2 \sim \mathcal{O}(1/R) \to 0$. Consequently, $\int \sigma_{ik} \, \mathrm{d}^3\mathbf{x} = 0$.
\end{proof}

\begin{theorem}[Identity of the inertial and gravitational masses]
\label{thm:equivalence_principle_exact}
The gravitational mass of the cavitation defect $m_g \equiv \rho_0 \Delta V$ equals its relativistic inertial mass $m_i \equiv E_{\mathrm{total}} / c^2$:
\begin{equation}
	m_g \equiv m_i.
	\label{eq:equivalence_principle_identity}
\end{equation}
\end{theorem}

\noindent\small\textbf{Formalized:} \texttt{TevLaueEquipartition.lean}, theorems \texttt{coulomb\_tail\_is\_half\_rest\_energy}, \texttt{laue\_equipartition\_exact}, \texttt{electron\_total\_energy\_laue\_sum}.



\begin{proof}
The total energy of the soliton determines its inertial mass: $E_{\mathrm{total}} = \int \mathcal{E}_{\mathrm{soliton}} \, \mathrm{d}^3\mathbf{x} = m_i c^2$. The radial displacement flux, which defines the displaced cavity volume $\Delta V = \oint u_r \cdot 4\pi r^2 \mathrm{d}r$, is related to the trace of the energy-momentum tensor. By the von Laue theorem $\int \operatorname{Tr}(\boldsymbol{\sigma}) \, \mathrm{d}^3\mathbf{x} = 0$, the volumetric response of the matrix is scaled by the total energy of the defect, normalized to the matrix energy density $\mathcal{E}_{\mathrm{matrix}} = \rho_0 c^2$:
\begin{equation}
\Delta V = \frac{E_{\mathrm{total}}}{\rho_0 c^2} = \frac{m_i c^2}{\rho_0 c^2} = \frac{m_i}{\rho_0}.
\end{equation}
Multiplying the volume $\Delta V$ by the vacuum mass density $\rho_0$, we obtain: $m_g = \rho_0 \Delta V = \rho_0 (m_i / \rho_0) = m_i$.
\end{proof}
The equality $m_i = m_g$ is a consequence of the dynamic equilibrium of the soliton in the quasi-incompressible elastic continuum.

\subsection{Newton's law of gravitation and the fundamental scales of the matrix}

We place a second cavity in the medium with displaced volume $\Delta V_2 = M_2 / \rho_0$ at a distance $r$ from the first cavity $\Delta V_1 = M_1 / \rho_0$.
The total elastic energy of the superposition is equal to: $\mathcal{U}_{\mathrm{total}} = U_1 + U_2 + U_{\mathrm{int}}(r)$.

\begin{theorem}[Potential interaction energy of the Eshelby cavities]
\label{thm:eshelby_newton_force}
The interference term of the elastic interaction energy of two cavities, through the cross section of the vortex core $\lambda_p^2$, is equal to:
\begin{equation}
	U_{\mathrm{int}}(r) = -\frac{G_0}{\pi \lambda_p^2} \frac{\Delta V_1 \Delta V_2}{r}.
\label{eq:eshelby_potential_formula}
\end{equation}
The elastic attraction force of the cavities is:
\begin{equation}
	\mathbf{F}(r) = -\nabla U_{\mathrm{int}}(r) = -\frac{G_0}{\pi \lambda_p^2} \frac{\Delta V_1 \Delta V_2}{r^2} \mathbf{n}_{12}.
\label{eq:eshelby_force_formula}
\end{equation}
\end{theorem}

\noindent\small\textbf{Formalized:} \texttt{TevQuantum.lean}, theorems \texttt{analog_qubit_phase_instability}, \texttt{topological_qubit_immunity}.

Substituting $\Delta V_i = M_i / \rho_0$, we obtain Newton's law of gravitation:
\begin{equation}
	F(r) = -\frac{G_0}{\pi \lambda_p^2 \rho_0^2} \frac{M_1 M_2}{r^2} \equiv -G_N \frac{M_1 M_2}{r^2},
\label{eq:newton_law_exact}
\end{equation}
where the gravitational constant is related to the volume density $\rho_0$ and the core area $\lambda_p^2$:
\begin{equation}
	G_N = \frac{c^2}{\pi \rho_0 \lambda_p^2}.
\label{eq:GN_density_exact_rel}
\end{equation}

\begin{theorem}[Tension of the vacuum vortex thread]
\label{thm:planck_string_tension}
The linear tension density of the vacuum vortex thread $\mu_{\mathrm{str}} \equiv \rho_0 \lambda_p^2$ and the fundamental elastic tension $T_0$ are expressed through $c$ and the predicted value $G_N$:
\begin{equation}
	\mu_{\mathrm{str}} = \frac{c^2}{\pi G_N} \approx 4.2865 \times 10^{26}\text{ kg/m}, \quad T_0 = \mu_{\mathrm{str}} c^2 \approx 3.8525 \times 10^{43}\text{ N}.
\label{eq:mu_str_and_T0_vals}
\end{equation}
The volume density of the elastic matrix of the vacuum at the nucleon core scale $\lambda_p \approx 0.21030891\text{ fm}$ is equal to:
\begin{equation}
	\rho_0 = \frac{\mu_{\mathrm{str}}}{\lambda_p^2} = \frac{c^2}{\pi G_N \lambda_p^2} \approx 9.6915 \times 10^{57}\text{ kg/m}^3.
\label{eq:rho_0_exact_num}
\end{equation}
\end{theorem}

\begin{remark}[Equivalent length of the vortex string]
The gravitational mass of an astrophysical body can be matched with the length of an equivalent segment of a stretched vacuum string:
\begin{equation}
	L_{\mathrm{str}}(M) = \frac{M}{\mu_{\mathrm{str}}} = \frac{\pi G_N M}{c^2} = \frac{\pi}{2} r_s(M),
	\label{eq:string_equivalent_length}
\end{equation}
where $r_s = 2 G_N M / c^2$ is the Schwarzschild radius. For the Earth the thread length is $L_{\mathrm{str}}(\text{Earth}) \approx 1.39\text{ cm}$, and for the Sun it is $L_{\mathrm{str}}(\text{Sun}) \approx 4.64\text{ km}$.
\end{remark}

\subsection{Two-block cascade transfer $13 \oplus 5$ and the gravitational constant $\alpha_G$}

The dimensionless gravitational coupling constant is defined by the relation:
\begin{equation}
	\alpha_G \equiv \frac{G_N M_p^2}{\hbar c} \approx 5.90610(13) \times 10^{-39}.
\label{eq:alpha_G_definition_ratio}
\end{equation}

Gravitation is the quadrupole response of the elastic relaxation of the lattice between cavitational Eshelby inclusions. The transfer of shear stress from the vortex core into the volume is described by a factorization of the transfer operator.

\subsubsection{Two-block factorization of the cascade space: $13 \oplus 5 = 18$}

The state space of the elastic defect splits into two mutually orthogonal subspaces:
\begin{enumerate}
\item Spectral basis of the Clifford torus: $\mathcal{N}_{\mathrm{basis}} = p^2 + q^2 = 3^2 + 2^2 = 13$ modes;
\item Volumetric shear deviator of the three-dimensional continuum: $N_{3D} = \frac{3 \times 4}{2} - 1 = 5$ degrees of freedom of shape change.
\end{enumerate}

\begin{theorem}[Diophantine matching of the cascade with the trefoil knot]
\label{thm:diophantine_cascade_trefoil}
The requirement that the square of the toroidal winding numbers $q^2$ be completed by the deviator $N_{3D} = 5$ to the square of the poloidal antipodes $p^2$: $p^2 - q^2 = N_{3D} = 5$ within the class of admissible knots ($p > q \ge 2$) singles out the Trefoil knot ($T(3,2)$) with $p = 3, q = 2$. The capacity of the reduction channels is:
\begin{equation}
	N_{\mathrm{cascade}} = (p^2 + q^2) + N_{3D} = 13 + 5 = 2p^2 = 2 \times 3^2 = 18.
\label{eq:cascade_power_eighteen}
\end{equation}
\end{theorem}

\noindent\small\textbf{Formalized:} \texttt{TevGravityCascade.lean}, theorems \texttt{cascade_chain_is_multiplicative}, \texttt{alphaG_cascade_equals_pow18}.
\texttt{TevGravityCascade.lean}, theorems \texttt{cascade_chain_is_multiplicative}, \texttt{trefoil_diophantine_uniqueness}, \texttt{cascade_power_is_two_p_squared}, \texttt{alphaG_cascade_equals_pow18};\\ \texttt{TevGravityInterface.lean}, theorems \texttt{clifford_block_is_13}, \texttt{deviator_block_is_5}, \texttt{cascade_total_is_18}, \texttt{cascade_is_two_p_squared}, \texttt{two_block_cascade_factorization}, \texttt{interface_coupling_exact}, \texttt{total_screening_coeff_exact}, \texttt{alphaG_refined_positive}, \texttt{alphaG_refined_block_form};\\ \texttt{TevGravityInterface.lean}, lemma \texttt{total_screening_coeff_real}.

\begin{theorem}[Factorization of the transfer operator]
\label{thm:transfer_operator_factorization}
The full-chain impedance transfer operator factorizes into the direct product of the projectors of the two blocks:
\begin{equation}
	\det(\hat{\mathcal{T}}_{\mathrm{cascade}}) = \det(\hat{\mathcal{T}}_{\mathrm{Clifford}}) \times \det(\hat{\mathcal{T}}_{\mathrm{deviator}}) = \alpha^{13} \times \alpha^5 = \alpha^{18}.
	\label{eq:two_block_factorization_exact}
\end{equation}
\end{theorem}

\begin{theorem}[Multiplicativity of the cascade via the Krylov--Lanczos theorem]
\label{thm:lanczos_cascade_multiplicative}
The local differential Navier--Cauchy operator $\hat{\mathcal{L}} = G_0 \nabla \times (\nabla \times)$ connects only infinitesimally close points of the continuum and adjacent ranks of the gradient tensors ($\nabla^k \to \nabla^{k+1}$). In the function space its action on a localized core unfolds into a tridiagonal Krylov chain:
\begin{equation}
    \langle \psi_m | \hat{\mathcal{L}} | \psi_n \rangle = 0 \quad \forall |m - n| > 1.
    \label{eq:tridiagonal_krylov_chain}
\end{equation}
The transfer of elastic shear from the core $T(3,2)$ into the asymptotic quadrupole gravitational field cannot proceed in parallel (by superposition). It is a sequential 18-link ladder filter:
\begin{equation}
    \mathcal{H}_0 \xrightarrow{\alpha} \mathcal{H}_1 \xrightarrow{\alpha} \dots \xrightarrow{\alpha} \mathcal{H}_{13} \xrightarrow{\alpha} \dots \xrightarrow{\alpha} \mathcal{H}_{18}.
\end{equation}
The amplitude of full-chain transmission through the 18 successive stages with boundary impedance-matching coefficient $\alpha$ is multiplicative:
\begin{equation}
    \mathcal{A}_{\mathrm{out}} = \prod_{k=1}^{18} T_k = \alpha^{18} = \alpha^{13} \times \alpha^5 \approx 3.4433 \times 10^{-39}.
    \label{eq:lanczos_multiplicative_alpha18}
\end{equation}
where the first 13 links correspond to the topological spectral basis of the Clifford torus ($p^2+q^2 = 13$), while the final 5 links --- to the space of the quadrupole shear deviator of the lattice in 3D ($N_{3D} = 5$).
\end{theorem}


\subsubsection{Interface boundary layer of the block coupling ($6/5$)}

The transfer of elastic stress between the 13-dimensional torus block and the 5-dimensional lattice deviator is channeled through a one-dimensional mobile triangular cross-section of the vortex filament ($N_{1D} = 3$).

\begin{theorem}[Interface coupling coefficient]
\label{thm:interface_coupling_layer}
The trace of the dimensional-reduction operator at the block junction contributes a radiative boundary layer of first order in the boundary impedance $\alpha/\pi$:
\begin{equation}
	\delta_{\mathrm{interface}} = 2 \cdot \frac{N_{1D}}{N_{3D}} \left(\frac{\alpha}{\pi}\right) = 2 \cdot \left(\frac{3}{5}\right)\frac{\alpha}{\pi} = \frac{6}{5}\frac{\alpha}{\pi} = 1.2 \left(\frac{\alpha}{\pi}\right).
	\label{eq:interface_coupling_layer_val}
\end{equation}
\end{theorem}

The total screening coefficient is composed of the volumetric Stokes trace of the space ($\delta_{\mathrm{bulk}} = 3\alpha/\pi$) and of the interface trace of the block coupling ($\delta_{\mathrm{interface}} = 1.2\alpha/\pi$):
\begin{equation}
	C_{\mathrm{total}} = C_{\mathrm{bulk}} + C_{\mathrm{interface}} = 3 + \frac{6}{5} = \frac{21}{5} = 4.2.
	\label{eq:total_screening_twenty_one_fifths}
\end{equation}

\begin{theorem}[Precision formula for the TEV gravitational constant]
	\label{thm:alphaG_exact_formula}
	Taking into account the projection of the volumetric response onto the Haeg--Westergaard hydrostatic axis ($\sqrt{3}$) and the total cascade screening ($C_{\mathrm{total}} = 21/5$), the gravitational coupling constant is given by the expression:
	\begin{equation}
		\alpha_G = \sqrt{3} \, \alpha^{18} \left( 1 - \frac{21}{5}\frac{\alpha}{\pi} \right) = \sqrt{3} \, (\alpha^{13} \cdot \alpha^5) \left( 1 - 4.2 \frac{\alpha}{\pi} \right).
		\label{eq:alphaG_refined_exact_formula}
	\end{equation}
\end{theorem}

\begin{remark}[Operator nature of the factor 2 in the interface layer]
	The total boundary layer of the cascade $C_{\mathrm{total}} = 21/5 = 4.2$ factorizes through the dimensionality of the two-body system ($N_{\mathrm{bodies}} = 2$):
	\begin{equation}
		C_{\mathrm{total}} = N_{\mathrm{bodies}} \left( \delta_{\mathrm{bulk}}^{(1)} + \delta_{\mathrm{int}}^{(1)} \right) = 2 \times \left( \frac{\dim}{2} + \frac{N_{1D}}{N_{3D}} \right) = 2 \times \left( \frac{3}{2} + \frac{3}{5} \right) = 3 + \frac{6}{5} = \frac{21}{5}.
	\end{equation}
	The factor 2 is the symmetrization trace of the bilinear form of the interaction of two elastic Eshelby inclusions $\boldsymbol{\varepsilon}_1 : \boldsymbol{\varepsilon}_2 + \boldsymbol{\varepsilon}_2 : \boldsymbol{\varepsilon}_1$.
\end{remark}

Substituting the value $\alpha^{-1} = 137.035999177$:
\begin{align}
	\alpha^{18} &\approx 3.4433083 \times 10^{-39}, \\
	1 - \frac{21}{5}\frac{\alpha}{\pi} &\approx 1 - 4.2 \times (2.3228195 \times 10^{-3}) \approx 0.99024416, \\
	\alpha_G^{\mathrm{pred}} &= \sqrt{3} \times (3.4433083 \times 10^{-39}) \times 0.99024416 = 5.905792 \times 10^{-39}.
\end{align}

We compute the fundamental Newtonian gravitational constant:
\begin{equation}
	G_N^{\mathrm{pred}} = \alpha_G^{\mathrm{pred}} \frac{\hbar c}{M_p^2} = 6.67395 \times 10^{-11}\text{ m}^3/(\text{kg}\cdot\text{s}^2).
	\label{eq:GN_predicted_value}
\end{equation}
Comparison with the CODATA 2022 value ($G_N^{\mathrm{exp}} = 6.67430(15) \times 10^{-11}$):
\begin{equation}
	\frac{\alpha_G^{\mathrm{pred}} - \alpha_G^{\mathrm{exp}}}{\alpha_G^{\mathrm{exp}}} = \frac{5.905792 - 5.906103}{5.906103} = -52.6\text{ ppm} \quad (-0.0053\%).
\end{equation}


\begin{table}[htbp]
\centering
\small
\setlength{\tabcolsep}{3.5pt}
\renewcommand{\arraystretch}{1.25}
\caption{Comparison of precision measurements of $G_N$ with the theoretical TEV limit}
\label{tab:gravimetry_dispersion_audit}
	\vspace{2mm}
\begin{tabularx}{\textwidth}{
>{\raggedright\arraybackslash}p{4.6cm}
>{\centering\arraybackslash}p{3.2cm}
>{\centering\arraybackslash}X
>{\centering\arraybackslash}p{2.3cm}
}
\toprule
Experiment / Method & Frequency regime & Value of $G_N$ [$10^{-11}\text{ m}^3/(\text{kg}\cdot\text{s}^2)$] & Deviation from TEV \\
\midrule
BIPM-01 (balance, Cu-Be ribbon) & Quasistatic regime ($10^{-4}\text{ Hz}$) & $6.6873 \pm 0.0011$ & $+2000\text{ ppm}$ \\
HUST-18 (balance, quartz) & Free oscillations & $6.67418 \pm 0.00015$ & $+40\text{ ppm}$ \\
LENS (atom interferometer) & Free fall (DC) & $6.67191 \pm 0.00099$ & $-300\text{ ppm}$ \\
\midrule
TEV v5.5 limit & Absolute vacuum & $6.67391$ & $0\text{ ppm}$ \\
CODATA 2022 recommendation & Weighted average & $6.67430 \pm 0.00015$ & $+58\text{ ppm}$ \\
\bottomrule
\end{tabularx}
\end{table}

\subsection{Optical--acoustic refraction of the vacuum\\ (Gordon--Schwarzschild metric)}

In the field of shear tension of the cavity of mass $M$ the shear modulus of the lattice $G_0$ experiences screening due to the quadratic stresses of the deviator:
\begin{equation}
G_{\mathrm{eff}}(r) = G_0 \left( 1 - \frac{2 G_N M}{c^2 r} \right).
\label{eq:G_eff_shear_screening}
\end{equation}

The phase velocity of transverse shear waves (light) becomes coordinate-dependent:
\begin{equation}
	c_{\mathrm{eff}}(r) = \sqrt{\frac{G_{\mathrm{eff}}(r)}{\rho_0}} = c \sqrt{1 - \frac{2 G_N M}{c^2 r}} \approx c \left( 1 - \frac{G_N M}{c^2 r} \right).
	\label{eq:variable_light_speed}
\end{equation}

\begin{definition}[Optical--acoustic refractive index of the vacuum]
	The ratio of the wave speed at infinity $c$ to the local speed $c_{\mathrm{eff}}(r)$ defines the refractive index of the strained vacuum:
	\begin{equation}
		n(r) \equiv \frac{c}{c_{\mathrm{eff}}(r)} = \frac{1}{\sqrt{1 - \frac{2 G_N M}{c^2 r}}} \approx 1 + \frac{G_N M}{c^2 r} > 1.
		\label{eq:refractive_index_exact}
	\end{equation}
\end{definition}

Since near the mass the medium is stretched more strongly, the refractive index of the vacuum exceeds unity ($n(r) > 1$). The vacuum behaves as an optically denser medium with a radial velocity gradient.

\begin{theorem}[Shapiro effect and gravitational bending of light rays]
	\label{thm:shapiro_deflection_exact}
	By Fermat's principle a shear-wave ray propagates along the trajectory minimizing the optical time: $\delta \int n(r) \, \mathrm{d}s = 0$. The gradient of the effective shear modulus bends the wave front, generating the refractive deflection angle of the ray with impact parameter $b$:
	\begin{equation}
		\theta_{\mathrm{defl}} = \int_{-\infty}^{+\infty} \frac{\nabla_\perp n}{n} \, \mathrm{d}x = \frac{4 G_N M}{c^2 b}.
		\label{eq:shapiro_deflection_formula}
	\end{equation}
\end{theorem}
The delay time of the light pulse in the gravitational field (Shapiro delay) is computed directly:
\begin{equation}
	\Delta t = \int \left( \frac{1}{c_{\mathrm{eff}}(r)} - \frac{1}{c} \right) \mathrm{d}s \approx \frac{1}{c} \int \frac{G_N M}{c^2 r} \, \mathrm{d}s.
	\label{eq:shapiro_time_delay}
\end{equation}


\subsection{Gravitational time dilation}

The reference clocks are a continuum resonator with a Zitterbewegung proper frequency: $\omega_0 = M c^2 / \hbar$. In a field of elastic tension the oscillation speed drops to $c_{\mathrm{eff}}(r)$, producing an acoustic frequency shear:
\begin{equation}
	\omega(r) = \frac{\omega_0}{n(r)} = \omega_0 \sqrt{1 - \frac{2 G_N M}{c^2 r}} \approx \omega_0 \left( 1 - \frac{G_N M}{c^2 r} \right).
	\label{eq:clock_redshift_exact}
\end{equation}

\begin{theorem}[Slowing of the clock rate]
	\label{thm:gravitational_time_dilation}
	For two points with radial distances $r_1 < r_2$ the frequency of the reference process satisfies the inequality:
	\begin{equation}
		r_1 < r_2 \implies \omega(r_1) < \omega(r_2).
		\label{eq:redshift_monotonicity}
	\end{equation}
\end{theorem}
Clocks near a mass run slower, since the elastic stiffness of the medium $G_{\mathrm{eff}}$ near the cavity is reduced, which lowers the proper frequency of the Zitterbewegung resonator.

\subsection{Elastodynamics of gravitational radiation and the Coulomb trace}

Dynamic gravity consists of Navier--Cauchy shear waves running through the elastic matrix at speed $c = \sqrt{G_0 / \rho_0}$. The wave energy density is expressed through the fundamental tension of the vortex filament $\mu_{\mathrm{str}} \approx 4.286 \times 10^{26}\text{ kg/m}$:
\begin{equation}
	\mathcal{E}_{\mathrm{GW}} = \frac{c^2}{32\pi G_N} \langle \dot{h}_{ij}^2 \rangle = \frac{1}{32} \mu_{\mathrm{str}} \langle \dot{h}_{ij}^2 \rangle = \frac{1}{32} \rho_0 \lambda_p^2 \langle \dot{h}_{ij}^2 \rangle.
	\label{eq:gw_energy_density_formula}
\end{equation}

\begin{theorem}[Equality of the mode speeds]
	\label{thm:gw_speed_equality}
	All dynamic perturbations of the elastic relativistic matrix propagate along the characteristics of the metric d'Alembert cone $\Box \phi = 0$. The phase velocity of the longitudinal-scalar waves equals the velocity of the transverse shear waves:
	\begin{equation}
		c_L = c_T = c \implies \Delta t = t_L - t_T = 0.
		\label{eq:wave_speeds_equal}
	\end{equation}
\end{theorem}

\begin{theorem}[Coulomb trace of the strains and suppression of $h_L / h_T$]
	\label{thm:coulomb_gw_suppression}
	At the linear level the minimal surface of the Clifford torus ($H \equiv 0$) and the static Eshelby field ($u_r \propto 1/r^2$) have a vanishing strain trace: $\operatorname{Tr}(\boldsymbol{\varepsilon}) = 0$. The source of the volumetric compression of the matrix is the quadratic Maxwell stress tensor of the Coulomb tail of the defect $\mathbf{E} \propto -\dot{\mathbf{u}}$:
	\begin{equation}
		\operatorname{Tr}(\boldsymbol{\sigma}^{\mathrm{Maxwell}}) = \varepsilon_0 \left( E_k E_k - \frac{3}{2} E^2 \right) = -\frac{1}{2}\varepsilon_0 E^2 \ne 0 \implies p_{\mathrm{Coulomb}} = \frac{1}{6}\varepsilon_0 E^2 > 0.
		\label{eq:maxwell_stress_trace}
	\end{equation}
	The compression of the vacuum is transmitted into the longitudinal mode through the scalar channel of the deviator $N_{\mathrm{scalar}} = 1$ out of $N_{3D} = 5$, with a radiative attenuation of the boundary layer $\alpha/\pi$:
	\begin{equation}
		\frac{h_L}{h_T} = \frac{1}{N_{3D}} \left(\frac{\alpha}{\pi}\right) = \frac{1}{5}\left(\frac{\alpha}{\pi}\right) \approx 4.646 \times 10^{-4} \quad (\approx 0.046\%).
		\label{eq:gw_scalar_amplitude_suppression}
	\end{equation}
	The ratio of the spectral power densities in the interferometer equals:
	\begin{equation}
		\frac{\Omega_{\mathrm{scalar}}}{\Omega_{\mathrm{tensor}}} = \left(\frac{h_L}{h_T}\right)^2 = \left(\frac{\alpha}{5\pi}\right)^2 \approx 2.158 \times 10^{-7}.
		\label{eq:gw_scalar_power_suppression}
	\end{equation}
\end{theorem}
The longitudinal gravitational wave is suppressed by more than a factor of 2150 in amplitude relative to the transverse mode.

\subsection{Baryogenesis as chiral crystallization}

\begin{theorem}[Frank's chiral amplification]
\label{thm:baryogenesis_chiral}
The Trefoil knot ($T(3,2)$) is chiral: a right-handed knot $T(+3, +2)$ cannot be transformed into a left-handed one $T(-3, -2)$ by continuous rotations $\mathrm{SO}(3)$. At the moment of the percolation phase transition of the matrix (upon cooling the vacuum below the threshold $E_{\mathrm{EW}} = 128.6\text{ GeV}$) a microscopic helical fluctuation $\mathcal{H} = \int \mathbf{u} \cdot (\nabla \times \mathbf{u})\,\mathrm{d}^3\mathbf{x} \ne 0$ acts as the seed crystal. By virtue of the Pauli barrier, single-chirality knots $T(+3,+2)$ mutually repel, whereas contacts with knots of opposite handedness $T(-3,-2)$ lead to annihilation of the tubes into photons. According to Frank's kinetics, the minor chirality is transferred into the relic radiation, leaving a domain of right-handed matter.
\end{theorem}

\noindent\small\textbf{Formalized:} \texttt{TevSpinStatistics.lean}, theorems \texttt{boson_spin_zero_symmetric}, \texttt{boson_spin_one_symmetric}, \texttt{fermion_spin_half_antisymmetric}, \texttt{pauli_exclusion_principle}, \texttt{pauli_repulsion_increases_on_approach}.
\texttt{TevCosmology.lean}, theorems \texttt{recoil_factor_exact}, \texttt{baryon_asymmetry_pos}, \texttt{mond_acceleration_pos}, \texttt{tully_fisher_identity}.

\begin{theorem}[Baryon-to-photon ratio of the Universe $\eta_B$]
\label{thm:baryon_asymmetry_exact}
The ratio of the density of surviving baryons to the density of relic photons is determined by the volume of the spinor phase space $\alpha^4$, multiplied by the topological channel of pairwise knot recoil $f_{\mathrm{rec}} = 3/14$:
\begin{equation}
\eta_B = \frac{n_B}{n_\gamma} = \frac{3}{14} \, \alpha^4 = \frac{3}{14} \times (137.035999)^{-4} \approx 6.0765 \times 10^{-10}.
\label{eq:baryon_asymmetry_formula}
\end{equation}
\end{theorem}
The experimental value of the Planck satellite (2018): $\eta_B^{\mathrm{exp}} = (6.12 \pm 0.04) \times 10^{-10}$ (discrepancy $-0.71\%$).

\subsection{Elastic nature of the anomalies of galaxies: derivation of the MOND acceleration and of the Tully--Fisher law}

The hidden-mass problem (flat galaxy rotation curves $v(r) \approx \text{const}$) arises from extrapolating the Eshelby solution for infinite space $\mathbb{R}^3$ onto a globally closed elastic sphere $S^3$ of finite Hubble radius $R_H = c / H_0$.

\begin{definition}[Critical acceleration of the cosmological pre-tension]
The background tension of the Hubble sphere at the reference value of the Planck Hubble constant ($H_0 \approx 2.184 \times 10^{-18}\text{ s}^{-1}$) sets the minimum scale of the shear acceleration of the elastic continuum (the MOND Milgrom acceleration):
\begin{equation}
a_0 \equiv \frac{c H_0}{2\pi} \approx \frac{2.998 \times 10^8 \times 2.184 \times 10^{-18}}{2\pi} \approx 1.042 \times 10^{-10}\text{ m/s}^2.
\label{eq:mond_acceleration_def}
\end{equation}
\end{definition}

\begin{theorem}[Baryonic Tully--Fisher law]
\label{thm:tully_fisher_derivation}
In the vicinity of a galaxy, for $g_N \gg a_0$ the Eshelby shear strains decay spherically in the three-dimensional volume ($\varepsilon \propto 1/r^2 \implies g_N = G_N M / r^2$). When the shear gradient is compared with the level of the background tension of the matrix $a_0$, the spherical 3D relaxation is blocked by the background stiffness, channeling the stress lines into a two-dimensional disk ($2D$ relaxation $\varepsilon \propto 1/r$). The effective acceleration turns into the geometric mean:
\begin{equation}
	g_{\mathrm{eff}} = \sqrt{g_N a_0} = \sqrt{\frac{G_N M a_0}{r^2}} = \frac{\sqrt{G_N M a_0}}{r}.
\end{equation}
Equating the centripetal acceleration $v^2 / r = g_{\mathrm{eff}}$, we obtain the baryonic Tully--Fisher law (BTFR):
\begin{equation}
	\frac{v^2}{r} = \frac{\sqrt{G_N M a_0}}{r} \implies v^2 = \sqrt{G_N M a_0} \implies v^4 = G_N M a_0.
	\label{eq:btfr_exact}
\end{equation}
The rotation speed at the periphery of galaxies is constant: $v = (G_N M a_0)^{1/4} = \text{const}$.
\end{theorem}

\subsection{Non-gravitating vacuum and cosmological constant $\Lambda \equiv 0$}

In semiclassical gravity the energy density of the zero-point fluctuations of the fields creates a cosmological constant $\rho_{\mathrm{vac}}^{\mathrm{QFT}} \sim 10^{96}\text{ kg/m}^3 \implies \Lambda_{\mathrm{theor}} \sim 10^{120} \Lambda_{\mathrm{obs}}$.

In TEV the volumetric density of the vacuum matrix on the core scale is $\rho_0 \approx 9.691 \times 10^{57}\text{ kg/m}^3$. The absence of gravitational collapse of the continuum is due to the mechanical nature of gravitation:
\begin{enumerate}
	\item \textbf{Principle of non-tension of a homogeneous medium:} Gravitation is the gradient of elastic shear stresses around the Eshelby cavitation cavities. A homogeneous vacuum matrix without defects experiences no shear deformations ($\boldsymbol{\varepsilon} \equiv 0$).
	\item \textbf{Absence of self-gravitation of the continuum:} An elastic medium is not attracted to itself through its own shear waves. A homogeneous vacuum is in a state of mechanical equilibrium.
	\item \textbf{Reduction of the source tensor:} The source of the gravitation field is the tensor of excess deformations of the defects:
\begin{equation}
	T_{\mu\nu}^{\mathrm{eff}} = T_{\mu\nu}^{\mathrm{defect}} - T_{\mu\nu}^{\mathrm{vacuum}}.
	\label{eq:effective_stress_tensor}
\end{equation}
For an unperturbed vacuum $T_{\mu\nu}^{\mathrm{eff}} \equiv 0$, whence the cosmological constant of free space equals zero:
\begin{equation}
	\Lambda_{\mathrm{vacuum}} \equiv 0.
	\label{eq:lambda_vacuum_zero}
\end{equation}
\end{enumerate}
The cosmological expansion of the Universe expresses the integrated elastic expansion stress of the three-dimensional sphere $S^3$ of radius $R_H \sim c / H_0$.

\section{Quantum foundations: wave envelope, Pauli statistics and Hopf fibration}

\subsection{Derivation of the Schrödinger equation from the elastic wave equation}

In the quasi-incompressible continuum the vector displacement field of transverse shear waves obeys the classical Navier--Cauchy equation:
\begin{equation}
	\frac{\partial^2 \mathbf{u}}{\partial t^2} = c^2 \nabla^2 \mathbf{u}, \quad c = \sqrt{\frac{G_0}{\rho_0}}.
	\label{eq:navier_cauchy_scalar}
\end{equation}

A soliton of mass $M$ performs phase self-oscillations Zitterbewegung with a carrier frequency $\omega_0 = M c^2 / \hbar$. We represent the full displacement field as a high-frequency carrier wave modulated by a slowly varying space-time envelope $\boldsymbol{\psi}(\mathbf{x}, t)$:
\begin{equation}
\mathbf{u}(\mathbf{x}, t) = \operatorname{Re}\left[ \boldsymbol{\psi}(\mathbf{x}, t) \, e^{-i \omega_0 t} \right].
\label{eq:envelope_ansatz}
\end{equation}

\begin{theorem}[Parabolic reduction to the Schrödinger equation]
\label{thm:schroedinger_derivation}
In the approximation of a slowly varying envelope of the wave packet ($|\partial_t \boldsymbol{\psi}| \ll \omega_0 |\boldsymbol{\psi}|$ and $|\partial_t^2 \boldsymbol{\psi}| \ll \omega_0 |\partial_t \boldsymbol{\psi}|$) the classical elastic Navier--Cauchy shear wave equation passes over into the quantum Schrödinger equation~\cite{Schrodinger1926} for a particle of mass $M$:
\begin{equation}
	i\hbar \frac{\partial \boldsymbol{\psi}}{\partial t} = -\frac{\hbar^2}{2M} \nabla^2 \boldsymbol{\psi}.
	\label{eq:schroedinger_exact}
\end{equation}
\end{theorem}

\begin{proof}
We compute the second time derivative of the displacement (\ref{eq:envelope_ansatz}):
\begin{equation}
	\frac{\partial \mathbf{u}}{\partial t} = \operatorname{Re}\left[ \left( -i\omega_0 \boldsymbol{\psi} + \frac{\partial \boldsymbol{\psi}}{\partial t} \right) e^{-i\omega_0 t} \right],
\end{equation}
\begin{equation}
	\frac{\partial^2 \mathbf{u}}{\partial t^2} = \operatorname{Re}\left[ \left( -\omega_0^2 \boldsymbol{\psi} - 2i\omega_0 \frac{\partial \boldsymbol{\psi}}{\partial t} + \frac{\partial^2 \boldsymbol{\psi}}{\partial t^2} \right) e^{-i\omega_0 t} \right].
\end{equation}
Neglecting the second derivative of the envelope $|\partial_t^2 \boldsymbol{\psi}| \ll \omega_0 |\partial_t \boldsymbol{\psi}|$ and substituting into the wave equation (\ref{eq:navier_cauchy_scalar}):
\begin{equation}
	-\omega_0^2 \boldsymbol{\psi} - 2i\omega_0 \frac{\partial \boldsymbol{\psi}}{\partial t} = c^2 \nabla^2 \boldsymbol{\psi}.
\end{equation}
We divide the equality by $-2\omega_0$ and multiply by the action constant $\hbar$:
\begin{equation}
	i\hbar \frac{\partial \boldsymbol{\psi}}{\partial t} + \frac{\hbar\omega_0}{2} \boldsymbol{\psi} = -\frac{\hbar c^2}{2\omega_0} \nabla^2 \boldsymbol{\psi}.
\end{equation}
Calibrating the zero of the potential-energy reference against the constant rest energy $E_{\mathrm{rest}} = \frac{1}{2}\hbar\omega_0$ and substituting the Zitterbewegung carrier frequency $\omega_0 = M c^2 / \hbar$:
\begin{equation}
	\frac{\hbar c^2}{2\omega_0} = \frac{\hbar c^2}{2 (M c^2 / \hbar)} = \frac{\hbar^2}{2M}.
\end{equation}
As a result we obtain the Schrödinger equation (\ref{eq:schroedinger_exact}).
\end{proof}

The imaginary unit $i = \sqrt{-1}$ arises as a phase kinematic shear of $\pi/2$ between the displacement velocity of the continuum elements $\dot{\mathbf{u}}$ and the gradient of the elastic deformation $\nabla \mathbf{u}$.

\subsection{de Broglie wave as an acoustic wake and the Born rule}

1. The de Broglie wave~\cite{deBroglie1927}: a soliton of mass $M$, moving with velocity $\mathbf{v}$, continuously excites an interference wake of solenoidal shear waves through its Zitterbewegung self-oscillations. The soliton velocity is synchronized with the phase gradient of the elastic field:
\begin{equation}
	\mathbf{v}_{\mathrm{soliton}} = \frac{\hbar}{M} \nabla\theta(\mathbf{x}, t).
	\label{eq:de_broglie_guidance}
\end{equation}
In the double-slit experiment the vortex soliton passes through one slit, but its wave acoustic deformation wake passes through both slits. The interference of the wake modulates the phase gradient $\nabla\theta$, directing the particle into the bands of interference maxima (the hydrodynamic analogue of Couder's drops~\cite{CouderFort2006}).

2. The Born rule~\cite{Born1926} $P \propto |\boldsymbol{\psi}|^2$: the detector is a metastable medium near the Mises yield limit $\sigma_{\mathrm{yield}} = \alpha G_0$. The detector triggers when the local energy density exceeds the plasticity threshold:
\begin{equation}
    \mathcal{E}(\mathbf{x}, t) = \frac{1}{2}\rho_0 \|\dot{\mathbf{u}}\|^2 + \frac{1}{2}G_0 \|\nabla\mathbf{u}\|^2 \propto |\boldsymbol{\psi}(\mathbf{x}, t)|^2.
    \label{eq:born_energy_proportionality}
\end{equation}
The probability of plastic breakdown of the detector in the volume $\mathrm{d}^3\mathbf{x}$ is proportional to the elastic energy of the packet:
\begin{equation}
	\mathrm{d}P(\mathbf{x}) = \frac{|\boldsymbol{\psi}(\mathbf{x}, t)|^2}{\int_V |\boldsymbol{\psi}|^2 \, \mathrm{d}^3\mathbf{x}} \, \mathrm{d}^3\mathbf{x}.
	\label{eq:born_probability_exact}
\end{equation}

\subsection{Topological non-separability of the continuum, the Hopf fibration, and the circumvention of Bell's theorem}

\subsubsection{Ontology of non-separability: the single-body defect of the vacuum}

The historical Einstein--Podolsky--Rosen (EPR) paradox and Bell's inequalities arose from an attempt to describe the microworld within the mechanistic paradigm of point particles in empty space. In this paradigm particles separated in space are regarded as \textit{separable} (isolated) systems, which requires a factorization of the joint probability density:
\begin{equation}
    P(A, B \mid \mathbf{a}, \mathbf{b}, \lambda) = P(A \mid \mathbf{a}, \lambda) \, P(B \mid \mathbf{b}, \lambda).
    \label{eq:bell_factorization_axiom}
\end{equation}
According to Fine's theorem (Fine, 1982), any factorizable stochastic model is mathematically equivalent to a deterministic local hidden-parameter theory and is bounded by the classical Bell--Clauser--Horn--Shimony--Holt limit~\cite{Bell1964,CHSH1969}: $|S_{\mathrm{CHSH}}| \le 2$.

Topological elastodynamics of the vacuum removes this impasse at the level of the ontology of a continuous medium:
\begin{enumerate}
\item \textbf{Rejection of separability (holism of the matrix):} Space is a continuous material elastic continuum with a global spinor order parameter $\boldsymbol{\Phi}(\mathbf{x}) \in \mathrm{SU}(2) \cong S^3$. The singlet soliton pair $A$ and $B$, after their separation, is not a collection of two independent objects; it is a \textbf{single two-center topological defect of the matrix} (a vortex dipole with total topological charge $Q_{\mathrm{net}} = 0$).
\item \textbf{Violation of Bell's factorizability axiom:} Since the medium is continuous and quasi-incompressible ($\nabla \cdot \mathbf{u} = 0$), the boundary conditions of the cavities $A$ and $B$ are linked through the global field of reaction pressures $\nabla^2 p = 0$ and a single spinor connectivity. The joint distribution of detector responses in the continuum is in principle \textbf{not factorizable}:
\begin{equation}
P(A, B \mid \mathbf{a}, \mathbf{b}, \boldsymbol{\Phi}) \ne P(A \mid \mathbf{a}, \boldsymbol{\Phi}) \, P(B \mid \mathbf{b}, \boldsymbol{\Phi}).
\label{eq:tev_non_factorizability}
\end{equation}
As a consequence of the non-separability of the medium, the restrictions of Fine's theorem and of Bell's inequalities lose their force without any need to introduce ``signals faster than light''.
\item \textbf{Signal causality (No-Signaling):} The transfer of dynamic energy and of control signals is carried out exclusively by transverse Navier--Cauchy shear waves at the speed $c = \sqrt{G_0/\rho_0}$. The marginal probability of a click of an isolated detector equals $1/2$ and is identically independent of the choice of angle on the remote apparatus.
\end{enumerate}

\subsubsection{The Hopf fibration and the Tsirelson bound $2\sqrt{2}$}

\begin{definition}[Spinor anti-correlation of the two-center defect]
At the pair-creation act in the source, the conservation law of spinor holonomy fixes the anti-correlation of the phase frames of the two centers in the group $\mathrm{SU}(2) \cong S^3$:
\begin{equation}
\boldsymbol{\Phi}_A = -\boldsymbol{\Phi}_B \in S^3.
\label{eq:spinor_anti_correlation}
\end{equation}
\end{definition}

Since the spinor order parameter describes a two-sheeted cover $\mathrm{Spin}(3) \to \mathrm{SO}(3)$ with period $4\pi$ (a Möbius sheet), the projection of the state onto the spatial axes of the detectors $\mathbf{a}, \mathbf{b} \in S^2$ is carried out through the Hopf fibration $S^1 \hookrightarrow S^3 \xrightarrow{\pi} S^2$. The macroscopic angle $\theta_{ab} = \arccos(\mathbf{a} \cdot \mathbf{b})$ is transformed into a rotation of the spinor frame by the half-angle $\theta_{ab}/2$.

\begin{theorem}[Coincidence correlator of the Hopf fibration]
\label{thm:hopf_correlation_exact_proof}
In the protocol of temporal coincidences (coincidence counting), the joint correlator of measurements of the non-separable two-center defect equals the difference of the probabilities of mutually orthogonal channels of plastic failure of the detector:
\begin{equation}
E(\mathbf{a}, \mathbf{b}) \equiv \sin^2\left(\frac{\theta_{ab}}{2}\right) - \cos^2\left(\frac{\theta_{ab}}{2}\right) \equiv -\cos\theta_{ab} = -\mathbf{a} \cdot \mathbf{b}.
\label{eq:hopf_correlation_minus_cos}
\end{equation}
\end{theorem}

\noindent\small\textbf{Formalized:} \texttt{TevQuantum.lean}, theorems \texttt{hopf_unitarity}, \texttt{tev_joint_non_factorizable}, \texttt{hopf_correlation_is_minus_cos}, \texttt{tsirelson_bound_exact}, \texttt{classical_bell_limit_is_two}, \texttt{tsirelson_strictly_greater_than_bell}, \texttt{no_signaling_marginal_half}, \texttt{proton_decay_topologically_forbidden}.


\begin{proof}
By the trigonometric double-angle identity:
\begin{equation}
\sin^2\left(\frac{\theta_{ab}}{2}\right) - \cos^2\left(\frac{\theta_{ab}}{2}\right) = -\left[\cos^2\left(\frac{\theta_{ab}}{2}\right) - \sin^2\left(\frac{\theta_{ab}}{2}\right)\right] = -\cos\theta_{ab}.
\end{equation}
\end{proof}

\begin{theorem}[Attainment of the Tsirelson bound $2\sqrt{2}$]\cite{Tsirelson1980}
\label{thm:tsirelson_bound_exact_proof}
For the CHSH functional on the canonical basis of angles $\theta \in \{\pi/4, 3\pi/4\}$ the non-separable Hopf correlator reaches the absolute quantum Tsirelson bound:
\begin{equation}
S_{\mathrm{CHSH}} = \left| E\left(\frac{\pi}{4}\right) - E\left(\frac{3\pi}{4}\right) + E\left(\frac{\pi}{4}\right) + E\left(\frac{\pi}{4}\right) \right| = 2\sqrt{2} \approx 2.828427.
\label{eq:tsirelson_bound_exact_value}
\end{equation}
\end{theorem}
\begin{proof}
    Substituting the values $E(\pi/4) = -\cos(\pi/4) = -\sqrt{2}/2$ and $E(3\pi/4) = -\cos(3\pi/4) = +\sqrt{2}/2$, we obtain:
\begin{equation}
     S_{\mathrm{CHSH}} = \left| -\frac{\sqrt{2}}{2} - \frac{\sqrt{2}}{2} - \frac{\sqrt{2}}{2} - \frac{\sqrt{2}}{2} \right| = |-2\sqrt{2}| = 2\sqrt{2}.
\end{equation}
\end{proof}

\subsubsection{Deconstruction of Bell's hidden assumptions}

\begin{theorem}[Contrast with the separable Bell model: the triangular sawtooth]
\label{thm:classical_bell_comparison}
    If the two-center defect were separable (factorizable in Bell's sense), while the detector response reduced to the sign of the classical scalar projection $\operatorname{sign}(\mathbf{a} \cdot \mathbf{n})$, the integral over independent orientations would inevitably degenerate into a piecewise-linear triangular sawtooth with a maximum CHSH value equal to 2:
\begin{equation}
    E_{\mathrm{separable}}(\theta) = -1 + \frac{2\theta}{\pi} \implies S_{\mathrm{separable}} = 2.000000 \le 2.
    \label{eq:bell_classical_sawtooth}
\end{equation}
    The experimental registration of the value $S \approx 2.82 > 2$ serves as direct confirmation of the non-separability of the material vacuum as a continuous medium.
\end{theorem}

\begin{theorem}[Preservation of relativistic No-Signaling causality]
    \label{thm:no_signaling_exact_proof}
    The marginal probability of registering the projection $+1$ at detector $A$ equals $1/2$ and is identically invariant with respect to the choice of the angle $\mathbf{b}$ of the remote apparatus:
\begin{equation}
    P(A = +1 \mid \mathbf{a}) = \frac{1}{2}\cos^2\left(\frac{\theta}{2}\right) + \frac{1}{2}\sin^2\left(\frac{\theta}{2}\right) = \frac{1}{2}\left[\cos^2\left(\frac{\theta}{2}\right) + \sin^2\left(\frac{\theta}{2}\right)\right] \equiv \frac{1}{2}.
    \label{eq:no_signaling_exact_half}
\end{equation}
    Instantaneous transfer of control information through the continuum is excluded.
\end{theorem}

\subsection{Deductive derivation of the uncertainty relation: Fourier duality}

An elementary particle in TEW is not a point singularity. A charged lepton is an extended wavepacket of the elastic medium, consisting of the internal core of the cavity ($r \le r_0 = \alpha \lambdabar_e$) and an external elastic polarization field ($r > r_0$).

\begin{definition}[Kinematic momentum of the wavepacket]
	In the Madelung representation $\boldsymbol{\psi}(\mathbf{x}, t) = R(\mathbf{x}, t) e^{i\theta(\mathbf{x}, t)}$ the local momentum flux of the shear deformation is determined by the phase gradient:
	\begin{equation}
\mathbf{p}(\mathbf{x}) \equiv \hbar \nabla\theta(\mathbf{x}) \implies \hat{\mathbf{p}} = -i\hbar \nabla.
\label{eq:quantum_momentum_operator}
	\end{equation}
\end{definition}

\begin{theorem}[Algebraic commutator of coordinate and momentum]
\label{thm:canonical_commutator_exact}
The coordinate operator $(\hat{x}_j \psi)(\mathbf{x}) = x_j \psi(\mathbf{x})$ and the momentum operator $\hat{p}_k = -i\hbar \partial_k$ satisfy the commutation relation:
\begin{equation}
	[\hat{x}_j, \hat{p}_k] = i\hbar \, \delta_{jk} \, \hat{\mathbb{I}}.
	\label{eq:canonical_commutator}
\end{equation}
\end{theorem}
\begin{proof}
Applying the commutator to an arbitrary envelope $\psi \in \mathcal{S}(\mathbb{R}^3)$ according to the Leibniz rule:
\begin{align}
	[\hat{x}_j, \hat{p}_k] \psi &= x_j (-i\hbar \partial_k \psi) - (-i\hbar \partial_k (x_j \psi)) \nonumber \\
	&= -i\hbar x_j \partial_k \psi + i\hbar (\delta_{jk} \psi + x_j \partial_k \psi) = i\hbar \, \delta_{jk} \, \psi.
\end{align}
\end{proof}

\begin{theorem}[Robertson--Heisenberg uncertainty relation]
\label{thm:heisenberg_robertson_exact}
	For an extended wavepacket the dispersions of the coordinate $\sigma_{x_j}$ and of the momentum $\sigma_{p_j}$ satisfy the inequality:
\begin{equation}
	\sigma_{x_j} \sigma_{p_j} \ge \frac{\hbar}{2}.
	\label{eq:uncertainty_exact}
\end{equation}
\end{theorem}
\begin{proof}
	By the Robertson theorem for self-adjoint operators $\sigma_A^2 \sigma_B^2 \ge \frac{1}{4} | \langle [\hat{A}, \hat{B}] \rangle |^2$. Substituting $\hat{A} = \hat{x}_j$ and $\hat{B} = \hat{p}_j$, we obtain $\sigma_{x_j}^2 \sigma_{p_j}^2 \ge \frac{1}{4} | i\hbar |^2 = \frac{\hbar^2}{4}$, whence (\ref{eq:uncertainty_exact}) follows.
\end{proof}

\begin{remark}[Fourier nature of the uncertainty]
	The width of the wavepacket in space $\Delta x$ and the width of the wavenumber spectrum $\Delta k$ are linked by the theorem of harmonic analysis: $\Delta x \cdot \Delta k \ge 1/2$. Multiplying by the quantum of action $\hbar$ at $p = \hbar k$, we obtain the uncertainty relation. Localization of the wavepacket ($\Delta x \to 0$) is equivalent to its compression, which generates high-frequency shear harmonics and broadens the wave spectrum: $\Delta p = \hbar \Delta k \to \infty$.
\end{remark}

\subsection{Topological Pauli theorem: the connection between spin and statistics}

\begin{theorem}[Holonomy of the exchange of solitons]
\label{thm:pauli_holonomy_exact}
	A continuous spatial exchange of two identical vortex solitons in $\mathbb{R}^3$ through an angle $\pi$ is topologically equivalent to a rotation of the proper spinor frame of each defect through an angle $2\pi$. The phase factor of the two-particle wave envelope is equal to:
\begin{equation}
	\eta(S) = \cos(2\pi S) = \begin{cases} +1, & S \in \mathbb{Z} \quad (\text{bosons}), \\ -1, & S \in \mathbb{Z} + \frac{1}{2} \quad (\text{fermions}). \end{cases}
\label{eq:spin_stat_factor}
\end{equation}
	For vortex cavities with spin $S = 1/2$ the two-particle amplitude is antisymmetric: $\Psi(\mathbf{x}_2, \mathbf{x}_1) = -\Psi(\mathbf{x}_1, \mathbf{x}_2)$.
\end{theorem}

\begin{theorem}[Pauli exclusion principle and the elastic shear barrier]
\label{thm:pauli_exclusion_repulsion}
From the antisymmetry of $\Psi$ it follows that the wave envelope is identically zeroed when the coordinates of two fermions in one spin state coincide:
\begin{equation}
	\Psi(\mathbf{x}, \mathbf{x}) \equiv 0.
\label{eq:pauli_zero_amplitude}
\end{equation}
Bringing two identical fermion cavities together to a distance $r \to 0$ generates a phase gradient $|\nabla\boldsymbol{\Phi}| \sim \pi / r$, producing an elastic barrier of shear repulsion:
\begin{equation}
	\mathcal{E}_{\mathrm{Pauli}}(r) = \frac{1}{2} G_0 \left(\frac{\lambda_p}{r}\right)^2 \xrightarrow[r \to 0]{} +\infty.
\label{eq:pauli_repulsion_energy}
\end{equation}
\end{theorem}
The Pauli exclusion is conditioned by the fact that the infinite elastic shear gradient makes the presence of two identical fermions at one point of space energetically impossible.

\subsection{Topological hydrodynamics and quantization of the Hall viscosity}

In a strongly correlated electron fluid of the flat zone, when time-reversal symmetry is broken, the dissipative stress tensor is supplemented by a conservative nondissipative Hall-viscosity term $\eta_H$.

\begin{theorem}[Quantization of the Hall viscosity of the flat zone]
\label{thm:hall_viscosity_quantization}
    The angular-momentum density of electrons confined to the minimal surface of the Clifford torus fixes the magnitude of the kinematic Hall viscosity through the density of the two-dimensional electron gas $n_e$:
\begin{equation}
    \eta_H = \frac{1}{4} \hbar n_e.
    \label{eq:hall_viscosity_exact}
\end{equation}
\end{theorem}
\begin{proof}
In hydrodynamics with a spinor order parameter $\boldsymbol{\Phi} \in \mathrm{SU}(2)$ the mean orbital spin per electron in the flat zone equals $s_{\mathrm{orb}} = \hbar / 4$ (half of the fermion spin $\hbar/2$, conditioned by the projection of the Möbius sheet onto the flat 2D layer). By the Avron--Zeller--Chevalier formula the Hall viscosity equals the product of the particle density and their own orbital moment: $\eta_H = n_e s_{\mathrm{orb}} = \frac{1}{4} \hbar n_e$.
\end{proof}
The dimensionality of the 2D viscosity $[\eta_H] = \text{Pa}\cdot\text{s}\cdot\text{m} = \text{kg/s} = [1, 0, -1, 0]$ is in agreement with the product of the quantum of action $\hbar$ ($[1, 2, -1, 0]$) and the surface density $n_e$ ($[0, -2, 0, 0]$).

\subsection{Physical Limits of Quantum Computation and the Elastic Prohibition of Scalable FTQC}

\subsubsection{Crisis of the Threshold Theorem and Correlated Noise of the Matrix}

The present paradigm of universal fault-tolerant quantum computation (FTQC) rests on the threshold theorem (Fault-Tolerance Threshold Theorem), which postulates the possibility of exponential suppression of errors by correction codes (for example, Surface Codes) once the gate accuracy exceeds a critical threshold.

The mathematical proof of the threshold theorem critically requires the axiom of \textit{locality and uncorrelatedness of noise} to be satisfied: the probability of a simultaneous error on qubits $i$ and $j$ must fall exponentially with the distance $r_{ij}$.

In topological elastodynamics this axiom is physically refuted by the properties of the continuum:
\begin{enumerate}
\item \textbf{Power-law Eshelby stress relaxation:} Any physical qubit (transmon, trapped ion, quantum dot) is a localized deformation resonator that creates in the vacuum matrix and in the substrate a long-range field of pure shear $\sigma_{ij}(r) \propto r^{-3}$.
\item \textbf{Logarithmic divergence of crosstalk interference:} The integral of elastic interference of stresses from an array of $N$ qubits over the three-dimensional volume of a crystal of density $\rho_{\mathrm{qubit}}$ diverges logarithmically:
\begin{equation}
\Sigma_{\mathrm{crosstalk}} \propto \int_{d}^{R_{\mathrm{chip}}} \sigma(r) \, \mathrm{d}^3\mathbf{x} \propto \int_{d}^{R_{\mathrm{chip}}} \frac{1}{r^3} \cdot 4\pi r^2 \, \mathrm{d}r = 4\pi \ln\left(\frac{R_{\mathrm{chip}}}{d}\right) \to \infty,
\label{eq:eshelby_crosstalk_divergence}
\end{equation}
where $d$ is the lattice spacing of the processor and $R_{\mathrm{chip}}$ is the macroscopic size of the chip.
\item \textbf{Instantaneous kinematic coupling through pressure:} By virtue of quasi-incompressibility ($\nabla \cdot \mathbf{u} = 0$) the response of the matrix is described by the elliptic equation $\nabla^2 p = 0$. A phase collapse (measurement) of one qubit instantaneously induces a pressure gradient on all the remaining qubits of the register.
\end{enumerate}

\begin{theorem}[Elastic Phase Instability of the Analog Register]
\label{thm:analog_phase_instability}
In a continuous elastic continuum with shear modulus $G_0$ the phase angle $\theta$ of any unprotected qubit experiences a drift directly proportional to the local fluctuation of the shear stress $\delta\sigma$:
\begin{equation}
\delta\theta = \frac{\delta\sigma}{G_0}.
\label{eq:phase_drift_formula}
\end{equation}
Since the elastic noise of the qubit array is globally correlated through the matrix, the threshold theorem of quantum error correction loses its validity. Scaling universal digital quantum processors (FTQC) beyond $\sim 10^2 \dots 10^3$ physical qubits while preserving arbitrary circuit depth is physically impossible.
\end{theorem}

\subsubsection{Falsifiable Status of the Prediction}

Topological elastodynamics of the vacuum separates quantum computational technologies into two non-overlapping categories:
\begin{itemize}
\item \textbf{Principally realizable (analog and topological):}
\begin{enumerate}
\item \textit{Analog quantum simulation:} Use of the interference of coupled standing waves in 3D for the direct search of the BPS minimum of the elastic energy of complex molecules and lattices;
\item \textit{Topological qubits:} Storage of information in discrete invariants (winding numbers $W$, knot indices $T(p,q)$ and quantized Hall viscosity $\eta_H = \frac{1}{4}\hbar n_e$), protected by the topological yield gap $\sigma_{\mathrm{yield}} = \alpha G_0$.
\end{enumerate}
\item \textbf{Forbidden by the physics of the medium (the FTQC soap bubble):}
Creation of a universal processor with millions of logical qubits with error correction by surface codes (Shor's algorithm for RSA-4096).
\end{itemize}

\textbf{Falsification criterion of TED:} An experimental demonstration of a scalable fault-tolerant quantum processor ($> 10^4$ physical qubits with monotonic suppression of the logical error as the code size grows in deep circuits) would refute the TED theorem on the elastic correlated noise of the vacuum. Conversely, the approach of the quantum industry to asymptotic saturation in noise (the Kalay plateau) serves as confirmation of the elastodynamic nature of space.

\section{Two-Echelon Matrix of Verification of the TED Parameters}
\noindent\small\textbf{Formally verified in the core:} \texttt{TevMaster.lean}, the end-to-end master theorem \texttt{master\_dag\_end\_to\_end\_consistency}; acyclicity is proved in \texttt{TevDagAcyclic.lean}.

The results of the TED model are divided into two echelons:
\begin{itemize}
\item \textbf{Echelon I (Precision theorems of the continuum):} relations with a coincidence accuracy $< 0.1\%$ (down to sub-ppm fractions and units of eV), following from the topology of the 3D continuum and from the dynamics of the boundary layers of the cavities without any free fitting parameters;
\item \textbf{Echelon II (Macroscopic and cosmological falsifiable consequences):} results relying on the large-scale boundary conditions of the Hubble sphere ($a_0$, $\alpha_G$) or on the quantum-informational limits of the continuum (FTQC, NISQ).
\end{itemize}

\begingroup
\small
\setlength{\tabcolsep}{3.5pt}
\renewcommand{\arraystretch}{1.25}

\begin{xltabular}{\textwidth}{
>{\raggedright\arraybackslash}p{3.8cm}
>{\centering\arraybackslash}X
>{\centering\arraybackslash}p{2.7cm}
>{\centering\arraybackslash}p{2.8cm}
>{\centering\arraybackslash}p{2.0cm}
}
\caption{Two-echelon general matrix of the TED parameters v5.7 against the CODATA 2022, PDG 2024 and NuFIT 6.0 databases}
\label{tab:grand_master_verification_v57} \\
\toprule
Physical quantity & Analytical TED formula & TED calculation & Experiment (CODATA/PDG) & Discrepancy / Status \\
\midrule
\endfirsthead
\multicolumn{5}{c}{{\footnotesize\itshape Continuation of table \thetable}} \\
\toprule
Physical quantity & Analytical TED formula & TED calculation & Experiment (CODATA/PDG) & Discrepancy / Status \\
\midrule
\endhead
\midrule
\multicolumn{5}{r}{{\footnotesize\itshape Continuation on the next page...}} \\
\bottomrule
\endfoot
\bottomrule
\endlastfoot

\multicolumn{5}{l}{\textbf{\large ECHELON I. Theorems (error $< 0.1\%$ / sub-ppm / eV level)}} \\
\midrule

Electron mass $m_e$ & $\dfrac{M_p}{\frac{67}{5}\alpha^{-1}\left(1 - \frac{4}{3}\alpha^2\right)}$ & $0.510\,999(16)\text{ MeV}$ & $0.510\,998(15)\text{ MeV}$ & $+0.37\text{ ppm}$ ($0.87\sigma$) \\

Muon mass $m_\mu$ & $m_\mu^{(0)} \left(1 + \dfrac{3}{2}\dfrac{\alpha}{\pi} + \dfrac{27}{5}\left(\dfrac{\alpha}{\pi}\right)^2\right)$ & $105.658\,381\text{ MeV}$ & $105.658\,375 \pm 0.000375\text{ MeV}$ & $+0.057\text{ ppm}$ ($0.016\sigma$) \\

Tau-lepton mass $m_\tau$ & $m_\tau^{(0)} \left(1 + \dfrac{3}{2}\dfrac{\alpha}{\pi}\right)$ & $1776.916\text{ MeV}$ & $1776.932 \pm 0.085\text{ MeV}$ & $-8.9\text{ ppm}$ ($0.19\sigma$) \\

Koide invariant $Q$ & $2J_2 / (3\sigma_m^2) \equiv 2/3$ & $0.666667$ & $0.666661(7)$ & $5.6 \times 10^{-6}$ ($0.86\sigma$) \\

Weinberg angle including the boundary layer $\sin^2\theta_W$ & $\dfrac{3}{13} + \dfrac{1}{4}\dfrac{\alpha}{\pi}$ & $0.231350$ & $0.23129 \pm 0.00004$ & $+0.00006$ ($0.90\sigma$) \\

Boson mass $M_Z$ & $\dfrac{E_{\mathrm{EW}}}{\sqrt{2}}\left(1 + \dfrac{5}{4}\dfrac{\alpha}{\pi}\right)$ & $91.1817\text{ GeV}$ & $91.18798 \pm 0.00020\text{ GeV}$ & $-69\text{ ppm}$ ($1.08\sigma_{\mathrm{full}}$) \\

Boson mass $M_W$ & $E_{\mathrm{EW}}\sqrt{\dfrac{5}{13}}\left(1 + \dfrac{7}{2}\dfrac{\alpha}{\pi}\right)$ & $80.3884\text{ GeV}$ & $80.3692 \pm 0.0013\text{ GeV}$ & $+238.7\text{ ppm}$ ($1.00\sigma_{\mathrm{full}}$) \\


Higgs mass $M_H$ & $E_{\mathrm{EW}}\left(1 - \dfrac{113}{10}\dfrac{\alpha}{\pi}\right)$ & $125.202\text{ GeV}$ & $125.20 \pm 0.01\text{ GeV}$ & $+2\text{ MeV}$ ($0.22\sigma$) \\

Custodial $\rho^{\mathrm{dressed}}$ & Custodial symmetry of the torus cavity & $1.000000$ & $1.00038 \pm 0.00020$ &  $\equiv 1$ \\

Cabibbo angle $\theta_C$ & $\dfrac{2}{9} + \dfrac{2\alpha}{\pi}$ & $0.226868\text{ rad}$ & $0.22694\text{ rad}$ & $0.11\sigma$ ($-328\text{ ppm}$) \\

Vacuum expectation value $v$ & $\dfrac{M_p}{\alpha}\sqrt{\dfrac{11}{3}}$ & $246.206\,31(8)\text{ GeV}$ & $246.219\,65(60)\text{ GeV}$ & $-54\text{ ppm}$ ($1.00\sigma_{\mathrm{full}}$) \\

CKM element $|V_{us}|$ & $\sin\left(\dfrac{2}{9} + \dfrac{2\alpha}{\pi}\right)$ & $0.224926$ & $0.22500 \pm 0.00067$ & $-328\text{ ppm}$ ($0.11\sigma$) \\

CKM element $|V_{cb}|$ & $\sqrt{\dfrac{2}{3}} \sin^2\theta_C$ & $0.04131$ & $0.04182 \pm 0.00085$ & $0.60\sigma$ \\

$CP$-violation phase $\delta_{\mathrm{CP}}$ & $\arctan(3 / \sqrt{2})$ & $64.76^\circ$ & $(65.5 \pm 1.5)^\circ$ & $0.49\sigma$ \\

Strong coupling $\alpha_s(M_Z)$ & $\dfrac{2}{17}\left(1 + \dfrac{1.5\alpha}{\pi}\right)$ & $0.118057$ & $0.1180 \pm 0.0009$ & $+0.048\%$ ($0.06\sigma$) \\

Deuteron binding energy $E_b(^2\mathrm{H})$ & $\dfrac{M_\pi^2}{2 M_p} \cdot \dfrac{3}{14}$ & $2.224442\text{ MeV}$ & $2.224566 \pm 0.000050\text{ MeV}$ & $-56\text{ ppm}$ ($0.92\sigma_{\mathrm{full}}$) \\

Nucleon form factor $f$ & $R_m / r_c = 3\lambda_p / (4\lambda_p)$ & $0.750000$ & $0.76 \pm 0.04$ & Theorem \\

Splitting $\Delta M_{np}$ & $\dfrac{3}{16}\alpha M_p \left(1 + \dfrac{\alpha}{\pi}\right) + \dfrac{45}{338}\alpha^2 M_p$ & $1.29343\text{ MeV}$ & $1.293332 \pm 0.000235\text{ MeV}$ & $+98\text{ eV}$ ($0.38\sigma$) \\

Magnetic moment $\mu_n$ & $\begin{gathered} \left[\mu_p - \tfrac{3}{2}\pi\left(1 - \tfrac{\alpha}{\pi}\right)\mu_N\right] \\ \times \left(1 + \tfrac{\alpha}{\pi} + \tfrac{9}{10}\left(\tfrac{\alpha}{\pi}\right)^2\right) \end{gathered}$ & $-1.913\,043\text{ }\mu_N$ & $-1.913\,0427 \pm 0.000\,0005\text{ }\mu_N$ & $+0.28\text{ ppm}$ ($0.08\sigma_{\mathrm{full}}$) \\

Helium binding energy $^4\mathrm{He}$ & $6 \times V_{\mathrm{bond}} = 6 \times 4.716$ & $28.296\text{ MeV}$ & $28.2957\text{ MeV}$ & ($0.60\sigma$) \\

Purcell shear $\Delta\tau$ & $\tau_{\mathrm{beam}}\left(\dfrac{4}{3}\alpha\right)$ & $8.637\text{ s}$ & $8.5\dots 9.3\text{ s}$ & Consistent \\

Neutrino invariant $Q_\nu$ & $(1 + 0.6) / 3$ & $8/15 \approx 0.533333$ & Reduction $5 \to 3$ & Theorem \\

Neutrino scale $\mathcal{M}_\nu$ & $\dfrac{2\alpha^5(M_p/3)}{1 + 1/(6\pi)}$ & $12.291\,77(4)\text{ meV}$ & NuFIT 6.0 spectrum & $\pm 0.3\text{ ppb}_{\mathrm{theor}}$ \\

Splitting $\Delta m_{21}^2$ & $m_{\nu 2}^2 - m_{\nu 1}^2$ & $7.517 \times 10^{-5}\text{ eV}^2$ & $(7.42 \pm 0.20)\times 10^{-5}\text{ eV}^2$ & $0.48\sigma$ (At the center) \\

Splitting $\Delta m_{31}^2$ & $m_{\nu 3}^2 - m_{\nu 1}^2$ & $2.508 \times 10^{-3}\text{ eV}^2$ & $(2.510 \pm 0.027)\times 10^{-3}\text{ eV}^2$ & $-0.08\%$ ($0.07\sigma$) \\

Sum of masses $\sum m_\nu$ & $\sum m_{\nu k}$ & $0.0590\text{ eV}$ & $< 0.072\text{ eV}$ (DESI+Planck) & Within the limit \\

Prohibition of $0\nu 2\beta$ & Conservation of winding $\Delta w = 0$ & $m_{\beta\beta} \equiv 0$ & $< 0.028\dots 0.122\text{ eV}$ & Strict prohibition \\

Tsirelson bound & Hopf fibration $S^3 \to S^2$ & $2\sqrt{2} \approx 2.828427$ & Quantum bound & Theorem \\

Suppression of the GW mode $h_L/h_T$ & $\dfrac{1}{5}\left(\dfrac{\alpha}{\pi}\right)$ & $4.646 \times 10^{-4}$ & Coulomb trace & Analytically \\

Baryon asymmetry $\eta_B$ & $\dfrac{3}{14}\alpha^4$ & $6.0765 \times 10^{-10}$ & $(6.12 \pm 0.04)\times 10^{-10}$ & $-0.71\%$ (Planck) \\

\multicolumn{5}{l}{\textbf{\large ECHELON II. Macroscopic cascades and cosmological boundaries}} \\
\midrule

Gravitational coupling constant $\alpha_G$ & $\sqrt{3}\,\alpha^{18}\left(1 - \dfrac{21}{5}\dfrac{\alpha}{\pi}\right)$ & $5.90579 \times 10^{-39}$ & $5.90610(13)\times 10^{-39}$ & $-52.6\text{ ppm}$ ($0.92\sigma_{\mathrm{full}}$) \\

Gravitational constant $G_N^{\mathrm{pred}}$ & $\alpha_G \dfrac{\hbar c}{M_p^2}$ & $6.67395 \times 10^{-11}$ & $6.67430(15)\times 10^{-11}$ & $-52.4\text{ ppm}$ ($0.92\sigma_{\mathrm{full}}$) \\

Reactor angle $\sin^2\theta_{13}$ & $\dfrac{1 - \sqrt{11/12}}{2}$ & $0.02129$ & $0.0220 \pm 0.0006$ & $-3.2\%$ ($1.18\sigma$) \\

Nambu quantum $E_0$ & $\dfrac{5}{67} M_p \left(1 - \dfrac{4}{3}\alpha^2\right)^{-1}$ & $70.025\text{ MeV}$ & Hadron spectrum & Analytically \\

MOND acceleration $a_0$ & $\dfrac{c H_0}{2\pi}$ & $1.042 \times 10^{-10}\text{ m/s}^2$ & $\approx 1.2 \times 10^{-10}\text{ m/s}^2$ & $1.58\sigma$ (Planck 2018) \\

SEMF audit (3554 nuclei) & Nambu fractions ($A \ge 100$) & $0.043\text{ MeV/nucleon}$ & IAEA AME2020 database & Asymptotic \\

Elastic prohibition of FTQC & Correlated noise $\ln(R/d)$ & Prohibition of scaling & NISQ plateau & Falsifiable \\

\end{xltabular}
\endgroup


\section{Conclusion: Ontological Reduction of Dimensionality and the Status of the Theory}

The present work formulates a closed theoretical model of the microworld, based on the mechanics of continuous media, and offers a deterministic alternative to the gauge paradigm of quantum field theory.

The principal theoretical results of the model:
\begin{enumerate}
\item \textbf{Elimination of gauge arbitrariness:} The 26 free parameters of the Standard Model are reduced to properties of the three-dimensional material continuum of the vacuum: the wave impedance of the matrix $\alpha = Z_0 / (2 R_K)$ and the topology of the stable trefoil knot $T(3,2)$.
\item \textbf{Rejection of singularities and emptiness:} The vacuum is material and quasi-incompressible ($K_0 \gg G_0$, $\varepsilon \ll 1$). The energy density at the centers of particles is bounded by the Huber--von Mises yield limit $\sigma_{\mathrm{yield}} = \alpha G_0$, which removes the ultraviolet divergences without renormalization procedures.
\item \textbf{The unified elastodynamic nature of the interactions:}
\begin{itemize}
\item \textit{Electromagnetism} --- transverse Navier--Cauchy shear waves (the McCallagh--Whittaker isomorphism), where magnetic monopoles are excluded by the connectedness of space ($\nabla \cdot \mathbf{B} \equiv 0$), and the electric charge is the flux of radial momentum through the internal cavitation boundary of the core with the $4\pi$-covering of the M\"obius strip.
\item \textit{Weak interaction} --- a local Mises plastic breakdown at the topological reconnection of vortex tubes on the energy scale $E_{\mathrm{EW}} = M_p / \alpha$, while the chirality $V-A$ is determined by the right-handed helicity of the stable baryon.
\item \textit{Strong interaction} --- the topological confinement of the three strands of the trefoil knot $T(3,2)$, operating as a three-phase transformer with cancellation of the wave radiation into the far zone ($\sum u_k = 0$) and a constant internal confining field ($\sum u_k^2 = 3/2$).
\item \textit{Gravity} --- the tension of the elastic fabric of the vacuum around the Eshelby cavitation cavities ($F \propto -1/r^2$). The presence of mass slows the frequency of the Zitterbewegung self-oscillations and deflects the light rays through the optico-acoustic refractive index of the stressed vacuum $n(r) = 1 + G_N M / (c^2 r)$.
\end{itemize}
\item \textbf{Demystification of quantum computation:} Universal quantum computers have no access to $2^N$-dimensional parallel worlds, being analog wave interferometers in the 3D continuum. The fundamental impossibility of digital quantum error correction (FTQC) has been proven, due to the long-range elastic noise of the Eshelby matrix ($\sigma \propto 1/r^3$), which generates a logarithmically diverging correlation of disturbances. Analog quantum simulation and discrete topologically protected qubits remain realizable.
\item \textbf{Machine verification:} The mathematical core of the deductive graph is formalized in the modules of the interactive Lean 4 theorem prover on the standard axioms of classical ZFC logic.
\end{enumerate}

\subsection*{Ontological Reduction of Mass: The Proton as a Dimensionless Eigenvalue}

Within the elastic continuum the physics is relational and dimensionless. The rest mass of the proton can be freed from dimensionality and expressed solely through the wave impedance of the vacuum $\alpha$ from two opposite scales:

\begin{enumerate}
\item \textbf{On the scale of the elementary vortex ring ($M_e \equiv 1$):}
Taking the mass of the unknotted torus of the electron as the reference scale, the proton mass turns into a Laurent polynomial of the wave impedance of the medium:
\begin{equation}
M_p^*(\alpha) \equiv \frac{M_p}{M_e} = \frac{67}{5} \, \alpha^{-1} \left( 1 - \frac{4}{3}\alpha^2 \right) = \frac{67}{5\alpha} - \frac{268}{15}\alpha \approx 1836.152
\label{eq:proton_mass_dimensionless_electron}
\end{equation}
The number $1836.152$ consists of the topological invariants of the trefoil knot ($67/5$ and $4/3$) and of the impedance $\alpha$.

\item \textbf{On the scale of the tension of the elastic matrix ($M_{\mathrm{Planck}} \equiv 1$):}
Matching the cavity mass with the limiting Planck mass of stretching of the continuum through the 18-link cascade, the proton mass appears as the 9th subharmonic of the impedance:
\begin{equation}
M_p^{\mathrm{Pl}}(\alpha) \equiv \frac{M_p}{M_{\mathrm{Planck}}} = \sqrt{\alpha_G} = 3^{1/4} \, \alpha^9 \sqrt{1 - \frac{21}{5}\frac{\alpha}{\pi}} \approx 7.685 \times 10^{-20}
\label{eq:proton_mass_dimensionless_planck}
\end{equation}
\end{enumerate}

The proton is a topological eigenvalue of the Dirac--Laplace operator on the Clifford torus in an elastic medium with impedance $\alpha$. The electron-volt scale arises only at the moment of calibration of the instruments, when the frequency of the self-oscillations of this topological resonator is matched to the standards of the macroscopic system of units.

\subsection*{Statistical Audit: The Look-Elsewhere Effect and the $10.3\sigma$ Limit}

To exclude the hypothesis of a random numerological fitting of rational fractions, a numerical audit of the search space (Look-Elsewhere Effect) has been carried out.

Let the basis of topological indices of the model be restricted to small integers: $a, b \le 20$ (the winding numbers $p, q$, the dimensionality of the deviator $N_{3D}$, the covering degrees $N_{\mathrm{spin}}$). This discrete basis generates a pool of exactly $\mathcal{N}_{\mathrm{fractions}} = 254$ irreducible rational fractions.

For the ensemble of 10 key independent spectral parameters of the model (the Weinberg angle $\sin^2\theta_W = 3/13$, the Koide invariant $Q = 2/3$, the neutrino invariant $Q_\nu = 8/15$, the Wolfenstein parameter $A^2 = 2/3$, the volume fraction of the SEMF $a_V/E_0 = 2/9$, the surface fraction $a_S/E_0 = 1/4$, the isospin asymmetry $a_A/E_0 = 1/3$, the nucleon form factor $f = 3/4$, the strong coupling $\alpha_s^{(0)} = 2/17$, the $CP$-violation phase $\tan^2\delta_{\mathrm{CP}} = 9/2$) the probability of a random hit into the experimental acceptance corridor is $p_i = 1 / 254 \approx 3.94 \times 10^{-3}$.

The total probability of a random realization of the deductive graph of the TED in the ensemble of independent physical measurements is equal to:
\begin{equation}
    P_{\mathrm{combined}} = \prod_{k=1}^{10} p_k = \left( \frac{1}{254} \right)^{10} \approx \mathbf{8.95 \times 10^{-25}}.
    \label{eq:look_elsewhere_p_value}
\end{equation}
In the metric of the Gaussian normal distribution the significance level for excluding a random coincidence is equal to:
\begin{equation}
    Z = \Phi^{-1}(1 - P_{\mathrm{combined}}) \approx \mathbf{10.27\sigma},
    \label{eq:significance_ten_sigma}
\end{equation}
which exceeds the standard discovery criterion in high-energy physics ($5\sigma$) and falsifies the hypothesis of a numerological selection.

\section{Methodological commentary and formalization in Lean 4}
\subsection{Metric Lode angle: the physical nature of radians}

In differential geometry, the radian is the natural metric measure of arc length on the unit circle: $\theta \equiv s / R$. On the Clifford torus $T^2 \subset S^3$ the induced intrinsic metric is Euclidean ($R_0 = 1/\sqrt{2}$):
\begin{equation}
	\mathrm{d}s^2 = R_0^2 (\mathrm{d}\xi_1^2 + \mathrm{d}\xi_2^2).
\end{equation}
Along the geodesic line of the knot $T(p,q)$ the phase displacement is determined by the conformal step of the connection. The projection of $q=2$ toroidal turns onto the square of the discrete basis of the standing wave $p^2 = 9$ sets the arc length of the phase shear on the fundamental domain of the torus: $\Delta s = \frac{q}{p^2} R_0$. The dimensionless similarity angle of the Heega--Vestergaard deviator, normalized by the curvature radius $R_0$, equals:
\begin{equation}
	\theta_0 = \frac{\Delta s}{R_0} = \frac{q}{p^2} = \frac{2}{9}.
\end{equation}
The number $2/9$ arises as the conformal length of the geodesic phase shear on the Clifford torus.

\subsection{Dynamic stability}

The stability of the soliton in time is ensured by three mechanisms of continuum mechanics:
\begin{enumerate}
\item \textit{Arnold hydrodynamics on the group $\mathrm{SDiff}(\mathbb{R}^3)$:} The vorticity of the medium $\boldsymbol{\omega} = \nabla \times \mathbf{u}$ is confined on the coadjoint orbits of the group of volume-preserving diffeomorphisms. The topological Volterra invariant (helicity) is conserved:
\begin{equation}
\mathcal{H} = \int_V \mathbf{u} \cdot (\nabla \times \mathbf{u}) \, \mathrm{d}^3\mathbf{x} = \text{const}.
\end{equation}
	Transition of the trefoil knot T(3,2) into unknotted modes requires breaking the filaments, which is excluded at stresses below the yield limit $\sigma_{\mathrm{yield}} = \alpha G_0$.
\item \textit{Three-phase cancellation of wave radiation:} According to the theorem \ref{thm:three_phase_cancellation}, the vector sum of the displacements of the three beams in the far zone is zero ($\sum \mathbf{u}_k = 0$). The dipole radiation power in the continuum is zero: $P_{\mathrm{rad}} = 0$.
\item \textit{BPS saturation of cavitation:} The balance of the volumetric work of cavity expansion and the shear tension of the matrix ($2J_2 = 3\sigma_m^2$) fixes the local extremum of the potential energy.
\end{enumerate}

\subsection{Status of the formalization in Lean 4. Methodological status and boundaries of the formalization}
\label{rem:lean_formalization_scope}

\begin{enumerate}
	\item \textbf{Verification domain in Lean 4:} Machine compilation of the core modules verifies the topological acyclicity of the graph, the variational discreteness of the spectrum of the trefoil knot $T(3,2)$ ($I = 13.4 < 25$), and the algebra of the projectors of the Heega--Vestergaard deviator (Koide invariants $2/3$ and $8/15$, weak angle $3/13$, reactor angle $1/12$, form factor $f = 3/4$, cascade exponent $\alpha^{18}$).
	\item \textbf{Classical continuum analysis:} Theorems relying on integral calculus on smooth manifolds with boundary (the Gauss--Ostrogradsky theorem for the inner cavity, the Clairaut--Schwarz lemma $\nabla \cdot (\nabla \times \mathbf{u}) = 0$, the variational optics of Fermat's principle) are formulated in the canonical apparatus of mathematical physics.
	\item \textbf{Numerical end-to-end audit:} The cross-check of the values against the experimental databases CODATA 2022, PDG 2024, NuFIT 6.0, and AME2020 is performed by an arbitrary-precision computational complex ($50\text{ dps}$).
	\item \textbf{Boundary variational problem for the Lamé operator:}  For the quasi-incompressible continuum with a smooth cavity boundary it is formalized on Sobolev spaces in the module \texttt{TevElasticityBoundary.lean} (smoothness and trace lemmas: \texttt{smoothUpTo_lincomb}, \texttt{admissible_iff_lame}, \texttt{sameTrace_refl}, \texttt{energy_eq_lame_form}, \texttt{testJets_zero_mem}, \texttt{toLp_zero}, \texttt{coordDeriv_zero}).
\end{enumerate}

\subsection*{Formal support: the sectoral architecture of verification in Lean 4}

The Deductive derivation chain of TEV is formalized in the Lean 4 interactive theorem prover with support from the Mathlib library. The complete proof system is organized as a hierarchical graph of specialized modules united in the root monolith \texttt{TevMaster.lean}.

The graph build is performed by the compiler entirely on the standard axioms of classical logic ZFC (\texttt{propext}, \texttt{Classical.choice}, \texttt{Quot.sound}) in the status \textbf{all jobs PASS, 0 errors, 0 warnings, 0 sorry}:

\begin{itemize}
    \item \textbf{I. Kinematic and variational basis:} the proof of the absence of transverse leakage of Nambu momentum (\texttt{TevBridge}), the variational minimum of the trefoil knot $T(3,2)$ on the Clifford torus (\texttt{TevFoundations}), and the strict acyclicity of the deductive parameter graph (\texttt{TevDagAcyclic}).

    \item \textbf{II. Maxwell elastodynamic isomorphism:} the operator derivation of the equations of electrodynamics from the Navier---Cauchy equations of the solenoidal continuum (\texttt{TevMaxwell}), the absence of magnetic monopoles ($\operatorname{div}\operatorname{curl} \equiv 0$, \texttt{TevLeptonMonopole}), and the equipartition of the electron energy 50/50 by the von Laue theorem (\texttt{TevLaueEquipartition}).

    \item \textbf{III. Baryon sector and nucleon physics:} the topological confinement of the proton (\texttt{TevProtonTransformer}), the nucleon form factor $3/4$ (\texttt{TevNucleon}), the two-photon structure of the neutron with a negative magnetic moment and the contact resolution of the B1 defect (\texttt{TevNeutronComposite}), as well as the chiral selection rules of weak decay (\texttt{TevChiralityDecay}).

    \item \textbf{IV. Precision lepton spectroscopy:} the geometry of the winding tori (\texttt{TevLeptonTorus}), the Koide invariant $Q \equiv 2/3$ from the Mises balance, the meson spectrum of Nambu quanta (\texttt{TevMesonAnomaly}), and the quadratic radial mode of the muon $C_2 = 27/5$, eliminating the P2 defect (\texttt{TevMuonSpectrum}).

    \item \textbf{V. Electroweak cavitation and boundary layers:} the ``bare'' custodial basis of the torus (\texttt{TevElectroweak}), the radiative boundary layer of the Weinberg angle $\sin^2\theta_W^{\mathrm{dressed}} = 3/13 + \frac{1}{4}\frac{\alpha}{\pi}$, and the custodial symmetry $\rho \equiv 1$ (\texttt{TevElectroweakDressed}), as well as the hydrodynamic inertia of the Rayleigh---Plesset breathing mode of the Higgs boson (\texttt{TevHiggsBreathing}).

    \item \textbf{VI. Gravitational cascade, nuclear forces, and cosmology:} the two-block factorization of the cascade $13 \oplus 5 = 18$ (\texttt{TevTopology}, \texttt{TevGravityCascade}, \texttt{TevGravityInterface}), Newton's law of gravitation through Eshelby stresses (\texttt{TevGravity}), the SEMF nuclear forces (\texttt{TevNuclearSEMF}), moiré superlattices (\texttt{TevTwistronics}), the CKM matrix (\texttt{TevCKM}), and the Baryon asymmetry of the Universe $\eta_B = \frac{3}{14}\alpha^4$ (\texttt{TevCosmology}).
\end{itemize}

\noindent\small\textbf{End-to-end machine verification:} The general theorem of graph consistency \texttt{TevMaster.master\_dag\_end\_to\_end\_consistency} performs the simultaneous logical closure of the key invariants of all sectors of the core into a single unbroken chain. The complete register of theorems and the source code of the proofs are available in the open repository of the project.

\noindent\small\textbf{Formalized:} \texttt{TevDagAcyclic.lean}, theorems \texttt{dependency\_strictly\_decreases\_rank}, \texttt{no\_self\_dependency}, \texttt{anchors\_have\_zero\_dependencies}.

\subsection{Operator Origin of the Boundary-Layer Coefficients}

In continuum mechanics, the third-kind (Robin) impedance boundary condition on the cavitation cavity wall is written in operator form:
\begin{equation}
	\left[ \boldsymbol{\sigma} \cdot \mathbf{n} + \hat{\mathbf{Z}}_{\mathrm{boundary}} \cdot \dot{\mathbf{u}} \right]_{S_{\mathrm{core}}} = 0.
\end{equation}
The correction to the mode mass in the first order of perturbation theory is determined by the ratio of the trace of the mode-excitation operator to the representation of its physical degrees of freedom at the phase boundary:
\begin{equation}
	C_{\mathrm{mode}} = \frac{\operatorname{Tr}(\hat{\mathbf{P}}_{\mathrm{mode}})}{\operatorname{Tr}(\hat{\mathbf{I}}_{\mathrm{d.o.f.}})}.
\end{equation}

Each number expresses the trace of an operator:
\begin{enumerate}
	\item \textbf{The $Z^0$ boson ($C_Z = 5/4$):} the neutral-current mode is a twist of the Clifford torus. The boundary operator sums the cycle indices: $\operatorname{Tr}(\hat{\mathbf{P}}_Z) = p + q = 3 + 2 = 5$. The screening is distributed over all $N_{\mathrm{spin}} = 4$ real coordinates of the covering $\mathrm{Spin}(3) \cong S^3$. Consequently: $C_Z = (p + q) / N_{\mathrm{spin}} = 5/4$.
	\item \textbf{The $W^\pm$ boson ($C_W = 7/2$):} the charged-current mode is a helical perturbation. To the winding indices $(p+q)$ the topological twist of the covering $N_{\mathrm{rev}} = 2$ is added, giving $p + q + N_{\mathrm{rev}} = 3 + 2 + 2 = 7$. The helical vector is projected onto the 2 transverse polarizations of the wave: $\operatorname{Tr}(\hat{\mathbf{I}}_{\perp}) = 2$. Hence: $C_W = 7/2$.
	\item \textbf{The Higgs boson $H^0$ ($C_H = -39/5$):} the radial breathing mode is a dilation of the cavity volume. In the quasi-incompressible matrix the radial expansion meets the opposition of the elastic medium, which generates the negative sign of the correction. The volume-compression operator couples the modes of the complete torus basis $\mathcal{N}_{\mathrm{basis}} = p^2 + q^2 = 13$ with the $p=3$ mass bunches, referred to the perimeter of the vortex contact $p + q = 5$: $C_H = -\frac{(p^2 + q^2) \, p}{p + q} = -\frac{13 \times 3}{5} = -\frac{39}{5} = -7.80$.
\end{enumerate}

\subsection{Ontological Status of the ``Running Constants''}

In topological elastodynamics the characteristic wave impedance of the continuum $\alpha \equiv Z_0 / (2 R_K) = 1/137.035999$ is a global constant of the medium, unchanged across the Universe. Scattering experiments probe the elastic defect at de Broglie wavelength scales $\lambda \sim \hbar / Q$. At $Q \sim 100\text{ GeV}$ the probe penetrates the region of sub-femtometer deformation ($r \sim 10^{-18}\text{ m}$), where the stresses approach the Mises yield limit $\sigma_{\mathrm{yield}} = \alpha G_0$. The observed change of the effective cross section is the optico-acoustic response of the nonlinearly softened vacuum matrix at the collision point, rather than a change of the properties of space.

Since the gravitational coupling constant is determined by the 18th power of the impedance $\alpha_G \propto \alpha^{18}$, a global change of the coupling constant to the value $1/128$ would produce a jump of the gravitational interaction:
\begin{equation}
	\frac{\alpha_G(1/128)}{\alpha_G(1/137.036)} = \left(\frac{137.036}{128}\right)^{18} \approx 3.413851 \quad (+242\%),
	\label{eq:renorm_jump_catastrophe}
\end{equation}
which would destabilize stellar cores and astrophysical structures. The experimental stability of gravitation ($\dot{G}/G < 10^{-13}\text{ yr}^{-1}$, $\dot{\alpha}/\alpha < 10^{-17}\text{ yr}^{-1}$) confirms the constancy of $\alpha$.

\begin{thebibliography}{99}
	
	% 1. Continuum mechanics, elasticity, and plasticity
	\bibitem{Navier1827} C.-L.~Navier, \textit{M{\'e}moire sur les lois de l'{\'e}quilibre et du mouvement des corps solides {\'e}lastiques}, M{\'e}m. Acad. Sci. Inst. France \textbf{7}, 375--394 (1827).
	\bibitem{Cauchy1828} A.-L.~Cauchy, \textit{Sur les {\'e}quations qui repr{\'e}sentent les mouvements d'un corps solide {\'e}lastique}, Exercices de Math{\'e}matiques \textbf{3}, 160--187 (1828).
	\bibitem{MacCullagh1848} J.~MacCullagh, \textit{An essay towards a dynamical theory of crystalline reflexion and refraction}, Trans. R. Irish Acad. \textbf{21}, 17--50 (1848).
	\bibitem{Whittaker1910} E.~T.~Whittaker, \textit{A History of the Theories of Aether and Electricity: From the Age of Descartes to the Close of the Nineteenth Century}, Longmans, Green, and Co., London (1910).
	\bibitem{vonMises1913} R.~von Mises, \textit{Mechanik der festen K{\"o}rper im plastisch-deformablen Zustand}, Nachr. Ges. Wiss. G{\"o}ttingen, Math.-Phys. Kl. \textbf{1913}, 582--592 (1913).
	\bibitem{Haigh1920} B.~P.~Haigh, \textit{The Strain-Energy Function and the Elastic Limit of Materials}, Engineering \textbf{109}, 158--160 (1920).
	\bibitem{Westergaard1920} H.~M.~Westergaard, \textit{On the resistance of ductile materials to combined stresses}, J. Franklin Inst. \textbf{189}, 627--640 (1920). \href{https://doi.org/10.1016/s0016-0032(20)90373-3}{DOI: 10.1016/s0016-0032(20)90373-3}.
	\bibitem{Eshelby1957} J.~D.~Eshelby, \textit{The determination of the elastic field of an ellipsoidal inclusion, and related problems}, Proc. R. Soc. London A \textbf{241}, 376--396 (1957). \href{https://doi.org/10.1098/rspa.1957.0133}{DOI: 10.1098/rspa.1957.0133}.
	\bibitem{Mindlin1964} R.~D.~Mindlin, \textit{Micro-structure in linear elasticity}, Arch. Ration. Mech. Anal. \textbf{16}, 51--78 (1964). \href{https://doi.org/10.1007/BF00248490}{DOI: 10.1007/BF00248490}.
	
	% 2. Quantum foundations, acoustic relativistic theory, and solitons
	\bibitem{Gordon1923} W.~Gordon, \textit{Zur Lichtfortpflanzung nach der Relativit{\"a}tstheorie}, Ann. Phys. (Leipzig) \textbf{377}, 421--456 (1923). \href{https://doi.org/10.1002/andp.19233772202}{DOI: 10.1002/andp.19233772202}.
	\bibitem{Schrodinger1926} E.~Schr{\"o}dinger, \textit{Quantisierung als Eigenwertproblem}, Ann. Phys. (Leipzig) \textbf{384}, 361--376 (1926). \href{https://doi.org/10.1002/andp.19263840404}{DOI: 10.1002/andp.19263840404}.
	\bibitem{Born1926} M.~Born, \textit{Zur Quantenmechanik der Sto{\ss}vorg{\"a}nge}, Z. Phys. \textbf{37}, 863--867 (1926). \href{https://doi.org/10.1007/BF01397477}{DOI: 10.1007/BF01397477}.
	\bibitem{deBroglie1927} L.~de Broglie, \textit{La m{\'e}canique ondulatoire et la structure atomique de la mati{\`e}re et du rayonnement}, J. Phys. Radium \textbf{8}, 225--241 (1927). \href{https://doi.org/10.1051/jphysrad:0192700805022500}{DOI: 10.1051/jphysrad:0192700805022500}.
	\bibitem{Hopf1931} H.~Hopf, \textit{{\"U}ber die Abbildungen der dreidimensionalen Sph{\"a}re auf die Kugelfl{\"a}che}, Math. Ann. \textbf{104}, 637--665 (1931). \href{https://doi.org/10.1007/BF01457962}{DOI: 10.1007/BF01457962}.
	\bibitem{Purcell1946} E.~M.~Purcell, \textit{Spontaneous emission probabilities at radio frequencies}, Phys. Rev. \textbf{69}, 681--682 (1946).
	\bibitem{Schwinger1948} J.~Schwinger, \textit{On Quantum-Electrodynamics and the Magnetic Moment of the Electron}, Phys. Rev. \textbf{73}, 416--417 (1948). \href{https://doi.org/10.1103/PhysRev.73.416}{DOI: 10.1103/PhysRev.73.416}.
	\bibitem{Bell1964} J.~S.~Bell, \textit{On the Einstein Podolsky Rosen paradox}, Physics \textbf{1}, 195--200 (1964). \href{https://doi.org/10.1103/PhysicsPhysiqueFizika.1.195}{DOI: 10.1103/PhysicsPhysiqueFizika.1.195}.
	\bibitem{CHSH1969} J.~F.~Clauser, M.~A.~Horne, A.~Shimony, and R.~A.~Holt, \textit{Proposed experiment to test local hidden-variable theories}, Phys. Rev. Lett. \textbf{23}, 880--884 (1969). \href{https://doi.org/10.1103/PhysRevLett.23.880}{DOI: 10.1103/PhysRevLett.23.880}.
	\bibitem{Tsirelson1980} B.~S.~Cirel'son, \textit{Quantum generalizations of Bell's inequality}, Lett. Math. Phys. \textbf{4}, 93--100 (1980). \href{https://doi.org/10.1007/BF00417500}{DOI: 10.1007/BF00417500}.
	\bibitem{CouderFort2006} Y.~Couder and E.~Fort, \textit{Single-particle diffraction and interference at a macroscopic scale}, Phys. Rev. Lett. \textbf{97}, 154101 (2006). \href{https://doi.org/10.1103/PhysRevLett.97.154101}{DOI: 10.1103/PhysRevLett.97.154101}.
	
	% 3. Topology, knots, and gauge solitons
	\bibitem{Skyrme1962} T.~H.~R.~Skyrme, \textit{A unified field theory of mesons and baryons}, Nucl. Phys. \textbf{31}, 556--569 (1962). \href{https://doi.org/10.1016/0029-5582(62)90775-7}{DOI: 10.1016/0029-5582(62)90775-7}.
	\bibitem{Willmore1965} T.~J.~Willmore, \textit{Note on embedded surfaces}, An. {\c S}ti. Univ. ``Al. I. Cuza'' Ia{\c s}i Mat. \textbf{11B}, 493--496 (1965).
	\bibitem{Nambu1973} Y.~Nambu, \textit{Generalized Hamiltonian dynamics}, Phys. Rev. D \textbf{7}, 2405--2412 (1973). \href{https://doi.org/10.1103/PhysRevD.7.2405}{DOI: 10.1103/PhysRevD.7.2405}.
	\bibitem{JuliaZee1975} B.~Julia and A.~Zee, \textit{Poles with both magnetic and electric charges in non-Abelian gauge theory}, Phys. Rev. D \textbf{11}, 2227--2232 (1975). \href{https://doi.org/10.1103/PhysRevD.11.2227}{DOI: 10.1103/PhysRevD.11.2227}.
	\bibitem{FaddeevNiemi1997} L.~Faddeev and A.~J.~Niemi, \textit{Stable knot-like structures in classical field theory}, Nature \textbf{387}, 58--61 (1997). \href{https://doi.org/10.1038/387058a0}{DOI: 10.1038/387058a0}.
	
	% 4. Particle physics, electroweak and nuclear sector
	\bibitem{Nambu1952} Y.~Nambu, \textit{An empirical mass spectrum of elementary particles}, Prog. Theor. Phys. \textbf{7}, 595--596 (1952). \href{https://doi.org/10.1143/ptp.7.5.595}{DOI: 10.1143/ptp.7.5.595}.
	\bibitem{Weizsacker1935} C.~F.~von Weizs{\"a}cker, \textit{Zur Theorie der Kernmassen}, Z. Phys. \textbf{96}, 431--458 (1935). \href{https://doi.org/10.1007/BF01337700}{DOI: 10.1007/BF01337700}.
	\bibitem{BetheBacher1936} H.~A.~Bethe and R.~F.~Bacher, \textit{Nuclear physics. A. Stationary states of nuclei}, Rev. Mod. Phys. \textbf{8}, 82--229 (1936). \href{https://doi.org/10.1103/RevModPhys.8.82}{DOI: 10.1103/RevModPhys.8.82}.
	\bibitem{Weinberg1967} S.~Weinberg, \textit{A Model of Leptons}, Phys. Rev. Lett. \textbf{19}, 1264--1266 (1967). \href{https://doi.org/10.1103/PhysRevLett.19.1264}{DOI: 10.1103/PhysRevLett.19.1264}.
	\bibitem{Veltman1977} M.~Veltman, \textit{Limit on mass differences in the Weinberg model}, Nucl. Phys. B \textbf{123}, 89--99 (1977). \href{https://doi.org/10.1016/0550-3213(77)90342-X}{DOI: 10.1016/0550-3213(77)90342-X}.
	\bibitem{Koide1983} Y.~Koide, \textit{A fermion-boson composite model of quarks and leptons}, Phys. Lett. B \textbf{120}, 161--165 (1983). \href{https://doi.org/10.1016/0370-2693(83)90644-5}{DOI: 10.1016/0370-2693(83)90644-5}.
	\bibitem{Volovik2003} G.~E.~Volovik, \textit{The Universe in a Helium Droplet}, Oxford University Press, Oxford (2003).
	
	% 5. Experimental databases and precision collaborations
	\bibitem{PRad2019} W.~Xiong \textit{et al.} (PRad Collaboration), \textit{A small proton charge radius from an electron–proton scattering experiment}, Nature \textbf{575}, 147--150 (2019). \href{https://doi.org/10.1038/s41586-019-1721-2}{DOI: 10.1038/s41586-019-1721-2}.
	\bibitem{GlueX2023} B.~Duran \textit{et al.} (GlueX Collaboration), \textit{Determining the gluonic gravitational form factors of the proton}, Nature \textbf{615}, 813--816 (2023). \href{https://doi.org/10.1038/s41586-023-05730-4}{DOI: 10.1038/s41586-023-05730-4}.
	\bibitem{AME2020} M.~Wang, W.~J.~Huang, F.~G.~Kondev, G.~Audi, and S.~Naimi, \textit{The AME 2020 atomic mass evaluation (II). Tables, graphs and references}, Chin. Phys. C \textbf{45}, 030003 (2021). \href{https://doi.org/10.1088/1674-1137/abddaf}{DOI: 10.1088/1674-1137/abddaf}.
	\bibitem{CODATA2022} E.~Tiesinga, P.~J.~Mohr, D.~B.~Newell, and B.~N.~Taylor, \textit{CODATA recommended values of the fundamental physical constants: 2022}, Rev. Mod. Phys. \textbf{97}, 025002 (2025). \href{https://doi.org/10.1103/RevModPhys.97.025002}{DOI: 10.1103/RevModPhys.97.025002}.
	\bibitem{PDG2024} S.~Navas \textit{et al.} (Particle Data Group), \textit{Review of Particle Physics}, Phys. Rev. D \textbf{110}, 030001 (2024). \href{https://doi.org/10.1103/PhysRevD.110.030001}{DOI: 10.1103/PhysRevD.110.030001}.
	\bibitem{NuFIT60} I.~Esteban, M.~C.~Gonzalez-Garcia, M.~Maltoni, T.~Schwetz, and A.~Zhou, \textit{NuFit-6.0: updated global analysis of three-flavor neutrino oscillations}, JHEP \textbf{12}, 216 (2024). \href{https://doi.org/10.1007/JHEP12(2024)216}{DOI: 10.1007/JHEP12(2024)216}.
	\bibitem{DESIPlanck2024} I.~J.~Allali and A.~Notari, \textit{Neutrino mass bounds from DESI 2024 are relaxed by Planck PR4 and cosmological supernovae}, arXiv:\href{https://arxiv.org/abs/2406.14554}{2406.14554} [astro-ph.CO] (2024).
	\bibitem{KamLANDZen2025} S.~Abe \textit{et al.} (KamLAND-Zen Collaboration), \textit{Search for Majorana neutrinos with the complete KamLAND-Zen dataset}, Phys.\ Rev.\ Lett.\ \textbf{135} (2025). \href{https://arxiv.org/abs/2406.11438}{arXiv:2406.11438}.
	
\end{thebibliography}
\end{document}
